% The Erdos-Straus equation through residual shells: the workbench and the author's notes (version 4).
% Prepared by Claude (Anthropic), model claude-opus-5-5 (Opus 5.5).
% Compile with lualatex (three times).
% The Erdos-Straus equation through residual shells: the workbench and the author's notes (version 4).
% Prepared by Claude (Anthropic), model claude-opus-5-5 (Opus 5.5).
% Compile with lualatex (three times).
\documentclass[11pt,a4paper]{article}
\usepackage[a4paper,margin=24mm,headheight=14pt]{geometry}
\usepackage{fontspec}
\directlua{luaotfload.add_fallback("readerfallback",{"TeX Gyre Pagella Math:mode=harf;","STIX:mode=harf;","DejaVu Serif:mode=harf;","FreeSerif:mode=harf;"})}
\setmainfont{TeX Gyre Pagella}[RawFeature={fallback=readerfallback}]
\setsansfont{TeX Gyre Heros}[Scale=MatchLowercase]
\setmonofont{DejaVu Sans Mono}[Scale=0.84]
\usepackage{amsmath,amsthm,mathtools}
\usepackage{graphicx}
\usepackage{unicode-math}
\setmathfont{TeX Gyre Pagella Math}
\usepackage{microtype}
\usepackage{booktabs,longtable,array,tabularx}
\usepackage{enumitem}
\setlist{itemsep=2pt,topsep=4pt,leftmargin=*}
\usepackage[dvipsnames]{xcolor}
\usepackage{fancyhdr}
\usepackage{titlesec}
\usepackage[hidelinks,bookmarksnumbered]{hyperref}
\hypersetup{pdftitle={The Erdos-Straus equation through residual shells: the workbench and the author's notes, their results, with proofs, and their bridges},
  pdfauthor={Claude (Anthropic), model claude-opus-5-5; workbench by The Clankers and contributors; notes by KokunoYumeto},
  pdfsubject={Paper on the Erdos-Straus workbench and the author's notes, version 4},
  pdfkeywords={Erdos-Straus equation, unit fractions, shells, divisor criterion, congruence covering, negative results}}
\usepackage{bookmark}

\titleformat{\section}{\normalfont\Large\bfseries}{\thesection}{0.8em}{}
\titleformat{\subsection}{\normalfont\large\bfseries}{\thesubsection}{0.7em}{}
\titlespacing*{\section}{0pt}{2.2ex plus 1ex}{1.2ex}
\titlespacing*{\subsection}{0pt}{1.8ex plus .8ex}{0.9ex}

\pagestyle{fancy}
\fancyhf{}
\fancyhead[L]{\small\itshape The Erdős--Straus equation through residual shells}
\fancyhead[R]{\small\thepage}
\renewcommand{\headrulewidth}{0.3pt}
\fancypagestyle{plain}{\fancyhf{}\fancyfoot[C]{\small\thepage}\renewcommand{\headrulewidth}{0pt}}

\theoremstyle{plain}
\newtheorem{theorem}{Theorem}[section]
\newtheorem{proposition}[theorem]{Proposition}
\newtheorem{lemma}[theorem]{Lemma}
\newtheorem{corollary}[theorem]{Corollary}
\newtheorem{introthm}{Theorem}
\renewcommand{\theintrothm}{\Roman{introthm}}
\theoremstyle{definition}
\newtheorem{definition}[theorem]{Definition}
\newtheorem{example}[theorem]{Example}
\theoremstyle{remark}
\newtheorem{remark}[theorem]{Remark}

% status line printed under a statement
\newcommand{\status}[1]{\par\smallskip\noindent{\small\color{black!70}\textsf{Status.}\ #1}\par\smallskip}
\newcommand{\negres}[2]{\par\smallskip\noindent\textbf{#1}\ #2\par}

% notation
\newcommand{\B}{\mathcal{B}}
\newcommand{\Zs}{\mathcal{Z}}
\newcommand{\C}{\mathbb{C}}
\newcommand{\R}{\mathbb{R}}
\newcommand{\Q}{\mathbb{Q}}
\newcommand{\ZZ}{\mathbb{Z}}
\newcommand{\NN}{\mathbb{N}}
\newcommand{\Ff}{\mathbb{F}}
\newcommand{\re}{\operatorname{Re}}
\newcommand{\im}{\operatorname{Im}}
\newcommand{\Res}{\operatorname*{Res}}
\newcommand{\sgn}{\operatorname{sgn}}
\newcommand{\Gr}{\operatorname{Gr}}
\newcommand{\cl}{\overline}
\newcommand{\note}[1]{\texttt{#1}}
\newcommand{\lbl}[1]{\textsf{#1}}

\setlength{\parskip}{3pt}
\setlength{\emergencystretch}{2.5em}
\setlength{\parindent}{0pt}

\newcommand{\T}{\mathcal{T}}
\newcommand{\Hs}{\mathcal{H}}
\newcommand{\lcm}{\operatorname{lcm}}
\newcommand{\TV}{\operatorname{TV}}
\newcommand{\tr}{\operatorname{tr}}
\newcommand{\Var}{\operatorname{Var}}


\begin{document}
\thispagestyle{plain}
\begin{center}
{\LARGE\bfseries The Erdős--Straus equation\\[2pt] through residual shells}\\[8pt]
{\Large The workbench and the author's notes: their results,\\ with proofs, and their bridges}\\[14pt]
{\large A paper prepared from an audit of the workbench and the notes\\ by Claude (Anthropic), model \texttt{claude-opus-5-5}}\\[8pt]
{\normalsize Workbench: The Clankers (KokunoYumeto and the contributors credited in the introduction); notes: KokunoYumeto}\\[4pt]
{\normalsize Audit 26--27 September 2026; version~4, 27 September 2026 (the text of version~3 with its framing corrected, and new results; Appendix~\ref{app:changes})}\\[4pt]
{\small Published at the author's request ahead of the author's review; the author may revise or withdraw it.}\\[2pt]
{\small Companion to the paper on the split-zero programme around the Riemann zeta function, version~3, in the Zenodo record with concept DOI \href{https://doi.org/10.5281/zenodo.22678085}{10.5281/zenodo.22678085}}
\end{center}

\begin{abstract}
\noindent
The Erdős--Straus conjecture asserts that $4/n=1/x+1/y+1/z$ has a solution in positive integers for every $n\ge2$. This paper sets out, with proofs, what two bodies of research on the conjecture have established, and adds new results. The first is the Erdős--Straus workbench of The Clankers: a 619-page archive, a supplement of 72 statements on bounded congruence transport, a 124-page reader and a continuation of 28 items. The second is the author's notes of April--June 2026.

The research is organised by the shell criterion, which reduces solvability at a prime $p$ to the occupancy of finitely many shells, each decided by a divisor congruence. On this basis the paper proves exact criteria for the first two shells, a family of residue barriers culminating in the width-one barrier of the notes, and a converse under an exponent condition. It determines the record primes of the least-class function up to $3\cdot10^9$, proves a general fixed-divisor construction of which the archive's four explicit families are instances, and shows that the archive's carrier census can never force the rows that are squares, $1/1024$ of its final domain. From the notes it proves that their terminal core is exactly the set of classes that their certificates cannot reach, and it derives a quadratic-residue sieve on a counterexample whose survivors have density zero.

The research also builds bridges from the equation to other areas, and four of the results of this paper hold in those areas independently of the equation. The notes place the shadow of their order-$7$ mixed mock vector, whose functions agree with Ramanujan's order-$7$ mock theta functions, on three weight-$\tfrac32$ unary theta series; the paper proves that these carry exactly the modular representation of the minimal model $M(2,7)$ after conjugation and a twist by the character of $\eta^3$, and that the analogous order-$5$ block carries the Fibonacci modular data. It classifies the orbits of the workbench's $(3,4,\infty)$ monodromy covers at every level by one invariant. And it shows that the modular flow of the state on $M_2(\C)$ with density proportional to $\operatorname{diag}(e^r,e^{-r})$ and the Lorentz boost of the notes' Lorentzian geometry are the two real forms of one complex one-parameter subgroup of $SL_2(\C)$. Each bridge is described with what is established about it, and the negative results are stated with their exact scope. The workbench does not claim a proof of the conjecture, and nothing here proves or disproves it.
\end{abstract}

\section*{Introduction}\phantomsection\label{sec:about}
\addcontentsline{toc}{section}{Introduction}

\subsection*{The problem and this record}

The conjecture has been verified for all $n\le10^{17}$ \cite{Salez}, it holds for almost all $n$ \cite{Vaughan}, and classical polynomial identities solve it outside the six classes of primes that are squares modulo $840$ \cite{ElsholtzTao,Mordell}. What remains is a question about those classes, and on them polynomial identities provably cannot help (Section~\ref{sec:problem}).

The workbench of The Clankers attacks this question through the shell criterion. Fix a prime $p$ and a candidate least denominator $a$; the rest of the equation is solvable exactly when $a^2$ has a divisor in one of two residue classes modulo $R=4a-p$. The workbench builds on this a large calculus of sufficient conditions, obstructions and residue censuses, and around it geometric and operator-theoretic constructions. The author's notes of April--June 2026 develop the same shell arithmetic in a vocabulary of supports and attach structural layers to it: Niemeier and umbral data, the class $6B$ of the Monster, Ogg's primes, Lorentzian geometry, Cayley--Dickson algebras and mock theta functions.

The workbench and the notes are research carried out with AI systems. The author used them to formulate and test many hypotheses quickly, across several directions at once. Some directions are finished theorems, some are programmes with proved parts, and some are proposals. Not every direction is expected to hold. An unfinished programme is still mathematics, however, and its proved parts stand on their proofs whatever becomes of the rest.
This paper is meant to let a reader see all of this without working through more than a thousand pages. It states what holds, with proofs or the exact location of a proof, what fails and in which scope, and how the research connects with other areas. The audit behind it had limited compute and did not examine every direction. Each direction is described on its own merit and for its bearing on the equation, and where the paper gives a view rather than a proof, the view is labelled as mine and its reasons are given.

\subsection*{Main results}

Throughout, $p$ is an odd prime. For $a>p/4$ with $p\nmid a$ put $R_a=4a-p$ and
\[
E_a=\{u: u\mid a^2,\ R_a\mid 4u+1\},\qquad M_a=\{u: u\mid a^2,\ R_a\mid u+a\};
\]
the shell $a$ is \emph{occupied} if $E_a\cup M_a\ne\emptyset$.

\begin{introthm}[Theorem~\ref{thm:shell}, Propositions~\ref{prop:groupbarrier} and~\ref{prop:groupconverse}]
(i) The equation is solvable at $p$ if and only if some shell $a$ with $p/4<a<p/2$ is occupied, and the ordered pairs $(y,z)$ completing a shell $a$ number exactly $2|E_a|+|M_a|$.
(ii) Let $p\equiv1\pmod4$, $a=(p+R)/4$ with $\gcd(R,pa)=1$, and let $H_a\le(\ZZ/R)^\times$ be the subgroup generated by the primes dividing $a$. If neither $-1$ nor $-4$ lies in $H_a$, the shell $a$ is empty. If $2v_\ell(a)\ge\operatorname{ord}_R(\ell)-1$ for every prime $\ell\mid a$, the shell is occupied if and only if $-1$ or $-4$ lies in $H_a$.
\end{introthm}

Part (i) is the archive's divisor-shell bijection with its layer count, completed here by the range of the least denominator. Part (ii) is the sharpest of a family of residue barriers: the Jacobi-symbol barrier of the archive, a group barrier, and the width-one barrier of the notes. It is a shell form of the quadratic-character mechanism behind the obstruction to polynomial identities, and it says what a proof must produce: for some $R$, a factorisation of $(p+R)/4$ whose prime factors generate $-1$ or $-4$ modulo $R$.

Following Ionascu and Wilson \cite{IonascuWilson}, let $h(p)$ be the least $i$ such that some solution at $p$ has least denominator $x\le(p+4i-1)/4$; a prime is a \emph{strict record} if $h$ is smaller at every smaller prime.

\begin{introthm}[Proposition~\ref{prop:records}]
The strict records of $h$ up to $3\cdot10^9$ are exactly $2$, $73$, $1129$, $1201$, $21{,}169$, $118{,}801$ and $8{,}803{,}369$, with $h=1$, $2$, $3$, $6$, $8$, $15$ and $27$. At $8{,}803{,}369$ the least occupied shell has $R=107$, and exactly one unordered solution has that least denominator.
\end{introthm}

These seven are the records printed by Ionascu and Wilson, recovered in the archive with one class label corrected; that there is no further record up to $3\cdot10^9$ is not in the archive or in their paper, and it was computed by two independent programs. The archive's four explicit families solve the equation on progressions that meet every hard class (Theorem~\ref{thm:families}), and they are instances of a general fixed-divisor construction (Proposition~\ref{prop:fixeddivisor}). Every carrier of the archive's census is of this kind, and none reaches a class that is a square, so the census can never force its square rows, $1/1024$ of its final domain (Proposition~\ref{prop:nosquare}, Corollary~\ref{cor:censuslimit}).

The author's notes contribute a terminal core and a system of certificates. Put $M_{23}=840\cdot11\cdot19\cdot23$, let $S_{840}$ be the six unit squares modulo $840$, and let $Q_q$ be the non-zero squares modulo $q$.

\begin{introthm}[Theorems~\ref{thm:termcore} and~\ref{thm:eightshifts}, Proposition~\ref{prop:densityzero}]
(i) Among the $23{,}760$ classes modulo $M_{23}$ that lie over $S_{840}$ and are units modulo $11$, $19$ and $23$, those on which no elementary certificate of the notes, with residual built from $3,5,7,11,19,23$, applies on any lift are exactly the $2970$ classes of $S_{840}\times Q_{11}\times Q_{19}\times Q_{23}$. The single identities of Salez's seven families leave $3520$ classes open at this level: the core and $550$ further classes.
(ii) A counterexample $p\equiv1\pmod{24}$ is a quadratic residue modulo every odd prime factor of $(p+1)(p+2)(p+3)(p+4)(p+7)(p+8)(2p+1)(3p+1)$.
(iii) The primes $p\equiv1\pmod{24}$ that satisfy the condition from $p+1$ alone have relative density zero among all primes $p\equiv1\pmod{24}$.
\end{introthm}

One half of (i) is proved here and the other is a finite computation, reproduced by an independent program. The conditions from $p+2$, $p+8$ and $2p+1$ are derived here from the notes' certificate branches and are not stated in the archive. At the record prime $8{,}803{,}369$ the other six shifts give no certificate, and those from $p+2$ and $2p+1$ certify it. With certificates of three further shapes, only $23$ primes $p\equiv1\pmod{24}$ below $10^7$ have no certificate of the kinds considered.

The notes attach to the cubic middle of the shell problem an explicit rank-three mixed mock modular object, whose functions agree with Ramanujan's order-$7$ mock theta functions through $q^{40}$ (a reading pass), and they place its shadow on three weight-$\tfrac32$ unary theta series. For $r\in\{1,5,11\}$ let $G_r(\tau)=\sum_{\ell\in\ZZ}v_r(\ell\bmod84)\,\ell\,q^{\ell^2/168}$ be the weight-$\tfrac32$ unary theta series on the notes' three support vectors $v_r$ (Section~\ref{ssec:mock}), let $\rho$ be the pair of matrices by which $(G_1,G_5,G_{11})$ transforms under $S$ and $T$, and let $\rho_{2,7}$ be the representation of $SL_2(\ZZ)$ on the characters of the minimal model $M(2,7)$.

\begin{introthm}[Theorem~\ref{thm:m27}]
There is a signed permutation matrix $P$ with $\overline{\rho(S)}=P\rho_{2,7}(S)P^{-1}$ and $\overline{\rho(T)}\,e(1/8)=P\rho_{2,7}(T)P^{-1}$. So $\bar\rho\otimes\varepsilon^3\cong\rho_{2,7}$, where $\varepsilon$ is the multiplier of Dedekind's $\eta$: the theta series of the notes' shadow, conjugated and twisted by the character of $\eta^3$, transform exactly as the characters of $M(2,7)$.
\end{introthm}

The identity is verified exactly in $\Q(\zeta_{168})$. For order $5$ the same method gives the Fibonacci modular data (Proposition~\ref{prop:fibonacci}). Finally, a package of the workbench builds covers of the projective line from the integral monodromy of the $S^6$ period system; they have a section exactly when the equation is solvable. Write a point of the affine hyperplane $H_D$ of the dual lattice modulo $D$ as $\hat\gamma+b\hat u+c\hat w+d\hat\delta$.

\begin{introthm}[Theorem~\ref{thm:orbits}]
For every level $D\ge1$ the orbits of the monodromy group on $H_D$ are the sets $\{\gcd(b,c,\gcd(D,6))=g\}$, $g\mid\gcd(D,6)$, and they coincide with the orbits of the full stabilizer of the isotropic vector $\gamma$.
\end{introthm}

\subsection*{What the results mean}

My assessment is that, taken together, the results locate the difficulty of the conjecture precisely, for the following reasons. No polynomial identity solves the equation on a residue class that is a square \cite{ElsholtzTao}, and no fixed-divisor certificate reaches such a class (Proposition~\ref{prop:nosquare}). The archive's carrier census runs into this obstruction: its carriers force $97.606\%$ of its final domain and never a square row (Corollary~\ref{cor:censuslimit}). What decides the square classes is the factorisation of the numbers $(p+R)/4$: by Theorem~I(ii) a shell can be occupied only if the prime factors of $(p+R)/4$ generate $-1$ or $-4$ modulo $R$, and under the exponent condition this is also sufficient. A proof of the conjecture therefore has to control, for some $R$, the multiplicative structure of $(p+R)/4$ (Question~\ref{q:global}).

The quadratic-residue sieve shows how restrictive this already is. A counterexample must be a quadratic residue modulo every odd prime factor of eight shifted forms and several more. The primes that satisfy the first condition alone have density zero, and only $23$ primes $p\equiv1\pmod{24}$ below $10^7$ have no certificate of the kinds considered. Whether infinitely many primes have none is Question~\ref{q:qrsurvivors}; a negative answer would prove the conjecture at all but finitely many primes.

The structural results are of a second kind. Theorem~IV places the shadow of the notes' order-$7$ vector inside the modular representation of the $(2,7)$ minimal model. Through the agreement of that vector with Ramanujan's order-$7$ mock theta functions, the first of which gives the $\widehat Z$ invariant of the Brieskorn sphere $-\Sigma(2,3,7)$ \cite{CFS}, the statement concerns these classical objects, and it is independent of the Erdős--Straus problem. A statement about the equation from such structures needs a typed map from the shell data to them; those maps are the subject of Section~\ref{sec:overlaps} and of Questions~\ref{q:modularobject} and~\ref{q:thin}.

\subsection*{The bridges}

In the author's view the connections this research makes between the equation and other areas are among its most valuable outcomes. Section~\ref{sec:overlaps} treats each with the typed map, what is established, what is not verified here, and my assessment. Four results of this paper belong to the other areas themselves and hold independently of the equation: Theorem~IV, Proposition~\ref{prop:fibonacci}, Theorem~V, and Proposition~\ref{prop:flowboost}, by which the modular flow of the state on $M_2(\C)$ with density proportional to $\operatorname{diag}(e^r,e^{-r})$ and the Lorentz boost of the notes' Lorentzian geometry are the compact and the non-compact real forms of one complex one-parameter subgroup of $SL_2(\C)$. A companion note takes up the question this raises for the Apollonian group, the standard thin group acting on the same Minkowski space \cite{ApollonianNote}. Other bridges are partly constructed, such as the split-zero support states and the umbral and Ogg structures of the residue layer, and the bridge to computability through the $\Pi_1$ form of the conjecture is exact. The edges to the author's other workbenches (the zeta programme, $S^6$, Navier--Stokes and Collatz) are listed with their status.

\subsection*{How this record was made}

This record is an audit, carried out on 26--27 September 2026, of the workbench's published documents and of the author's notes of April--June 2026. The workbench is the work of The Clankers, the collective author name of KokunoYumeto and the contributors credited in the workbench's \texttt{README.md} and \texttt{CONTRIBUTIONS.md}:
\begin{itemize}
\item u/CommonCareful3149, for the mod-$107$ arithmetic-progression observation;
\item u/UmbrellaCorp\_HR, for arithmetic, positivity-code and prime-production contributions and for the central-angle and inversive-circle theorems of the archive's §6;
\item u/CivQ17, for the diagrams behind result R05;
\item u/Far-Lifeguard519, for the initiating computational investigation;
\item u/TextBackground496 and u/JGPTech, for source ideas;
\item the owner, KokunoYumeto (u/lepthymo), for the research directions and examples, for directing and developing the investigations with AI, and for the proof and source-fidelity requirements.
\end{itemize}
The archive records language and computational models as contributing-model provenance for particular calculations, not as authors. The notes are the author's.

The audit was carried out by Claude (Anthropic), model \texttt{claude-opus-5-5}, which prepared this paper (\emph{claude-ab} below). Independent Claude instances, each of which wrote none of the text, contributed as follows: a first-pass audit of the workbench; referee passes, which read for overstatement and for understatement and missed generality and found several of the results stated here, each credited at its statement; and six reading passes over the author's notes. Every result found by another instance was verified by claude-ab before it was applied, except the statements marked \emph{(a reading pass)}: those were checked by program in one of the reading passes and not re-derived here, and they count as \emph{not verified here}.

Every statement has one of four statuses:
\begin{itemize}
\item \emph{verified};
\item \emph{not verified here}, which implies nothing in either direction;
\item \emph{my assessment}, a view in the first person with its reasons, introduced by those words;
\item \emph{proved negative}, stated with its counterexample or derivation; the negative results on the workbench's mechanisms are collected in Section~\ref{sec:negative}, and those on the notes are stated where they occur.
\end{itemize}
The status lines under the results make the first two precise. \textsf{Proved here} means the proof is given in full; \textsf{Generalised here} means that the workbench's statement is proved here under weaker hypotheses; \textsf{Audited} means the workbench's proof was read in full, every step checked, and independent checks run; \textsf{Audited (certificate)} means the finite certificate on which a proof rests was recomputed independently and the deduction from it read; \textsf{Reproduced} means the finite content of a statement was recomputed by an independent program and its general proof not re-derived. These are forms of \emph{verified}. \textsf{Audited in part} is verified in the parts named. \textsf{Located} means the statement and its proof were found and are described, and it is a form of \emph{not verified here}. The workbench's own labels (proved, certificate-checked, conjectural, withdrawn) are quoted where they differ.

The workbench was read at commit \texttt{d20b32e} of \texttt{KokunoYumeto/erdos-straus-foundation}, which also holds the 488-page reader of 7 August and its source, and in the Zenodo collection 10.5281/zenodo.20401937, whose record 10.5281/zenodo.22678971 holds the 619-page archive and its source. Locators have the form \emph{archive §$n$, p.~$m$} for the 619-page archive, \emph{supplement §$n$} for the 86-page supplement, and file names and line numbers for the repository; the notes are cited by title and theorem number. Version~3 of this paper and the audit notes cited as \note{NN\_} were withdrawn from the branch \texttt{claude/claude-ab-grind-20260925} on 27 September 2026 because of their framing. They remain, unchanged, in the branch's history at commit \texttt{a90d9e5b}, in the folder \texttt{contrib/claude-ab/grind\_20260925/}, where the folders cited as \note{checks/\ldots} also are. The scripts on which this paper's statements rest are republished in \texttt{contrib/claude-ab/rewrite/checks/es/} on the same branch (Appendix~\ref{app:verification}).

This is version~4. It keeps the theorems, proofs and citations of version~3 and corrects its framing; like the author's latest record, it does not treat the drafts that record supersedes. It adds the results marked as new in version~4: Proposition~\ref{prop:densityzero}, Theorem~\ref{thm:m27}, Propositions~\ref{prop:fibonacci}, \ref{prop:dossier}, \ref{prop:torsor} and~\ref{prop:flowboost}, and Theorem~\ref{thm:orbits} with its corollary. Re-deriving the statements of version~3 from their sources changed some of them. In particular, the identification of the notes' shadow with $M(2,7)$ holds in the form of Theorem~IV, where version~3 said that it fails, and the Lorentz-boost proposition of the notes' Lorentzian section, which version~3 listed as an exception, is correct. Appendix~\ref{app:changes} lists every change.

\subsection*{Organization}

Section~\ref{sec:problem} states the problem and what is known. Section~\ref{sec:corpus} maps the workbench's documents, with the exact scope of the supplement's no-two-shears theorem and the quartic encoding of solutions. Sections~\ref{sec:core} and~\ref{sec:families} contain the core arithmetic and the records, families and census. Section~\ref{sec:items} treats the September continuation and the monodromy covers, and Section~\ref{sec:notes} the author's notes, including the mock theta results. Section~\ref{sec:negative} collects the proved negative results with their scope, Section~\ref{sec:overlaps} the bridges and the directions they open, and Section~\ref{sec:open} the questions and what was not checked. The appendices give the catalogue of statements, the verification, and the changes of version~4.


\tableofcontents
\clearpage

\section{The problem, what is known, and what the workbench claims}\label{sec:problem}

\subsection{The equation}

The Erdős--Straus conjecture asserts that for every integer $n\ge2$ there are positive integers $x,y,z$ with
\[
\frac4n=\frac1x+\frac1y+\frac1z .
\]
Repeated denominators are allowed; for example $4/5=1/2+1/4+1/20$.
It suffices to treat primes: $4/2=1/1+1/2+1/2$, and if $4/p=1/x+1/y+1/z$ then $4/(mp)=1/(mx)+1/(my)+1/(mz)$.

\subsection{What is known}

The following results frame everything below. Each is cited from the literature; none is re-proved here, except where Section~\ref{sec:core} proves a special case.
\begin{itemize}
\item \emph{The classes modulo $840$.} Explicit polynomial identities solve the equation for every $n$ in a primitive class $r\bmod 840$ unless $r$ is a square modulo $840$. So for primes only the classes
\[
p\equiv1,\,121,\,169,\,289,\,361,\,529\pmod{840}
\]
remain. Elsholtz and Tao attribute this to Obláth and to Mordell \cite{ElsholtzTao,Mordell}; Proposition~\ref{prop:mod840} below proves the reduction for primes.
\item \emph{Why identities do not cover the square classes.} Elsholtz and Tao \cite{ElsholtzTao} call a primitive class $n\equiv r\pmod q$ \emph{solvable by polynomials} if there are polynomials $P_1,P_2,P_3$ which take positive integer values for all sufficiently large $n$ in the class and satisfy $4/n=1/P_1(n)+1/P_2(n)+1/P_3(n)$. A primitive class with $r$ a square modulo $q$ is not solvable by polynomials. They cite Mordell and Schinzel for this, and trace the reciprocity argument behind it to Schinzel and to Yamamoto (\cite{ElsholtzTao}, §1).
\begin{itemize}
\item Section~\ref{ssec:jacobi} gives the workbench's shell form of this obstruction, with proof.
\item Proposition~\ref{prop:nosquare} gives a form that applies to the carriers of the archive's census.
\end{itemize}
\item \emph{Density.} The number of $n<N$ for which the equation has no solution is at most $N\exp(-c\log^{2/3}N)$ for an absolute $c>0$ and $N$ large (Vaughan \cite{Vaughan}, as stated by Elsholtz and Tao \cite{ElsholtzTao}).
\item \emph{Counting.} Elsholtz and Tao bound the number $f(n)$ of solutions and its averages, and classify solutions at primes as Type~I ($p$ divides exactly one denominator) or Type~II ($p$ divides two) \cite{ElsholtzTao}.
\item \emph{Verification.} The conjecture has been checked for all $n\le10^{14}$ (Swett, cited in \cite{ElsholtzTao}) and for all $n\le10^{17}$ (Salez \cite{Salez}); a 2025 preprint reports the bound $10^{18}$ \cite{MihneaDumitru}.
\end{itemize}

\subsection{What the workbench claims}

The repository's README opens by saying that the project claims neither a proof nor a counterexample. The research state calls it an ongoing programme, not a completed proof. The archive's status section (p.~4) names the remaining global obligation: for every relevant prime, at least one marked shell class must be shown to be occupied.
In the parts read here, the workbench's scope statements match what its proofs establish. Where a statement needed a correction, this paper gives the correction at the statement (Section~\ref{sec:negative} and Appendix~\ref{app:catalogue}).

\subsection{The approach in one paragraph}

Fix a prime $p$ and a candidate least denominator $a$ with $p/4<a<p$. The remaining equation $1/y+1/z=4/p-1/a$ has solutions exactly when $a^2$ has a divisor in one of two prescribed residue classes modulo $R=4a-p$ (Theorem~\ref{thm:shell}). The workbench calls $a$ a \emph{shell} and $R$ its \emph{residual}. For $p\equiv1\pmod4$ the residual is $\equiv3\pmod4$, and the shells in play are $R=3,7,11,\dots$ up to about $p$.
Most of the workbench's arithmetic is organised around three kinds of tool:
\begin{itemize}
\item sufficient conditions under which a shell is occupied, stated as congruence conditions on $p$ or on the prime factors of $a$;
\item obstructions under which a shell is empty;
\item bounded \emph{carrier} censuses, which count how much of a finite space of residue data these conditions cover.
\end{itemize}
Around this arithmetic core the documents also develop geometric and operator-theoretic constructions (inversive circles, Lorentz lifts, Boolean carriers, an auxiliary torus, a Busy-Beaver map). Section~\ref{sec:corpus} places each of them in the documents, and Section~\ref{sec:overlaps} describes the bridges they open.

\section{The documents, and what each contains}\label{sec:corpus}

\subsection{The editions and how they relate}

\begin{center}\small
\begin{tabularx}{\linewidth}{@{}>{\raggedright\arraybackslash}p{0.19\linewidth}>{\raggedright\arraybackslash}p{0.23\linewidth}>{\raggedright\arraybackslash}X@{}}
\toprule
Edition & Where & Contents and relation to the others\\
\midrule
Project reader, 488 pages, 7 August 2026 & repository root, with source and Lean archives & The earlier edition. The README says it is preserved and is not the newest reader.\\
Cumulative archive, 619 pages; title page dated 30 August 2026 & Zenodo 22666493 and 22678971, file 00 & The 488-page reader continued. Its main source has the same section skeleton, now 60{,}561 lines instead of 57{,}452. Fourteen satellite files add §§10--19, the $Q=83$ step of §24 (Theorem~24.241) and the carrier refinements §§24.1--24.3. It begins from the source of the earlier Zenodo record 21845035, version~69 (archive p.~3). It is the front file of the Zenodo collection.\\
Bounded congruence transport, 67 pages, 8 September & Zenodo 22666493, files 02--03 & The 58-statement supplement.\\
Expanded supplement, 86 pages, 9 September & Zenodo 22678971, files 07--08; repository \texttt{research/\allowbreak expanded-2026-09-08/} & The same, with 14 further statements: 72 statements in 15 source files, with checkers (§\ref{ssec:supplement}).\\
ES/Fable reader, 124 pages, 20 September & repository \texttt{research/\allowbreak incoming/\allowbreak es-fable-zeta-bridge-\allowbreak 20260920/} & A coordinate bridge from Erdős--Straus solutions to the Fable cover and a weighted conductor of the zeta programme (§\ref{ssec:crosswalk}).\\
Continuation, 17--23 September & repository README, \emph{Start reading} (28 items); \texttt{research/\allowbreak incoming/} & 41 dated packages; the results bench R01--R10 and the integration arguments X1--X7 (§\ref{ssec:continuation}).\\
\bottomrule
\end{tabularx}
\end{center}

These editions have different contents; the repository preserves them rather than overwriting them (README, edition table).

\subsection{The 619-page archive, section by section}\label{ssec:archive}

Page numbers are those of the PDF. The column \emph{Kind} says how a section relates to the equation:
\begin{itemize}
\item \emph{direct}: it proves or tests solvability;
\item \emph{literature}: it audits a published approach to the equation;
\item \emph{structural}: a construction whose connection to solvability is, in the archive's words (p.~4), an interface still to be supplied by an explicit typed map. Section~\ref{sec:overlaps} describes these connections.
\end{itemize}
The status column refers to this paper.

\begin{center}\small
\begin{longtable}{@{}>{\raggedright\arraybackslash}p{0.06\linewidth}>{\raggedright\arraybackslash}p{0.08\linewidth}>{\raggedright\arraybackslash}p{0.45\linewidth}>{\raggedright\arraybackslash}p{0.13\linewidth}>{\raggedright\arraybackslash}p{0.17\linewidth}@{}}
\toprule
§ & pp. & Content & Kind & Status here\\
\midrule
\endhead
front & 3--5 & Research status: the census figures of §24, the four families, the record-shell results, and an explicit list of non-claims & summary & read\\
1 & 6--16 & Scope, road map, mathematical status, attribution, lineage & summary & read in part\\
2 & 16--18 & The discriminant-five identity $(2n+1)^2-5=4(n^2+n-1)$ and its root involution, from u/CivQ17 & structural & located\\
3 & 18--20 & The shell $p_*=8{,}803{,}369$, $R_*=107$, reached through Busy-Beaver values ($S(4)=107$, $\mathrm{Space}(5)=12289$) & direct (one shell) & $R=107$ is the least occupied shell at $p_*$ (Proposition~\ref{prop:records})\\
4 & 20--22 & Alpöge's constant-Jacobian map $F:\C^3\to\C^3$ with a three-point fibre & structural & see \S\ref{sec:overlaps}\\
5 & 22 & Withdrawn constructions (a cube-root interpolation on the Fable fibre), with the point of failure & structural & read\\
6 & 22--27 & Central-angle and inversive-circle theorems of u/UmbrellaCorp\_HR; a Lorentz reflection & structural & located\\
7 & 27--31 & The complete fibre of $F$ over $(-\tfrac14,0,0)$ and a marked-factor lift & structural & see \S\ref{sec:overlaps}\\
8 & 31--38 & The singular $C_1$--$C_2$--$C_3$ support object; dihedral actions & structural & located\\
9 & 38--93 & The exact divisor-shell bijection (Theorem~9.1, for every odd prime) and its layer count (Theorem~9.64, p.~90); then directed blade transport and a quarter-radius inversion & direct (Theorems 9.1 and 9.64), structural (the rest) & the bijection and the count re-proved here (Theorem~\ref{thm:shell}(b))\\
10 & 93--98 & Residual sieve modulo $840$: an occupied shell with $R=3$ or $R=7$ at every prime $p\equiv1\pmod{24}$ outside the six hard classes (Theorem~10.6); the Jacobi-symbol barrier (Theorem~10.9) & direct & proved here (Propositions~\ref{prop:survivors}, \ref{prop:jacobi})\\
11 & 98--100 & Elsholtz--Tao Type~I and Type~II coordinates mapped to the residual shell & literature & located\\
12 & 100--106 & Adjacent programmes (López, Dyachenko, Bradford): what each proves, and where a claimed bridge fails & literature & located\\
13 & 106--114 & Mordell's constructions and Rosati's theorem; Salez's seven charts; Bright--Loughran surfaces; the Vaughan boundary & literature & located\\
14 & 114--119 & López's structural solutions: exact divisor coordinates and a quantifier repair & literature & located\\
15 & 119--124 & The Star--Kneser programme: a one-way forcing theorem and a retracted converse & direct & located\\
16 & 124--127 & The mod-$107$ deficit progression, following u/CommonCareful3149 & direct & progression and its corollary reproduced (§\ref{ssec:records})\\
17 & 127--135 & Two-fraction criteria; Monks--Velingker; Ionascu--Wilson classes, with the printed record list recovered and one printed class label corrected (Proposition~17.6); Zelator's prime-square classification & literature & records recomputed (§\ref{ssec:records})\\
18 & 135--148 & The residual-eleven support ideal, its colons, and the paired $R=7/R=11$ affine theorem & direct & located\\
19 & 148--157 & Strict record shells: an open conjecture and the boundary example $p=806{,}521$ & direct & records and deficits reproduced (§\ref{ssec:records})\\
20 & 157--179 & The shell $p=12289$, $R=11$, $a=3075$: counts $(3,6,3)$, its dihedral orbit, and the Star--Kneser quotient certificate & direct (one shell) & counts reproduced (below the table); $a$ is the least occupied shell; the rest located\\
21 & 179--200 & The split-zero semiring $G(R)=R\sqcup\{\tau\}$ and the $C_1$ carrier & structural & see \S\ref{sec:overlaps}\\
22 & 200--201 & Rigidified denominator towers over $\Ff_1$ & structural & located\\
23 & 201--250 & Mirror completion of the Boolean carrier; the $12289$ orbit matrix & structural & located\\
24 (i) & 250--298 & The opening of §24: Tomita--Takesaki transport and $C_3$ modular data (Theorems 24.1--24.44) & structural & located\\
24 (ii) & 299--604 & The rest of §24, from Theorem~24.45: shell arithmetic.
\begin{itemize}
\item Star--Kneser forcing (Theorems 24.45--24.48, pp.~299--301);
\item endpoint divisor fibres (Theorems 24.50--24.63, pp.~303--312);
\item coupled and adjacent shells, and the base of the census (Theorem~24.100, p.~335);
\item the carriers for $R=11$ to $R=31$, including the constant-factor carriers (Lemma~24.218, p.~422) and the four families (Theorem~24.219, p.~423);
\item the $R=71$, $R=191$ and $Q=83$ steps (Theorems 24.233, 24.237 and 24.241, pp.~436--444);
\item §§24.1--24.3 (pp.~451--477): the $R=167$ and $R=239$ carrier theorems and the refresh of the $R=167$ coefficient;
\item §24.4 (pp.~478--604): a return to the $R=23$ lineage, interleaved with further structural constructions.
\end{itemize}
& direct (with structural material in §24.4) & four families proved here (Theorem~\ref{thm:families}); census located (§\ref{ssec:census})\\
25 & 605--616 & A map from Busy-Beaver extrema to finite shell carriers & structural & see \S\ref{sec:overlaps}\\
\bottomrule
\end{longtable}
\end{center}

At $p=12289$ the shell $a=3075$ has
\[
E_a=\{41,1845,15375\},\qquad M_a=\{5,225,1875,5043,42025,1891125\}.
\]
So $(|E_a|,|M_a|,|E_a|)=(3,6,3)$, and exactly six unordered solutions have least denominator $3075$ (recomputed here).

So about 250 of the 619 pages treat solvability directly:
\begin{itemize}
\item §§3, 9 (Theorems 9.1 and 9.64), 10, 15, 16 and 18--20, about 70 pages;
\item §24 from p.~299 to the end of §24.3 (p.~477), 179 pages.
\end{itemize}
Parts of the 127 pages of §24.4 also do. The literature audits of §§11--14 and 17 add about 35 pages. The rest are structural constructions, whose connection to the equation the archive describes as an interface still to be supplied by an explicit typed map (p.~4).

\subsection{The bounded-transport supplement (72 statements)}\label{ssec:supplement}

The supplement starts from the primes $p=12h+1$ and all shells $3h+1\le a\le9h$ of Theorem~\ref{thm:shell}. Its sections are:
\begin{enumerate}[label=\arabic*.]
\item original shells, coordinates and integer reconstruction, with the first-two-shell sieve and an all-shell no-hit sieve;
\item local congruences with all integral lifts retained, with an integral cyclic bridge;
\item bounded divisibility and simultaneous congruences: unary Boolean transport, a shared-variable Chinese-remainder construction credited to a colleague, raw-scale hit incidence, and a boundary-swap completion;
\item exact cross-shell transport at a fixed rational ratio;
\item a primitive Euclidean successor with its full arithmetic domain;
\item arithmetic localisation along the original norm subdivision, with a reading of Connes' primitive intersections;
\item a strictly decreasing denominator with exact arithmetic return;
\item the Navier--Stokes auxiliary cover: the torus matrix $J=\begin{pmatrix}3&1\\1&5\end{pmatrix}$ of the announced Navier--Stokes manuscript, its covering fibres and character corrections, and sampling and reconstruction estimates.
\end{enumerate}
The supplement's own README states that its results prove or disprove neither the Erdős--Straus conjecture nor the Riemann hypothesis nor the Yang--Mills mass gap, and that they do not validate the Navier--Stokes construction.

\status{Reproduced in part. A second-pass audit, which stopped before it wrote its report, recomputed the finite content of the statements of Sections 1--7. Its scripts, which are independent of the workbench's code, found every checked statement correct:
\begin{itemize}
\item the gluing and Chinese-remainder theorems on random systems and on all shells of small primes;
\item the ratio-transport domain on all $57{,}882$ ordered shell pairs of the primes $p\equiv1\pmod{12}$ below $260$;
\item the shear and no-two-shears theorems on all shells of $131$ primes below $4000$;
\item the Connes marked successor on all primes $p\equiv1\pmod{12}$ below $3\cdot10^5$;
\item the fixed-$y$ successor on all $3{,}114$ source witnesses below $1500$;
\item the orchard partition on all $42$ admissible triples below $1300$;
\item the boundary four-candidate law on all $72{,}282$ raw records below $400$;
\item the unary box, the input/output incidence and the raw-scale theorems on their stated ranges.
\end{itemize}
One of its checks, on the smallest member of a boundary orbit, reported mismatches. The script expected the pattern $(\min(A,B),\max(A,B))$ in every case. The supplement's own table (\texttt{boundary\_swap\_completion.tex}, lines 296--309) gives a one-point orbit's own point as its smallest member, which is what the script found. The mismatch is therefore in the script, not in the supplement. The same pass also looked at scope, and its search led to Proposition~\ref{prop:twoshears} below. The general proofs were not re-derived, except that of Section~5's no-two-shears theorem; the auxiliary-torus Section 8 was not checked. The scripts are in \note{checks/es\_reader/second\_pass\_partial/}.}

Section 5 of the supplement studies the two unimodular shears $L(m,n)=(m+n,n)$ and $J(m,n)=(m,m+n)$ as maps between adjacent shells. It encodes the solutions with first denominator $a$ by the set
\[
 \mathcal W_a=\{(m,n)\in\ZZ_{>0}^2:\ m\mid pa,\ n\mid pa,\ \gcd(m,n)=1,\ R_a\mid m+n\},
\]
which is in bijection with the ordered solutions $(a,y,z)$ through $y=pak/n$, $z=pak/m$, $k=(m+n)/R_a$ (supplement §1, first lemma; this is Theorem~\ref{thm:shell}(b) in other coordinates). A \emph{shear edge} is a step $(a,m,n)\mapsto(a+1,L(m,n))$ or $(a,m,n)\mapsto(a+1,J(m,n))$ with $(m,n)\in\mathcal W_a$ and the image in $\mathcal W_{a+1}$. The supplement proves, for $p=12h+1$ and shells $3h+1\le a\le9h$, that no two shear edges compose (Theorem \emph{No two consecutive positive Euclidean shell steps}, \texttt{publication.tex}, line~933). Its proof uses the hypothesis on $p$ at one step only, and there only through $p\equiv1\pmod3$.

\begin{proposition}[the no-two-shears theorem, exact scope]\label{prop:twoshears}
Let $p$ be a prime, and consider all shells $a$ with $p/4<a<p$.
\begin{enumerate}[label=(\alph*)]
\item If $p\not\equiv2\pmod3$, no two shear edges compose: there is no path $(a,m,n)\mapsto(a+1,m',n')\mapsto(a+2,m'',n'')$ of shear edges.
\item Such paths exist at some primes $p\equiv2\pmod3$. At $p=131$ one is $(33,1,33)\mapsto(34,1,34)\mapsto(35,1,35)$, with $R_a=1,5,9$. At $p=1613$, where $p\equiv1\pmod4$, one is $(404,1,404)\mapsto(405,1,405)\mapsto(406,1,406)$, with $R_a=3,7,11$. Proposition~\ref{prop:sheargraph} says exactly at which primes.
\end{enumerate}
\end{proposition}
\status{\textsf{Generalised here}, with the workbench's proof. The supplement states (a) for $p\equiv1\pmod{12}$ and $3h+1\le a\le9h$ only. Checked on the full graph of every prime below $2000$ (script \note{esr\_twoshears.py}: no path for $p\equiv1\pmod3$; $14$ paths for $p\equiv2\pmod3$), and by a complete search up to $p\le10^7$ that uses the shape of an edge proved below: it finds $19{,}558$ paths of two $L$-edges, all at primes $p\equiv2\pmod3$, and as many of two $J$-edges.}
\begin{proof}
\emph{The shape of one edge} (the supplement's classification, whose proof needs only $a+1<p$). Let $(a,m,n)\mapsto(a+1,m+n,n)$ be an $L$-edge. Then $n$ divides $pa$ and $p(a+1)$, whose greatest common divisor is $p$. If $n=p$, then $p\nmid m$ by coprimality, and $m+p$ divides $p(a+1)$ and is prime to $p$, so $m+p\mid a+1$; but $m+p>p>a+1$. So $n=1$. If $p\mid m$, then $m+1$ is prime to $p$ and exceeds $a+1$, and yet divides $p(a+1)$, which is impossible; so $p\nmid m$ and $m\mid a$. Then $m+1\le a+1<p$, so $m+1\mid a+1$. Finally $R_a\mid m+1$ and $R_{a+1}=R_a+4\mid m+2$. For a $J$-edge, exchange $m$ and $n$; the exchange preserves every $\mathcal W_a$.

(a) An $L$-edge ends at $(a+1,m+1,1)$, whose first coordinate is at least $2$; a $J$-edge starts at a pair whose first coordinate is $1$. So an $L$-edge is never followed by a $J$-edge, and by the exchange a $J$-edge is never followed by an $L$-edge. Suppose two $L$-edges pass through $(a+i,m+i,1)$ for $i=0,1,2$. By the shape of each edge, $m+i\mid a+i$ and $R_a+4i\mid m+i+1$ for $i=0,1,2$. Put $K=a-m$. Then $m+i\mid K$ for each $i$, and one of $m,m+1,m+2$ is divisible by $3$, so $3\mid K$. Also $R_a+4i$ divides $4(m+i+1)-(R_a+4i)=4m+4-R_a=:C$, and one of $R_a,R_a+4,R_a+8$ is divisible by $3$, so $3\mid C$. But $C=4m+4-4a+p=p+4-4K\equiv p+1\pmod3$. Hence $p\equiv2\pmod3$. Two $J$-edges are excluded by the exchange.

(b) At $p=131$: $R_{33}=1$ divides $1+33$; $R_{34}=5$ divides $35$; $R_{35}=9$ divides $36$; each coordinate divides $pa$, and each pair is coprime. At $p=1613$ (prime): $3\mid405$, $7\mid406=7\cdot58$ and $11\mid407=11\cdot37$. Each step is $J(1,n)=(1,n+1)$.
\end{proof}
So the hypothesis $p\equiv1\pmod{12}$ of the supplement can be weakened to $p\not\equiv2\pmod3$. The next proposition shows that no weaker congruence hypothesis suffices, although the conclusion fails only at a minority of the primes $p\equiv2\pmod3$ (Remark~\ref{rem:sheardensity}). The primes $p\equiv2\pmod3$ are solvable in any case (Section~\ref{sec:core}); both propositions concern the transport graph, not solvability.

\begin{proposition}[the shear graph]\label{prop:sheargraph}
Let $p$ be a prime. Consider the directed graph whose vertices are the triples $(a,m,n)$ with $p/4<a<p$ and $(m,n)\in\mathcal W_a$, and whose arrows are the shear edges.
\begin{enumerate}[label=(\alph*)]
\item Every vertex has at most one outgoing and at most one incoming edge. So every connected component is a directed path.
\item \emph{Diagonal paths.} Let $R=4a-p$, and let $k\ge1$ with $a+k<p$. The vertices $(a+i,a+i,1)$, $0\le i\le k$, form a path of $L$-edges if and only if $R+4i$ divides $p+4$ for every $0\le i\le k$.
\item If $(a+i,m+i,1)$, $0\le i\le k$, is a path of $L$-edges and $R=4a-p$, then $R+4i$ divides $p+4$ for every $0\le i\le k-1$. In particular a path of two edges forces $R(R+4)\mid p+4$.
\item Every prime $p\equiv131\pmod{180}$ has a path of two edges. More generally, let $q\ge1$ and let $c$ be prime to $q$, with $c\equiv2\pmod3$ if $3\mid q$ (these are exactly the classes that contain primes $\equiv2\pmod3$). Then infinitely many primes $p\equiv c\pmod q$ have a path of two edges. So if Proposition~\ref{prop:twoshears}(a) holds at all but finitely many primes of a union of classes modulo $q$, then none of those classes contains primes $\equiv2\pmod3$.
\item For every $k\ge1$, infinitely many primes have a path of length $k$.
\end{enumerate}
\end{proposition}
\status{Proved here. Found by the referee pass of version~1; each part checked here. Part (a) was confirmed on the full graph of every prime below $1200$. A complete enumeration below $10^7$ through (c) and the shape of an edge finds:
\begin{itemize}
\item $19{,}558$ two-edge paths at $17{,}425$ primes;
\item $R(R+4)\mid p+4$ on every one of those paths;
\item all $13{,}822$ primes $p\equiv131\pmod{180}$ among the primes with a path.
\end{itemize}
The least primes with a path of length $1$, $2$, $3$, $4$ and $5$ are $11$, $131$, $3461$, $109{,}391$ and $1{,}183{,}451$. The referee gave $3{,}533{,}141$ for length $5$; it has such a path, but it is not the least.}
\begin{proof}
(a) By the shape of an edge, an $L$-edge starts at a vertex with $n=1$ and a $J$-edge at one with $m=1$. So the only possible successors of $(a,m,n)$ are $(a+1,m+1,1)$, when $n=1$, and $(a+1,1,n+1)$, when $m=1$. Both are possible only if $m=n=1$. Then $(1,1)\in\mathcal W_a$ forces $R_a\mid2$, so $R_a=1$. Either successor would then need $R_{a+1}=5$ to divide $3$, which is false.
The only possible predecessors of $(a+1,m',n')$ are $(a,m'-1,1)$, by an $L$-edge when $n'=1$, and $(a,1,n'-1)$, by a $J$-edge when $m'=1$. Both are possible only if $m'=n'=1$, and then $m'-1=0$ is not a positive coordinate.
Finally, $a$ increases along edges, so there are no cycles.

(b) The triple $(a+i,a+i,1)$ satisfies the divisibility and coprimality conditions of $\mathcal W_{a+i}$. It lies in $\mathcal W_{a+i}$ exactly when $R+4i\mid a+i+1$. As $R+4i$ is odd and $4(a+i+1)=(R+4i)+(p+4)$, this holds exactly when $R+4i\mid p+4$. Consecutive vertices are joined by an $L$-edge, since $L(a+i,1)=(a+i+1,1)$.

(c) Let $0\le i\le k-1$ and put $K=a-m$. The edge from the $i$-th vertex gives $m+i+1\mid a+i+1$, hence $m+i+1\mid K$. The $i$-th vertex lies in $\mathcal W_{a+i}$, so $R+4i\mid m+i+1$. Hence $R+4i\mid K$. Also $R+4i$ divides $4(m+i+1)-(R+4i)=4m+4-R$. Adding $4K$, it divides $4a+4-R=p+4$. For two edges, $R$ and $R+4$ are odd with difference $4$, hence coprime, so $R(R+4)\mid p+4$.

(d) If $p\equiv131\pmod{180}$, then $p\equiv3\pmod4$, so $R=1$ occurs, at $a=(p+1)/4$. Also $45\mid p+4$, so $1$, $5$ and $9$ divide $p+4$, and (b) applies with $k=2$.
In general, choose an odd $R$ with the following properties:
\begin{itemize}
\item $R\equiv-c\pmod4$ if $4\mid q$;
\item no prime $\ell\ge5$ dividing $q$ divides $R(R+4)(R+8)$;
\item the one of $R,R+4,R+8$ that is divisible by $3$ is not divisible by $9$.
\end{itemize}
Such an $R$ exists by the Chinese remainder theorem: each prime $\ell\ge5$ excludes three classes, and modulo $9$ the classes $0$, $1$ and $5$ are excluded.
Put $m'=R(R+4)(R+8)$. Its three factors are odd with pairwise differences $4$ and $8$, hence pairwise coprime, and $\gcd(m',q)\mid3$.
The congruences $p\equiv c\pmod q$, $p\equiv-4\pmod{m'}$ and $p\equiv-R\pmod4$ are compatible: modulo $3$ both of the first two give $p\equiv2$, and modulo $4$ the first and the third agree by the choice of $R$. Together they define a class prime to its modulus, which contains infinitely many primes by Dirichlet's theorem. Each such prime $p>R+8$ has the diagonal path of (b) with $a=(p+R)/4$ and $k=2$.
If $3\nmid q$, every class prime to $q$ contains primes $\equiv2\pmod3$ (by the Chinese remainder theorem and Dirichlet's theorem). If $3\mid q$, a class contains such primes exactly when $c\equiv2\pmod3$.

(e) Take $R=1$. The modulus $\operatorname{lcm}(1,5,9,\dots,4k+1)$ is odd, so the congruences $p\equiv3\pmod4$ and $p\equiv-4$ modulo it define a class prime to its modulus. Each prime $p>(4k+1)/3$ in it has, by (b), a diagonal path of length $k$ from $a=(p+1)/4$.
\end{proof}

\begin{remark}[how often the theorem fails]\label{rem:sheardensity}
Among the primes $p\equiv2\pmod3$, those with a path of two edges have upper relative density at most $0.335$. So the conclusion of Proposition~\ref{prop:twoshears}(a) holds on a set of primes $p\equiv2\pmod3$ of lower relative density at least $0.665$. Below $10^7$ it fails at $17{,}425$ of the $332{,}384$ primes $p\equiv2\pmod3$ ($5.24\%$); the smallest primes $p\equiv2\pmod3$ at which it holds are $2,5,11,17,23,29,41,47,53,59,71,83,\dots$.
\end{remark}
\begin{proof}
By Proposition~\ref{prop:twoshears}(a) and Proposition~\ref{prop:sheargraph}(c), a prime with a path of two edges satisfies:
\begin{itemize}
\item $p\equiv2\pmod3$;
\item for some odd $R$ with $R(R+4)\le p+4$, both $p\equiv-4\pmod{R(R+4)}$ and $p\equiv-R\pmod4$.
\end{itemize}
Fix $R$ and put $m=R(R+4)$. By Dirichlet's theorem, the second condition picks out a set of relative density $\delta_R=(1+[3\mid m])/(2\varphi(m))$ among the primes $p\equiv2\pmod3$.
For large $R$ the bound must be uniform in $x$:
\begin{itemize}
\item \emph{For $4m\le x^{1/2}$.} The Brun--Titchmarsh inequality $\pi(x;Q,b)\le2x/(\varphi(Q)\log(x/Q))$ for $Q<x$ \cite{MontgomeryVaughan}, with $Q=4m$, bounds the number of such primes up to $x$ by $2x/(\varphi(m)\log x)$. This is $(4+o(1))/\varphi(m)$ times the number of primes $p\equiv2\pmod3$ up to $x$.
\item \emph{For $4m>x^{1/2}$.} Each class has at most $x/m+1$ members up to $x$, and $R\le\sqrt{x+4}$. The total over these $R$ is $O(x^{3/4})$.
\end{itemize}
Since $\sum_R1/\varphi(R(R+4))$ converges, letting $x\to\infty$ and then the cut between small and large $R$ tend to infinity bounds the upper relative density by $\sum_{R\text{ odd}}\delta_R$.
Numerically, the partial sum over $R<2\cdot10^5$ is $0.334275$. The rest is less than $2\cdot10^{-4}$, by $n/\varphi(n)<e^\gamma\log\log n+3/\log\log n$ for $n\ge3$ \cite{RosserSchoenfeld}. So the sum is less than $0.335$.
\end{proof}

\subsection{The 124-page ES/Fable reader}\label{ssec:crosswalk}

It encodes a solution of $4/p=1/x+1/y+1/z$ as the binary quartic
\[
H_{\mathrm{ES}}(U,V)=-S^{-1}(U-pV)(U-xV)(U-yV)(U-zV),\qquad S=p+x+y+z,
\]
with normalised coefficients $u_0=-1/S$, $u_1=1$, $u_2=-\bigl(p(x+y+z)+xy+yz+zx\bigr)/S$, $u_3=\bigl(p(xy+yz+zx)+xyz\bigr)/S$ (which equals $5xyz/S$ on solutions) and $u_4=-pxyz/S$. It then transports the quartic to the Fable cover and to a weighted conductor of the zeta programme.

\begin{proposition}\label{prop:quartic}
\begin{enumerate}[label=(\alph*)]
\item On solutions, $p=-5u_4/u_3$ and
\[
625u_0u_4^3-125u_3u_4^2+25u_2u_3^2u_4-4u_3^4=0 .
\]
\item Conversely, let $u_0,u_2,u_3,u_4$ be complex numbers with $u_0u_3u_4\ne0$ that satisfy this relation. Put
\[
S=-1/u_0,\quad p=-5u_4/u_3,\quad P=u_3S/5,\quad \sigma_1=S-p,\quad \sigma_2=-u_2S-p\sigma_1 .
\]
Then the roots $x,y,z$ of $T^3-\sigma_1T^2+\sigma_2T-P$ are non-zero and satisfy $4/p=1/x+1/y+1/z$. The quartic $-S^{-1}(U-pV)(U-xV)(U-yV)(U-zV)$ has the coefficients $u_0,1,u_2,u_3,u_4$.
\end{enumerate}
\end{proposition}
\status{Proved here. The README of the package states (a), and \note{21\_} (Proposition~21.6) also checked the relation. Part (b) was found by the referee pass of version~1 and is proved here.}
\begin{proof}
(a) Write $\sigma_1=x+y+z$, $\sigma_2=xy+yz+zx$ and $P=xyz$. The equation is $4P=p\sigma_2$, so the third elementary symmetric function of $p,x,y,z$ is $p\sigma_2+P=5P$; this gives $u_3=5P/S$, and $p=-5u_4/u_3$ follows. Substituting and multiplying by $S^4/625$, the left side becomes $P^3\bigl(p^3-p^2S+p(p\sigma_1+\sigma_2)-4P\bigr)=P^3(p\sigma_2-4P)=0$.

(b) With the stated definitions, $u_0=-1/S$, $u_4=-pu_3/5$, $u_2=-(p\sigma_1+\sigma_2)/S$ and $u_3=5P/S$. Substituting, and using $S=p+\sigma_1$, gives
\[
625u_0u_4^3-125u_3u_4^2+25u_2u_3^2u_4-4u_3^4=\frac{u_3^3}{S}\bigl(5p^3-5p^2S+5p(p\sigma_1+\sigma_2)-20P\bigr)=\frac{5u_3^3}{S}\,(p\sigma_2-4P).
\]
So the relation gives $p\sigma_2=4P$. Now $xyz=P\ne0$, since $u_3\ne0$, and $p\ne0$, since $u_4\ne0$. Hence $1/x+1/y+1/z=\sigma_2/P=4/p$.
For the quartic, the coefficients of $U^4$, $U^3V$, $U^2V^2$, $UV^3$ and $V^4$ in $-S^{-1}(U-pV)(U-xV)(U-yV)(U-zV)$ are:
\begin{itemize}
\item $-1/S=u_0$;
\item $(p+\sigma_1)/S=1$;
\item $-(p\sigma_1+\sigma_2)/S=u_2$;
\item $(p\sigma_2+P)/S=5P/S=u_3$;
\item $-pP/S=u_4$.
\end{itemize}
\end{proof}
So the relation is the equation itself, written in these coordinates: by (a) and (b), it cuts out exactly the quartics that come from triples, possibly complex, with $4/p=1/x+1/y+1/z$. The Fable and conductor maps built on it are \textsf{located}; their comparison with the zeta programme is in the zeta audit (\note{47\_}, overlap OV-13).

\subsection{The continuation of 17--23 September}\label{ssec:continuation}

The README's \emph{Start reading} list has 28 items with proof packages in \texttt{research/incoming/}. The results bench R01--R10 and the integration arguments X1--X7 accompany them.
\begin{itemize}
\item The first-pass audit \note{43a\_} lists them all.
\item Section~\ref{sec:items} states the eleven items audited there.
\item Section~\ref{sec:negative} gives their negative results.
\end{itemize}

\subsection{The author's notes}\label{ssec:notescorpus}

The author's notes of April--June 2026 are separate from the workbench's documents, and Section~\ref{sec:notes} treats them. Their arithmetic is the shell calculus of Section~\ref{sec:core} together with the classical identity families. Their structural layers are Niemeier and umbral data, the class $6B$, Ogg's primes, Lorentzian geometry, Cayley--Dickson algebras and mock theta functions. The archive's §25 is drawn from one of them.


\section{The core arithmetic, with proofs}\label{sec:core}

Throughout, $p$ is an odd prime and a \emph{solution} is a triple of positive integers with $4/p=1/x+1/y+1/z$. The workbench mostly works with $p=12h+1$ ($h\ge1$). The results of this section, except Propositions~\ref{prop:jacobi}, \ref{prop:groupbarrier} and~\ref{prop:groupconverse}, were first stated in the reader of four workbenches \cite{WorkbenchReader}, whose Erdős--Straus section was refereed three times (Appendix~\ref{app:verification}); the audit note is \note{43\_} and the first-pass report \note{43a\_}.

\subsection{The shell criterion}

For an integer $a>p/4$ not divisible by $p$ put $R_a=4a-p$, $S_a=pa$ and
\[
E_a=\{u: u\mid a^2,\ R_a\mid 4u+1\},\qquad M_a=\{u: u\mid a^2,\ R_a\mid u+a\}.
\]
A \emph{shell} $a$ is \emph{occupied} if $E_a\cup M_a\neq\emptyset$.

\begin{theorem}[shell criterion]\label{thm:shell}
\begin{enumerate}[label=(\alph*)]
\item If $x\le y\le z$ is a solution, then $p/4<x\le3p/4$. More precisely, $x<p/2$ and $x<y$, unless $p\equiv3\pmod4$ and $(x,y,z)=\bigl(\frac{p+1}2,\frac{p+1}2,\frac{p(p+1)}4\bigr)$. For $p=12h+1$ this gives $3h+1\le x\le 6h$.
\item For each integer $a>p/4$ not divisible by $p$ (in particular each $a$ with $p/4<a<p$), the ordered pairs $(y,z)$ of positive integers with $1/a+1/y+1/z=4/p$ number exactly $2|E_a|+|M_a|$. A pair with $y=z$ occurs if and only if $R_a=1$ (possible only for $p\equiv3\pmod4$ and $a=(p+1)/4$); otherwise $|M_a|$ is even and the unordered pairs number $|E_a|+\tfrac12|M_a|$. For $p\equiv1\pmod4$ this is always the case.
\item Hence the equation is solvable at $p$ if and only if some shell $a$ with $p/4<a<p/2$ is occupied.
\end{enumerate}
\end{theorem}
\status{Proved here. Where each part is in the workbench:
\begin{itemize}
\item \emph{The bijection of (b).} Archive Theorem~9.1 (\emph{Exact divisor-shell bijection}, p.~38) states it for every odd prime $p$ and every $R\equiv-p\pmod4$ with $\gcd(R,pa)=1$. That is exactly the generality of (b), since $\gcd(R,pa)=1$ holds if and only if $p\nmid a$.
\item \emph{The count $2|E_a|+|M_a|$.} Archive Theorem~9.64 (p.~90) splits the divisors by their $p$-adic layer and proves by an involution that the two outer layers have equal size. The targets $-4a^2$, $-a$ and $-4^{-1}$ are the archive's table (394) (§10, pp.~93--94, stated there for $R\equiv3\pmod4$).
\item \emph{The version for $p=12h+1$.} The supplement's ES-01/ES-02 (§1), audited in the first pass, state the criterion for $p=12h+1$ and $a\in[3h+1,9h]$.
\end{itemize}
This paper adds (a), with its exceptional solution at $p\equiv3\pmod4$; the analysis of the swap in (b), that is, when $y=z$ occurs, the parity of $|M_a|$ and the unordered count; and (c).
The halved range at $p\equiv1\pmod{12}$, with $x$ unique, is already a reduction step of the workbench's item~22, checked by reading in the first pass (\note{43a\_} §2.22), and items 6 and 7 work in it; the workbench's later packages attribute these coordinates to Elsholtz and Tao \cite{ElsholtzTao}, §2. Checked against a brute-force count of the ordered pairs on all $4{,}726$ pairs $(p,a)$ with $p<260$ an odd prime and $p/4<a<p$ (a pair with $y=z$ exactly in the $29$ cases with $R_a=1$) and on all $4{,}303$ pairs with $p<120$, $p/4<a\le3p$ and $p\nmid a$; on all $4{,}398$ shells of the primes \mbox{$p\equiv1\pmod{12}$} below $700$, and in the first pass on all $8{,}124$ shells below $1000$; (a) against all $2{,}259$ solutions at the odd primes below $400$; every prime \mbox{$p\equiv1\pmod{12}$} below $30{,}000$ has an occupied shell.}

\begin{proof}
(a) The other two terms are positive, so $1/x<4/p$; and $x$ is the least denominator, so $4/p\le3/x$. Hence $p/4<x\le3p/4$.
Suppose $x>p/2$, so $x\ge(p+1)/2$ as $p$ is odd, and $x<p$. Then $2/y\ge1/y+1/z=4/p-1/x>2/p$, so $y<p$. From $4xyz=p(xy+yz+zx)$ the prime $p$ divides $xyz$, and as $0<x\le y<p$ it divides $z=pz'$; then $xy(4z'-1)=pz'(x+y)$ with $p\nmid xy$, so $4z'-1=pk$ with $k\ge1$. Put $N=pk+1=4z'$; then $1/x+1/y=4k/N$, that is, $(4kx-N)(4ky-N)=N^2$ with both factors positive. But $4kx-N\ge2k(p+1)-N=N+2(k-1)$, and likewise for $y$; so $k=1$ and $x=y=N/2=(p+1)/2$, with $z=p(p+1)/4$, which needs $p\equiv3\pmod4$. Finally, if $x<p/2$ and $y=x$, then $1/x+1/y>4/p$, which is impossible. For $p=12h+1$ the integers in $(p/4,p/2)$ are $3h+1,\dots,6h$.

(b) Write $R=R_a$, $S=S_a$. Then $R\ge1$ is odd, because $p$ is. Since $p\nmid a$, $\gcd(a,R)=\gcd(a,p)=1$ and $\gcd(p,R)=\gcd(p,4a)=1$, so $S$ is a unit modulo $R$.
The equation $1/y+1/z=R/S$ is equivalent to $(Ry-S)(Rz-S)=S^2$, and both factors are positive because $1/y$ and $1/z$ are smaller than $R/S$. So the ordered pairs correspond bijectively to ordered factorisations $S^2=dd'$ into positive divisors with $d\equiv d'\equiv-S\pmod R$, through $y=(d+S)/R$, $z=(d'+S)/R$. If $d\mid S^2$ and $d\equiv-S$, then $d$ is a unit and $d'=S^2/d\equiv S^2(-S)^{-1}=-S$; so the pairs are counted by $\#\{d\mid S^2: d\equiv-S\pmod R\}$.
Since $\gcd(p,a)=1$, each divisor of $S^2=p^2a^2$ is uniquely $d=p^iv$ with $i\in\{0,1,2\}$ and $v\mid a^2$. Using $4a\equiv p\pmod R$:
for $i=2$, $p^2v\equiv-pa\iff pv\equiv-a\iff 4pv\equiv-p\iff R\mid 4v+1$, so these $v$ form $E_a$;
for $i=1$, $pv\equiv-pa\iff R\mid v+a$, so these $v$ form $M_a$;
and $d\mapsto S^2/d$ maps the layer $i=2$ bijectively onto the layer $i=0$ while preserving the congruence. The total is $2|E_a|+|M_a|$.
The swap $(y,z)\mapsto(z,y)$ is $d\mapsto S^2/d$; it maps each layer $i$ to the layer $2-i$, so a fixed point lies in the layer $i=1$, where it is $v\mapsto a^2/v$ with only possible fixed point $v=a$. That needs $R\mid 2a$, i.e. $R\mid2$ (as $\gcd(R,a)=1$), i.e. $R=1$. If $R\ne1$ the swap has no fixed point, $|M_a|$ is even, $y\ne z$, and the unordered pairs number $|E_a|+\tfrac12|M_a|$. For $p\equiv1\pmod4$, $R\equiv-p\equiv3\pmod4$, so $R\ne1$.

(c) For $p\equiv1\pmod4$ it follows from (a) and (b). For $p\equiv3\pmod4$ both sides hold: the shell $a=(p+1)/4$ lies in $(p/4,p/2)$ and has $R_a=1$, so $u=a$ lies in $M_a$.
\end{proof}

The workbench also gives the exact bookkeeping around this criterion:
\begin{itemize}
\item an explicit bijection between a tagged set and the ordered solutions, with its inverse read off from the $p$-adic valuation of $Ry-S$ (archive (526), p.~151; supplement §1);
\item the global statement that a counterexample prime would satisfy the finite, computable condition $E_a=M_a=\emptyset$ for every $a$.
\end{itemize}
Both were re-derived in the first pass (\note{43a\_} §1.2).

\subsection{The first two shells}

\begin{proposition}[the residual-3 shell]\label{prop:r3}
Let $a=3h+1=(p+3)/4$, so $R_a=3$. Let $P_1=\prod(2v_\ell(a)+1)$ over the primes $\ell\equiv1\pmod3$ dividing $a$, and $P_2$ the same product over the primes $\ell\equiv2\pmod3$. Then $E_a=M_a=\{u\mid a^2: u\equiv2\pmod3\}$ and $|E_a|=|M_a|=P_1(P_2-1)/2$. So a solution with least denominator $a$ exists if and only if $a$ has a prime factor $\equiv2\pmod3$, and there are then exactly $3P_1(P_2-1)/4$ unordered solutions with least denominator $a$.
\end{proposition}
\status{Proved here; the workbench's residual-3 sieve (supplement §1, \texttt{first\_two\_shell\_sieve.tex}; archive Theorem~10.3, \emph{Exact $R=3$ factor criterion}, p.~94, for $p\equiv1\pmod{24}$), audited. Checked on all 4{,}472 primes $p\equiv1\pmod{12}$ below $2\cdot10^5$, and in the first pass on all 19{,}564 below $10^6$.}

\begin{proof}
Since $a\equiv1\pmod3$, $3\nmid a$. Modulo $3$ both $4u+1$ and $u+a$ are $\equiv u+1$, so both conditions say $u\equiv2\pmod3$. A divisor $u=\prod\ell^{e_\ell}$ of $a^2$ has $u\equiv(-1)^{\sum e_\ell}\pmod3$, the sum over the primes $\ell\equiv2\pmod3$. Their exponent vectors ($0\le e_\ell\le2v_\ell(a)$) number $P_2$, and the difference between the numbers with even and with odd sum is $\prod_\ell\sum_{e=0}^{2v_\ell}(-1)^e=1$; so $(P_2-1)/2$ have odd sum, and the other primes contribute the factor $P_1$. By Theorem~\ref{thm:shell}(b) the unordered solutions containing $a$ number $\tfrac32P_1(P_2-1)/2$, and $a=3h+1$ is the least denominator of each of them because it is the smallest possible one. Finally $P_2>1$ exactly when $a$ has a prime factor $\equiv2\pmod3$. (The count is an integer: $a\equiv1\pmod 3$ makes the total exponent of those primes even, so $P_2\equiv1\pmod4$.)
\end{proof}

\begin{proposition}[the residual-7 shell]\label{prop:r7}
Let $a=3h+2=(p+7)/4$, so $R_a=7$ and $7\nmid a$. Let $n_3$ be the total exponent in $a$ of the primes $\equiv3\pmod7$ and $n_{24}$ that of the primes $\equiv2$ or $4\pmod 7$. The shell is occupied if and only if $a$ has a prime factor $\equiv5$ or $6\pmod7$, or $n_3\ge3$, or $n_3\ge1$ and $n_{24}\ge1$.
\end{proposition}
\status{Proved here; the workbench's residual-7 sieve, audited, together with its count as coefficients of $\prod_\ell\sum_{e=0}^{2v_\ell(a)}T^{\lambda(\ell)e}$ in $\ZZ[T]/(T^6-1)$. Checked on all 4{,}472 primes $p\equiv1\pmod{12}$ below $2\cdot10^5$ (and below $10^6$ in the first pass).}

\begin{proof}
The exterior condition $7\mid4u+1$ is $u\equiv-4^{-1}\equiv5=3^5$, and the middle condition is $u/a\equiv-1=3^3\pmod7$; in base $3$ the discrete logarithms of $1,\dots,6$ are $0,2,1,4,5,3$.
\emph{Sufficiency.} A prime $\ell\equiv5$ dividing $a$ gives $u=\ell$; a prime $\ell\equiv6$ gives $u=a\ell$. If $n_3\ge3$, a product of primes $\equiv3$ dividing $a^2$ with total exponent $5$ gives $u\equiv3^5$. If primes $t\equiv3$ and $\ell\in\{2,4\}\pmod 7$ divide $a$, then $u=a\ell t$ (when $\ell\equiv2=3^2$) or $u=a\ell/t$ (when $\ell\equiv4=3^4$) divides $a^2$ and has $u/a\equiv3^3$.
\emph{Necessity.} If none of the conditions holds, every prime factor of $a$ is $\equiv1,2,3$ or $4$, and either $n_3=0$ or ($1\le n_3\le2$ and $n_{24}=0$). If $n_3=0$, $a$ and all divisors of $a^2$ lie in the subgroup $\{1,2,4\}$ of squares, which contains neither $5$ nor $-1$. Otherwise $a\equiv3^{n_3}$ and a divisor of $a^2$ is $\equiv3^e$ with $0\le e\le2n_3\le4$; the exterior target needs $e\equiv5\pmod6$ and the middle one $e-n_3\equiv3\pmod6$ with $|e-n_3|\le2$, both impossible.
\end{proof}

The shells with $R=3$ and $R=7$ exist exactly for the primes $p\equiv1\pmod4$, with $a=(p+3)/4$ and $a=(p+7)/4$ (\textsf{Generalised here}; checked on all $1{,}611$ such primes below $3\cdot10^4$). Proposition~\ref{prop:r7} holds for all of them, since its proof never uses $p\equiv1\pmod3$. Proposition~\ref{prop:r3} holds as stated for $p\equiv1\pmod{12}$ only. For $p\equiv5\pmod{12}$, $a=(p+3)/4\equiv2\pmod3$, so modulo $3$ the middle condition reads $u\equiv1$; then $E_a=\{u\mid a^2:u\equiv2\}$ and $M_a=\{u\mid a^2:u\equiv1\}$, with $|E_a|=P_1(P_2-1)/2$ and $|M_a|=P_1(P_2+1)/2$ by the parity count of the proof above, and there are exactly $|E_a|+\frac12|M_a|=P_1(3P_2-1)/4$ unordered solutions with least denominator $a$ (an integer: the total exponent in $a$ of the primes $\equiv2\pmod3$ is odd, so $P_2\equiv3\pmod4$). That shell is always occupied ($1\in M_a$), and its occupancy criterion holds trivially, since $a\equiv2\pmod3$ has a prime factor $\equiv2\pmod3$; those primes are solvable anyway, since $p\equiv2\pmod3$.

\begin{proposition}[the two-shell survivors]\label{prop:survivors}
Let $p\equiv1\pmod{12}$ be prime. Both shells $a=(p+3)/4$ and $a=(p+7)/4$ are unoccupied if and only if the following three conditions hold:
\begin{itemize}
\item $p\equiv1\pmod{24}$;
\item every prime factor of $(p+3)/4$ is $\equiv1\pmod3$;
\item every prime factor of $(p+7)/4$ is $\equiv1,2$ or $4\pmod7$.
\end{itemize}
Every such $p$ is a square modulo $7$ and satisfies $p\equiv\pm1\pmod5$. So it lies in one of the six hard classes $p\equiv1,121,169,289,361,529\pmod{840}$.
\end{proposition}
\status{Proved here. The conclusion modulo $840$ is archive Corollary~10.7 (p.~95), the contrapositive of its Theorem~10.6; the first pass stated the congruence mod $24$ as an observation on its computation. Checked on all primes $p\equiv1\pmod{12}$ below $10^6$. Of these, exactly $989$ have both shells unoccupied. All $989$ lie in the six hard classes, and each class occurs; the least in the classes $1,121,169,289,361,529$ are $2521$, $12721$, $2689$, $1129$, $1201$ and $3049$.}
\begin{proof}
If $p\equiv13\pmod{24}$, then $(p+3)/4$ is even and $2\equiv2\pmod3$, so the first shell is occupied (Proposition~\ref{prop:r3}). If $p\equiv1\pmod{24}$, the first shell is unoccupied exactly when every prime factor of $(p+3)/4$ is $\equiv1\pmod3$. Also $(p+7)/4$ is even, so $n_{24}\ge1$, and Proposition~\ref{prop:r7} says that the second shell is occupied exactly when $(p+7)/4$ has a prime factor $\equiv3$, $5$ or $6\pmod7$, a non-residue. If it has none, $(p+7)/4$ is a square modulo $7$, and so is $p\equiv4\cdot\frac{p+7}4$.
For the congruence modulo $5$:
\begin{itemize}
\item if $p\equiv2\pmod5$, then $5\mid(p+3)/4$ and $5\equiv2\pmod3$, so the first shell is occupied;
\item if $p\equiv3\pmod5$, then $5\mid(p+7)/4$ and $5$ is a non-residue modulo $7$, so the second shell is occupied;
\item $p=5$ is not $\equiv1\pmod{24}$.
\end{itemize}
So $p\equiv\pm1\pmod5$, and $p\equiv1,2,4\pmod7$, and $p\equiv1\pmod{24}$. These are the unit squares modulo $840$.
\end{proof}

\subsection{Which primes remain}

The workbench proves solvability for $p\equiv2\pmod 3$, $p\equiv3\pmod4$ and $p\equiv13\pmod{24}$ (scope identities that it does not advertise as new) by the families $(p,\frac{p+1}3,\frac{p(p+1)}3)$, $(\frac{p+1}4,\frac{p(p+1)}2,\frac{p(p+1)}2)$ and $(a,\frac{pa}2,pa)$ with $a=\frac{p+3}4$ (even when $p\equiv13\pmod{24}$); so its own reduction leaves the primes $p\equiv1\pmod{24}$ (archive Proposition~10.2). Its later packages restrict to the \emph{hard} primes $p\equiv1,121,169,289,361,529\pmod{840}$, which are exactly the unit squares modulo $840$ (a group of order $1\cdot1\cdot2\cdot3=6$). The reduction to them is classical (Obláth and Mordell; see Elsholtz--Tao \cite{ElsholtzTao} and Mordell \cite{Mordell}). The 22 September packages use it without proof.

The archive proves a stronger form (Theorem~10.6, \emph{Non-square modulo 840 theorem}, p.~95). For every prime $p\equiv1\pmod{24}$ outside the six hard classes, the shell with $R=3$ or the shell with $R=7$ is occupied. Its proof uses explicit divisors in these two shells; Proposition~\ref{prop:survivors} above is the same statement.
The classical proof uses shells up to $R=15$ and five identities. It is given here for comparison.

\begin{proposition}[the reduction modulo $840$]\label{prop:mod840}
Let $p\equiv1\pmod{24}$ be prime, and let $(A,B,R)$ be $(1,2,7)$, $(2,1,7)$, $(1,1,7)$, $(1,2,15)$ or $(2,1,15)$ according as $p\equiv3$, $5$, $6\pmod7$ or $p\equiv2$, $3\pmod5$. Then $D=(p+R)/(4AB)$ and $C=(A+pB)/R$ are positive integers and
\[
\frac4p=\frac1{ABD}+\frac1{ACD}+\frac1{pBCD}.
\]
So the equation is solvable at every prime $p\equiv1\pmod{24}$ that is not a square modulo $840$.
\end{proposition}
\status{Proved here (the reduction is classical); the first pass covered the $18$ non-square unit classes $\equiv1\pmod{24}$ by the same kind of families. The identities hold for all integers in the stated classes, not only for primes. Checked on all primes $p\equiv1\pmod{24}$ below $3\cdot10^5$.}
\begin{proof}
With $RC=A+pB$ and $4ABD=p+R$: $\frac1{ABD}+\frac1{ACD}+\frac1{pBCD}=\frac{pC+pB+A}{pABCD}=\frac{C(p+R)}{pABCD}=\frac{p+R}{pABD}=\frac4p$. Since $p\equiv1\pmod8$, $8$ divides $p+7$ and $p+15$, so $4AB\mid p+R$; $R\mid A+pB$ is the stated congruence modulo $7$ or $5$, together with $A+pB\equiv A+B\equiv0\pmod3$ when $R=15$. The primes $p\equiv1\pmod{24}$ left over have $p\equiv1,2,4\pmod7$ and $p\equiv\pm1\pmod5$, which are exactly the unit squares modulo $840$.
\end{proof}


\subsection{The Jacobi-symbol barrier}\label{ssec:jacobi}

\begin{proposition}[the Jacobi-symbol barrier]\label{prop:jacobi}
Let $p$ be an odd prime, $R\equiv3\pmod4$, $a=(p+R)/4$ an integer, and $\gcd(R,pa)=1$. If the Jacobi symbol $(q/R)$ equals $1$ for every prime $q\mid a$, then the shell $a$ is unoccupied. The converse is false: at $p=5$, $R=7$, $a=3$ the only prime factor of $a$ is a non-residue modulo $7$, and yet the shell is unoccupied.
\end{proposition}
\status{Proved here; archive Theorem~10.9 (\emph{Jacobi-symbol barrier}, p.~96) and Proposition~10.10 (\emph{The Jacobi barrier has no converse}, pp.~96--97), whose proofs are correct as read.}
\begin{proof}
The hypothesis $\gcd(R,pa)=1$ gives $p\nmid a$: otherwise $p$ would divide $4a-p=R$. So, by the proof of Theorem~\ref{thm:shell}(b), the shell is occupied exactly when some divisor $d$ of $S^2=p^2a^2$ satisfies $d\equiv-pa\pmod R$. Let $\chi=(\cdot/R)$ be the Jacobi symbol, a character on the units modulo $R$. The hypothesis gives $\chi(a)=1$, and $p\equiv4a\pmod R$ gives $\chi(p)=\chi(4)\chi(a)=1$. So $\chi(d)=1$ for every divisor $d$ of $p^2a^2$, while $\chi(-pa)=\chi(-1)=(-1)^{(R-1)/2}=-1$ because $R\equiv3\pmod4$. So no divisor lies in the target class.
For the converse: the squares modulo $7$ are $1,2,4$, so $(3/7)=-1$; the divisors of $15^2=225$ are $\equiv1,2,3,4,5\pmod7$, and the target $-15\equiv6$ is not among them.
\end{proof}

So at a prime $p\equiv1\pmod4$, a shell with residual $R$ can be occupied only if $a=(p+R)/4$ has a prime factor that is a non-residue modulo $R$ in the Jacobi sense. The argument is a quadratic-character argument: every divisor inherits the residue condition. This is the kind of mechanism behind the obstruction to polynomial identities recalled in Section~\ref{sec:problem}. The barrier explains why the workbench's occupancy criteria, such as Propositions~\ref{prop:r3} and~\ref{prop:r7}, are phrased through non-residue prime factors of $a$.

The quadratic character can be replaced by the whole group that the divisors generate. This gives a sharper barrier, and a converse when the exponents of $a$ are large enough. Both were found by the referee pass of version~1 of this paper; the proofs below were checked here. Part (c) of Proposition~\ref{prop:groupbarrier}, added in version~3, is from the author's notes and the archive (see its status).

\begin{proposition}[the group barrier]\label{prop:groupbarrier}
Let $p$ be an odd prime, $R\ge3$ with $R\equiv-p\pmod4$, $a=(p+R)/4$ an integer, and $\gcd(R,pa)=1$. Let $G_a\le(\ZZ/R)^\times$ be the subgroup generated by $4$ and the primes dividing $a$.
\begin{enumerate}[label=(\alph*)]
\item If $-1\notin G_a$, the shell $a$ is unoccupied.
\item This contains Proposition~\ref{prop:jacobi}. For prime $R\equiv3\pmod4$ the two hypotheses are equivalent. For composite $R$ the barrier (a) is strictly stronger: at $p=53$, $a=26$, $R=51$ one has $G_a=\langle2\rangle=\{1,2,4,8,13,16,26,32\}$, which does not contain $-1$, while $(2/51)=-1$.
\item \emph{(The width-one barrier.)} Let $H_a\le G_a$ be the subgroup generated by the primes dividing $a$. If neither $-1$ nor $-4$ lies in $H_a$, the shell $a$ is unoccupied. This contains (a), and it is strictly stronger. At $p=433$, $R=91$, $a=131$ (a prime), $H_a=\langle40\rangle=\{1,27,40,53,66,79\}$ contains neither $90=-1$ nor $87=-4$. Yet $-1\in G_a$, and $(131/91)=-1$, so neither (a) nor Proposition~\ref{prop:jacobi} applies; the shell is empty. When every prime factor of $R$ is $\equiv3\pmod4$, (a) and (c) are equivalent.
\end{enumerate}
\end{proposition}
\begin{proof}
(a) Every divisor $u$ of $a^2$, and $a$ itself, lie in $G_a$. The two targets $-4^{-1}$ and $-a$ lie in $G_a$ exactly when $-1$ does, because $4,a\in G_a$. So if $-1\notin G_a$, no divisor hits a target.
(b) If $(q/R)=1$ for every prime $q\mid a$, then $G_a$ lies in the kernel of the Jacobi character, which does not contain $-1$ when $R\equiv3\pmod4$. For prime $R\equiv3\pmod4$ the group $(\ZZ/R)^\times$ is cyclic of order $2\cdot\frac{R-1}2$ with $\frac{R-1}2$ odd. Its element $-1$ is the only element of order $2$. So $-1\notin G_a$ exactly when $G_a$ has odd order, that is, when $G_a$ lies in the subgroup of squares, which has odd order $\frac{R-1}2$. That is the hypothesis of Proposition~\ref{prop:jacobi}. At $R=51$: $2$ has order $2$ modulo $3$ and $8$ modulo $17$, hence order $8$ modulo $51$, with the powers listed. Also $13=2^6$ and $4=2^2$ modulo $51$, so $G_{26}=\langle2\rangle$. Finally $(2/51)=(2/3)(2/17)=-1$.
(c) Every divisor of $a^2$, and $a$, lie in $H_a$. The targets $-4^{-1}$ and $-a$ lie in $H_a$ exactly when $-4$, respectively $-1$, does. If $-1\notin G_a=\langle4,H_a\rangle$, then $-1\notin H_a$, and $-4\notin H_a$ since otherwise $-1=-4\cdot4^{-1}\in G_a$; so (c) contains (a).

Now let every prime factor of $R$ be $\equiv3\pmod4$. Then each factor $(\ZZ/q^k)^\times$ is cyclic of order $\equiv2\pmod4$. So $(\ZZ/R)^\times=P\times O$, where the $2$-part $P$ is elementary abelian and $O$ has odd order $|O|$.
\begin{itemize}
\item The element $4$ is a square, so it lies in $O$. Write $-1=(\varepsilon,1)$.
\item If $-1=4^jh$ with $h\in H_a$, then $h=(\varepsilon,w)$ for some $w\in O$.
\item So $h^{|O|}=(\varepsilon,1)=-1$ lies in $H_a$.
\end{itemize}
Hence $-1\in G_a$ implies $-1\in H_a$, and (a) and (c) are equivalent. The example was computed directly; there $91=7\cdot13$ has the factor $13\equiv1\pmod4$.
\end{proof}
\status{Part (c) is Theorem~3.1 of the author's note \emph{Normalized divisor supports and the residual pre-Niemeier datum} (Section~\ref{sec:notes}), and the \emph{reachable targets} of archive Theorem~15.5. It was added to this paper in version~3. Among the shells $R<3p$ of the primes $p\equiv1\pmod{24}$ below $3000$, it kills $124$ empty shells on which (a) is silent; every one of these $R$ has a prime factor $\equiv1\pmod4$. The equivalence for $R$ with all prime factors $\equiv3\pmod4$ was found by the referee of version~3.}

\begin{proposition}[a converse under an exponent condition]\label{prop:groupconverse}
In the setting of Proposition~\ref{prop:groupbarrier}, let $H_a$ be the subgroup generated by the primes dividing $a$. Suppose that $2v_\ell(a)\ge\operatorname{ord}_R(\ell)-1$ for every prime $\ell\mid a$.
\begin{enumerate}[label=(\alph*)]
\item The residues modulo $R$ of the divisors of $a^2$ are exactly the elements of $H_a$.
\item The shell $a$ is occupied if and only if $-1\in H_a$ or $-4\in H_a$ (equivalently $-4^{-1}\in H_a$, since $H_a$ is a group).
\item If moreover $R$ is a prime $\equiv3\pmod4$, the shell is occupied if and only if some prime factor of $a$ is a non-residue modulo $R$.
\end{enumerate}
\end{proposition}
\begin{proof}
(a) A divisor of $a^2$ is $\prod\ell^{e_\ell}$ with $0\le e_\ell\le2v_\ell(a)$. Under the hypothesis, $e_\ell$ runs through a complete set of residues modulo $\operatorname{ord}_R(\ell)$, so $\ell^{e_\ell}$ runs through all of $\langle\ell\rangle$. The products then fill the subgroup generated by these cyclic subgroups, which is $H_a$.
(b) The targets are $-4^{-1}$ and $-a$, and $-a\in H_a$ exactly when $-1\in H_a$, since $a\in H_a$.
(c) If every prime factor of $a$ is a residue, Proposition~\ref{prop:jacobi} applies. Suppose some prime factor is a non-residue. The group $(\ZZ/R)^\times$ is cyclic of order $R-1=2\cdot\frac{R-1}2$ with $\frac{R-1}2$ odd. So every subgroup of odd order lies in its unique subgroup of order $\frac{R-1}2$, which is the subgroup of squares. Since $H_a$ contains a non-square, its order is even. It therefore contains the unique element of order $2$, which is $-1$, and the shell is occupied by (b).
\end{proof}
\status{Both proved here. Checked on all $123{,}538$ shells with $R\ge3$ of the primes $5\le p<1500$: the group barrier holds on every one of them. On $9{,}462$ of them, all with composite $R\equiv3\pmod4$, it applies where the Jacobi barrier does not. The criterion of Proposition~\ref{prop:groupconverse}(b) is correct on all $677$ shells that satisfy the exponent condition. For prime $R\equiv3\pmod4$ the two barriers agree on all $18{,}157$ shells of the primes $p\equiv1\pmod4$ below $1500$. The counterexample $p=5$, $R=7$, $a=3$ to the converse of Proposition~\ref{prop:jacobi} fails the exponent condition: $\operatorname{ord}_7(3)=6>2v_3(3)+1=3$.}

\section{Least classes, record primes, explicit families, and the carrier census}\label{sec:families}

\subsection{The least-class function and its record primes}\label{ssec:records}

The archive (§17.4 and §19, pp.~131 and 148--157) puts, for $i\ge1$, following Ionascu and Wilson \cite{IonascuWilson},
\[
\mathcal C_i=\Bigl\{n\ge2:\ \frac4n=\frac1x+\frac1y+\frac1z\ \text{with } x\le\frac{n+4i-1}4\Bigr\},\qquad h(n)=\min\{i:n\in\mathcal C_i\}.
\]
An odd prime $p$ is a \emph{strict record} if $h(q)<h(p)$ for every prime $q<p$.

\begin{lemma}\label{lem:leastclass}
Let $p\equiv1\pmod4$. The least denominator occurring in any solution at $p$ is the least occupied shell $a$. With $R=4a-p$ one has $h(p)=(R+1)/4$. For $p\equiv3\pmod4$, $h(p)=1$, and $h(2)=1$.
\end{lemma}
\status{Proved here.}
\begin{proof}
If $a$ is occupied, some solution has $a$ as a denominator, so the least denominator of that solution is at most $a$. Conversely, the least denominator $x$ of a solution is itself an occupied shell: it exceeds $p/4$ (Theorem~\ref{thm:shell}(a)) and is completed by the other two denominators. So the least occupied shell equals the least denominator over all solutions. A solution lies in $\mathcal C_i$ exactly when $4x-p\le4i-1$. For $p\equiv1\pmod4$ the residual $4x-p$ is $\equiv3\pmod4$, so the least $i$ is $(R+1)/4$. For $p\equiv3\pmod4$ the shell $(p+1)/4$ is occupied (Theorem~\ref{thm:shell}(c)). For $p=2$, take $4/2=1/1+1/2+1/2$ with $x=1\le5/4$.
\end{proof}

\begin{corollary}\label{cor:hle2}
If a prime $p$ does not lie in one of the six hard classes $1,121,169,289,361,529$ modulo $840$, then $h(p)\le2$. Hence every strict record after $73$ lies in a hard class.
\end{corollary}
\status{Proved here; found by the referee pass of version~1. It is the archive's Theorem~10.6 and Proposition~10.2 read through Lemma~\ref{lem:leastclass}. An independent computation of $h(p)$ for all $144{,}449{,}537$ primes below $3\cdot10^9$ finds no prime with $h(p)\ge3$ outside the hard classes.}
\begin{proof}
For $p=2$ and $p\equiv3\pmod4$, $h(p)=1$ by Lemma~\ref{lem:leastclass}. Let $p\equiv1\pmod4$.
\begin{itemize}
\item If $p\equiv5\pmod{12}$, the shell $(p+3)/4$ is occupied (Section~\ref{sec:core}, after Proposition~\ref{prop:r7}).
\item If $p\equiv13\pmod{24}$, the same shell is occupied (Proposition~\ref{prop:survivors}, first line of the proof).
\item If $p\equiv1\pmod{24}$ is not hard, Proposition~\ref{prop:survivors} shows that the shell with $R=3$ or the shell with $R=7$ is occupied.
\end{itemize}
In each case $R\le7$, so $h(p)\le2$. A strict record after $73$, where $h=2$, has $h\ge3$.
\end{proof}

\begin{proposition}[the record primes]\label{prop:records}
The strict records $p\le8{,}803{,}369$, with $h(p)$ and the residual $R$ of the least occupied shell, are exactly
\[
(2,1,2),\ (73,2,7),\ (1129,3,11),\ (1201,6,23),\ (21169,8,31),\ (118801,15,59),\ (8803369,27,107).
\]
For the odd records $R=4h(p)-1$; for $p=2$ the least denominator is $1$, so $R=4\cdot1-2=2$. Moreover $h(806521)=h(118801)=15$, so $806{,}521$ is not a strict record. There is no strict record between $8{,}803{,}369$ and $3\cdot10^9$. Below $3\cdot10^9$, no prime other than $8{,}803{,}369$ has $h\ge22$. The primes with $h=19$, $20$ or $21$ are
\begin{itemize}
\item $h=21$: $287{,}567{,}281$, $651{,}412{,}441$, $2{,}188{,}875{,}529$;
\item $h=20$: $778{,}103{,}881$, $1{,}749{,}233{,}641$, $2{,}535{,}613{,}921$, $2{,}809{,}257{,}889$;
\item $h=19$: $230{,}089{,}729$, $793{,}592{,}809$.
\end{itemize}
\end{proposition}
\status{Reproduced here: every prime up to $8{,}803{,}369$ was computed by Lemma~\ref{lem:leastclass} and Theorem~\ref{thm:shell}, using the residues of the divisors of $a^2$ modulo $R$ (script \note{esr\_records.py}, under three seconds).
\begin{itemize}
\item \emph{Source of the list.} Ionascu and Wilson print the record list \cite{IonascuWilson}. The archive (Proposition~17.6, pp.~131--132) recovers every printed record prime from the definition. It corrects one printed class label: $1201\in\mathcal C_6\setminus\mathcal C_5$, where the source prints $\mathcal C_4$. Its table (IW-RC) (p.~133) gives $R=2$ for $p=2$.
\item \emph{The range to $3\cdot10^9$.} It was found by the referee pass of version~1 with its own program, and recomputed here by an independent segmented program. The two agree on every record, on every prime with $h\ge19$, and on the histogram of $h$.
\item \emph{The shell at $p_*$.} At $p_*=8{,}803{,}369$ the least occupied shell has $R=107$, the shell of archive §3.
\end{itemize}}

For the five records whose residual $R$ is a safe prime ($R=2m+1$ with $m$ prime), and whose shell $a=(p+R)/4$ has exactly one non-residue prime factor modulo $R$, to exponent one, the archive defines a \emph{deficit} set $D\subset\ZZ/m$ from the discrete logarithms of the prime factors of $a$ (§19.1). It states the \emph{strict-record deficit-box conjecture} (§19.2): at such a record with non-empty deficit, $D$ is a difference-proper generalised arithmetic progression of rank at most two, and two induced \emph{odd sheets} cover $\ZZ/m$.
The deficits of the first three rows are empty. The two non-empty ones are proved to be a rank-two rectangle $\{11,14,15,18\}\subset\ZZ/29$ at $(118801,59)$ and the progression $\{23+42j:0\le j\le9\}\subset\ZZ/53$ at $(8803369,107)$.
The boundary example $p=806{,}521$ satisfies every hypothesis except strict recordness, and both conclusions fail there. So the hypothesis cannot be dropped.

\status{The conjecture is open, as the archive says. Its finite evidence is \textsf{reproduced} here from the definitions (script \note{esr\_deficits.py}, with the primitive root $g=2$ that the archive uses for $R=59$ and $R=107$):
\begin{itemize}
\item the deficits of the rows $p=73$, $1129$ and $1201$ are empty, and the row $p=21169$ is outside the domain because $m=15$ is composite;
\item at $p=118{,}801$: $a=29715=3\cdot5\cdot7\cdot283$, $D=\{11,14,15,18\}$, $D-D=\{3r+4s:|r|,|s|\le1\}$ with nine elements, offsets $u=11$ and $v=17$, and the two odd sheets cover $\ZZ/29$;
\item at $p=8{,}803{,}369$: $a=2{,}200{,}869=3^2\cdot11^2\cdot43\cdot47$, $D=\{23+42j:0\le j\le9\}$, $D-D=\{42k:|k|\le9\}$ with nineteen elements, offsets $29$ and $23$, and the sheets cover $\ZZ/53$; also the corollary of archive §16, that the middle target $-1$ is reached by exactly two exponent vectors and the exterior target not at all;
\item at $p=806{,}521$: $a=201{,}645=3^2\cdot5\cdot4481$, a deficit of $14$ points with $D-D=\ZZ/29$, and the odd sheets miss exactly $\{1,4,8,12,16,20,24,27\}$.
\end{itemize}
Another primitive root multiplies $D$ by a unit of $\ZZ/m$ (archive §19.1). The script confirms that the sizes of $D$ and $D-D$, and the number of missed points, do not change.
The failure of conclusion (i) at $806{,}521$ follows from $D-D=\ZZ/29$ by the archive's counting argument. A difference-proper progression of rank one with $14$ terms has exactly $27$ differences. A proper one of rank two has side lengths $\{2,7\}$, so a difference-proper one would have $3\cdot13=39>29$ distinct differences. That argument was read and is correct.
By Theorem~\ref{thm:shell}(b), the counts at $p_*=8{,}803{,}369$ ($|E_a|=0$, $|M_a|=2$) mean that exactly one unordered solution has least denominator $2{,}200{,}869$, namely the archive's (405) (p.~97), from the middle divisor $u=11^2$:
\[
\frac4{8803369}=\frac1{2200869}+\frac1{181085300330}+\frac1{3293760527702370}.
\]
The archive's Lean module covers the $R=107$ row only, as the archive states.}

\subsection{Fixed-divisor families, and four progressions on which the equation is solved}\label{ssec:families}

Theorem~\ref{thm:shell}(b) turns a single divisor into a solution. If the divisor is a fixed multiple of the cofactor $h=a/c$, the solution is given by polynomials in $n$ on a whole residue class. The archive states this for primes, as its constant-factor carriers (Lemma~24.218, p.~422). The form below holds for all $n$. It was found by the referee pass of version~1 and checked here.

\begin{proposition}[fixed-divisor families]\label{prop:fixeddivisor}
Let $R\ge3$ be odd, and let $c,s\ge1$ with $\gcd(c,R)=1$ and $s\mid c^2$. Let $n\ge1$ with $n\equiv-R\pmod{4c}$, and put $a=(n+R)/4$ and $h=a/c$, both integers.
\begin{enumerate}[label=(\roman*)]
\item \emph{Exterior.} If $n\equiv-cs^{-1}\pmod R$, then $4/n=1/a+1/y+1/z$ with the positive integers
\[
y=\frac{n(nsh+a)}R,\qquad z=\frac{c^2h/s+na}R .
\]
\item \emph{Middle.} If $n\equiv-4s\pmod R$, then $4/n=1/a+1/y+1/z$ with the positive integers
\[
y=\frac{n(a+s)}R,\qquad z=\frac{n(a+a^2/s)}R .
\]
\end{enumerate}
Each case is one residue class modulo $4cR$. If $\gcd(cR,35)=1$, that class meets all six hard classes modulo $840$ or none of them. It meets all six exactly when it is $\equiv1$ modulo $\gcd(4cR,24)$.
\end{proposition}
\begin{proof}
Put $S=na$. In both cases $n$ is a unit modulo $R$, because $\gcd(cs,R)=1$. So is $a$, because $4a\equiv n$. Hence $S$ is a unit modulo $R$. The proof of Theorem~\ref{thm:shell}(b) uses only this. A positive divisor $d$ of $S^2$ with $d\equiv-S\pmod R$ therefore gives positive integers $y=(d+S)/R$ and $z=(S^2/d+S)/R$ with $1/y+1/z=R/S$. Then $1/a+1/y+1/z=(n+R)/(na)=4/n$.
\begin{itemize}
\item In case (i), take $d=n^2u$ with $u=sh$. Here $u\mid c^2h^2=a^2$, and $d\equiv-S$ means $nu\equiv-a$, that is, $4u\equiv-1\pmod R$. That holds because $c(4sh+1)=s(n+R)+c\equiv sn+c\equiv0\pmod R$. Then $S^2/d=a^2/u=c^2h/s$.
\item In case (ii), take $d=nu$ with $u=s$. Here $s\mid c^2\mid a^2$, and $d\equiv-S$ means $s\equiv-a$, that is, $4s\equiv-n\pmod R$. Then $S^2/d=na^2/s$.
\end{itemize}
Since $R$ is odd and prime to $c$, the congruences modulo $4c$ and modulo $R$ define one class modulo $4cR$ (Chinese remainder theorem). If $\gcd(cR,35)=1$, then $\gcd(4cR,840)=\gcd(4cR,24)$. A class $r$ modulo $4cR$ meets a class $t$ modulo $840$ exactly when $r\equiv t$ modulo $\gcd(4cR,840)$. The six hard classes are all $\equiv1\pmod{24}$.
\end{proof}
\status{Proved here. For primes it is archive Lemma~24.218; the archive's proof of its Theorem~24.219 identifies the four families below as these middle and exterior channel maps.}

\begin{theorem}[four unconditional families]\label{thm:families}
For $k\ge0$ put $g_1=31k+30$, $g_2=31k+15$, $m_3=23k+15$, $m_4=23k+19$, $w_3=1116k+727$ and $w_4=4464k+3683$. Then $4/n_i=1/a_i+1/b_i+1/c_i$ with positive integers
\begin{center}\small
\begin{tabular}{@{}cllll@{}}
\toprule
$i$ & $n_i$ & $a_i$ & $b_i$ & $c_i$\\
\midrule
1 & $248k+209$ & $2g_1$ & $2n_1(k+1)$ & $2n_1g_1(k+1)$\\
2 & $248k+89$ & $2g_2$ & $n_2(2k+1)$ & $2n_2g_2(2k+1)$\\
3 & $17112k+11137$ & $186m_3$ & $12n_3m_3w_3$ & $124m_3w_3$\\
4 & $17112k+14113$ & $186m_4$ & $6n_4m_4w_4$ & $31m_4w_4$\\
\bottomrule
\end{tabular}
\end{center}
No primality of $n_i$ is needed.
\end{theorem}
\status{Proved here by direct expansion: $4a_ib_ic_i-n_i(b_ic_i+a_ic_i+a_ib_i)$ is the zero polynomial in $k$ for each $i$ (sympy), and the entries are positive for $k\ge0$; exact fractions agree for $k\le200$ (script \note{esr\_families.py}). Archive Theorem~24.219, p.~423.
\begin{itemize}
\item \emph{As fixed-divisor families.} The four families are instances of Proposition~\ref{prop:fixeddivisor}. With $(R,c,s)=(31,2,2)$ and $(31,1,1)$ in case (ii), and $(23,186,18)$ and $(23,186,36)$ in case (i), where $h=m_3$ and $h=m_4$, its formulas give exactly the table (checked by expansion). Two of the families are halves of larger instances (archive Proposition~24.222(ii)--(iii), p.~425):
\begin{itemize}
\item family 2 is half of the class $n\equiv89\pmod{124}$, which is $(31,1,1)$ itself;
\item family 3 is half of the class $n\equiv2581\pmod{8556}$, which is $(23,93,9)$ in case (i).
\end{itemize}
\item \emph{Rosati.} The archive shows that these triples are specialisations of Rosati's two 1954 parametrisations. It claims no independence from the classical parametrisations (archive abstract, p.~2, and §13.1). That comparison is located, not checked.
\end{itemize}}
\begin{proof}
All displayed linear factors have positive constant terms, so every entry is a positive integer for $k\ge0$. The identity $4a_ib_ic_i=n_i(b_ic_i+a_ic_i+a_ib_i)$ is a polynomial identity in $k$, verified by expanding both sides. Dividing by $n_ia_ib_ic_i$ gives the decomposition.
\end{proof}

Each progression meets all six hard classes modulo $840$. This is the last part of Proposition~\ref{prop:fixeddivisor}, and directly: since $\gcd(248,840)=8$ and $\gcd(17112,840)=24$, the numbers $n_1,n_2$ run through every class $\equiv1\pmod8$ modulo $840$, and $n_3,n_4$ through every class $\equiv1\pmod{24}$. Every hard class is $\equiv1\pmod{24}$, being the square of a unit.
Each $n_i(0)$ is prime to the step of its progression, so by Dirichlet's theorem every progression contains infinitely many primes in each hard class. Examples are the strict records $1201=n_1(4)\equiv361$ and $21169=n_2(85)\equiv169\pmod{840}$.
So the families solve the equation at infinitely many primes in every class that the reduction modulo $840$ leaves open (script \note{esr\_families.py} lists the least prime for each family and class).
This does not conflict with the obstruction to polynomial identities recalled in Section~\ref{sec:problem}. That obstruction is modulo $q$, and it applies only when $r$ is a square modulo $q$: here $11137\equiv5\pmod{23}$, and $5$ is not a square modulo $23$. The same holds for the other rows, since $14113\equiv14\pmod{23}$, $209\equiv23\pmod{31}$ and $89\equiv27\pmod{31}$ are non-squares. Proposition~\ref{prop:nosquare} below shows that this is forced: no fixed-divisor family covers a square class.

\subsection{The carrier census}\label{ssec:census}

The largest part of the archive's arithmetic is a sequence of \emph{sufficient-carrier censuses}. It occupies §18 and §24 from the base of the census (Theorem~24.100, p.~335) to the end of §24.3 (p.~477).
\begin{itemize}
\item \emph{Domain.} The census concerns only the primes $p\equiv25\pmod{72}$ whose shells $R=7$ and $R=11$ are both empty (archive Corollary~18.17 and Theorem~24.100). Such a prime lies in one of $30$ classes modulo $11088=16\cdot9\cdot7\cdot11$, the base rows $B_1\cup B_{25}$. The primes $p\equiv1$ and $p\equiv49\pmod{72}$ with these two shells empty are outside the census. Below $3\cdot10^6$ there are $2{,}316$ primes $p\equiv1\pmod{24}$ with both shells empty; $829$ are $\equiv1$, $682$ are $\equiv25$ and $805$ are $\equiv49\pmod{72}$ (computed here).
\item \emph{Rows.} The archive considers residue data of $p$ modulo a growing product of moduli. We call each such datum a \emph{row}. The moduli are $11088$, $5$, $19$, $23$, $31$ and $47$, then $71$, $191$, $83$, $167$ and $239$.
\item \emph{Forced rows.} A row is \emph{forced} when an explicit construction shows that, for every prime $p$ in the branch with that row, one of a finite list of shells is occupied. An example is the implication~(2272) on p.~431: for a prime $p\equiv25\pmod{72}$ whose shells $R=7$ and $R=11$ are empty and whose row lies in the domain but outside the survivor set, some shell with $R\in\{15,19,23,31,47\}$ is occupied.
\item \emph{Survivors.} The remaining rows are \emph{survivors}. Nothing is concluded for them.
\end{itemize}
The archive's front section (pp.~3--5) reports the following chain.
\begin{center}\footnotesize
\begin{tabular}{@{}>{\raggedright\arraybackslash}p{0.19\linewidth}rrr@{}}
\toprule
Step & Domain (rows) & Forced & Survivors\\
\midrule
six shells after the unit lift mod $71$ & $4{,}590{,}432{,}000$ & $4{,}212{,}008{,}640$ & $378{,}423{,}360$\\
$R=191$ & $71{,}900{,}438{,}400$ & $24{,}762{,}018{,}720$ & $47{,}138{,}419{,}680$\\
$Q=83$ & $3{,}865{,}350{,}413{,}760$ & $934{,}648{,}428{,}742$ & $2{,}930{,}701{,}985{,}018$\\
$R=167$ & $486{,}496{,}529{,}512{,}988$ & $115{,}851{,}891{,}806{,}170$ & $370{,}644{,}637{,}706{,}818$\\
$R=239$ & $88{,}213{,}423{,}774{,}222{,}684$ & $20{,}524{,}419{,}041{,}579{,}843$ & $67{,}689{,}004{,}732{,}642{,}841$\\
refresh of the $R=167$ coefficient & & $39{,}025{,}685{,}947{,}842$ & $67{,}649{,}979{,}046{,}694{,}999$\\
\bottomrule
\end{tabular}
\end{center}
Cumulatively $2{,}757{,}919{,}929{,}517{,}785{,}001$ of the $2{,}825{,}569{,}908{,}564{,}480{,000}$ rows of the final hard domain are forced, a proportion of $0.97606\ldots$.
The table is internally consistent (checked here):
\begin{itemize}
\item each domain is the previous survivor count times $q-1$, for the next modulus $q=191,83,167,239$;
\item the final domain is $4{,}590{,}432{,}000\cdot190\cdot82\cdot166\cdot238$;
\item forced plus survivors equals the domain in every row;
\item the fractions $69643/75900$ and $6257/75900$ of the first row are exact.
\end{itemize}

\status{Located. The arithmetic consistency of the reported numbers was checked here; the carrier theorems and the enumerations behind them were not re-run.
The archive states the following about its own checks.
\begin{itemize}
\item The enumerations are replayed by independent Python programs.
\item Bounded Lean modules check the literal census rows, totals, fractions, target residues and candidate orders, but not the affine and gcd reconstructions, carrier membership or the positive-denominator witnesses (pp.~3--5).
\item The ranking that selected $239$ after $167$ was corrected: on the full ten-coordinate object, $59$ ranks first and $239$ second. The $R=239$ step is kept as a valid non-greedy continuation, not as an optimal one (p.~4).
\end{itemize}
The archive calls these \emph{bounded sufficient-carrier censuses} and states that they are not shell-exhaustion, density, endpoint or optimality theorems (p.~2). A survivor row is only outside the carriers' images; it is not shown to be shell-empty.}

\paragraph{What the census reaches.}
Every carrier of the census is a fixed-divisor family in the sense of Proposition~\ref{prop:fixeddivisor}. On its row it fixes the residual $R$, a divisor $c$ of $a$ that the row determines, and a divisor $s$ of $c^2$, and uses $u=sh$ or $u=s$:
\begin{itemize}
\item for $R=15,19,23,31$ and for the factor-$31$ carrier at $R=23$, the maps listed in the proof of archive Theorem~24.228 (pp.~432--433), with $c=2,1,6$ or $12,2$ and $186$;
\item for $R=47$ and later, the sets $E_q(c)=\{-cs^{-1}\}$ and $M_q(c)=\{-4s\}$ over the divisors $s\mid c^2$, with $c=\gcd(a,N)$ read off the row (archive (2268)--(2269), p.~431, and the coefficients $c_Q$, pp.~462--463).
\end{itemize}
In each case $4cR$ divides the modulus of the row. The obstruction of Schinzel and Yamamoto therefore applies, in the following form.

\begin{proposition}[no fixed-divisor family covers a square class]\label{prop:nosquare}
In the setting of Proposition~\ref{prop:fixeddivisor}, the class modulo $4cR$ on which (i) or (ii) applies contains no square modulo $4cR$. So it contains no class that is a square modulo a multiple of $4cR$.
\end{proposition}
\begin{proof}
Suppose $n$ lies in the class and is a square modulo $4cR$.
\begin{itemize}
\item \emph{The prime $2$.} $n$ is odd, so $n\equiv1\pmod4$ and $R\equiv-n\equiv3\pmod4$. If $2\mid c$, then $8\mid4c$, so $n\equiv1\pmod8$, $R\equiv7\pmod8$ and $(2/R)=1$.
\item \emph{Odd primes $\ell\mid c$.} $-R\equiv n$ is a non-zero square modulo $\ell$, so $(-R/\ell)=1$. By reciprocity, $(\ell/R)=(R/\ell)(-1)^{(\ell-1)/2}=(-R/\ell)=1$, since $(R-1)/2$ is odd.
\end{itemize}
Hence $(c/R)=1$, and $(s/R)=1$ because $s\mid c^2$. But $n$ is a unit square modulo $R$, so $(n/R)=1$. In case (ii), $n\equiv-4s$ gives $(n/R)=(-1/R)(s/R)=-1$. In case (i), $n\equiv-cs^{-1}$ gives $(n/R)=(-1/R)(c/R)(s/R)=-1$. Either way this is a contradiction.
\end{proof}
\status{Proved here; found by the referee pass of version~1. It agrees with the archive's own residue sets (checked here):
\begin{itemize}
\item at $R=47$ the forced residues $F_{36}$ are exactly the $23$ non-residues modulo $47$, and $F_{18}$ consists of $21$ of them;
\item the unforced sets contain every quadratic residue of their modulus: $S_5=\{1,4\}$, $A_{25}=Q_{23}$, and $S_{19}$ and $S^\sharp_{31}$, whose excluded residues $15,18\bmod19$ and $15,23,27,29,30\bmod31$ are all non-residues.
\end{itemize}}

\begin{corollary}[the square rows of the census]\label{cor:censuslimit}
No row that is a square modulo the modulus of the census is ever forced by its carriers. All $30$ base rows are squares modulo $11088$. So the rows whose coordinates at the ten odd primes $5,19,23,31,47,71,191,83,167,239$ are all non-zero squares are never forced. They number
\[
\frac{2{,}825{,}569{,}908{,}564{,}480{,}000}{2^{10}}=2{,}759{,}345{,}613{,}832{,}500,
\]
that is, $1/1024$ of the final domain and $4.08\%$ of its $67{,}649{,}979{,}046{,}694{,}999$ survivors. Each further odd prime modulus added to the census halves this share, but no finite chain of such carriers removes it.
\end{corollary}
\begin{proof}
A forced row lies in a class of Proposition~\ref{prop:fixeddivisor} modulo $4cR$, and $4cR$ divides the modulus of the row. Proposition~\ref{prop:nosquare} applies. The base rows were checked directly: each is a square of a unit modulo $11088$ (script \note{esr\_referee\_checks.py}). The domain is $30\cdot\prod(q-1)$ over the ten odd primes, and exactly half of the units modulo each odd prime $q$ are squares.
\end{proof}

The survivors make up $2.39\%$ of the final domain. Each of the four lifts forced between $23\%$ and $35\%$ of its domain ($R=191$: $34.4\%$; $Q=83$: $24.2\%$; $R=167$: $23.8\%$; $R=239$: $23.3\%$). By Corollary~\ref{cor:censuslimit}, a residue census built from such carriers does not reach the square rows. Carriers that reach them are not fixed-divisor families; one kind of candidate uses the prime factors of $(p+R)/4$ (compare Proposition~\ref{prop:jacobi} and Theorem~I). The archive itself names the removal of the survivors as the open global obligation (p.~4).

\section{The September continuation}\label{sec:items}

The continuation of 17--23 September is organised as the 28-item \emph{Start reading} list of the README, with proof packages in \texttt{research/incoming/}. The workbench describes the results of its 22 September release as written proofs with Python replays, without a Lean build or independent human review (\note{43a\_} §2, citing \texttt{DAILY\_RESULTS\_20260922}).

\subsection{The audited items}

The 28 items of the README's \emph{Start reading} list (17--23 September) are listed in \note{43a\_} §2, with their packages. Those audited:
\begin{itemize}
\item \emph{Item 6} (audited). With $L_a=8\tau(a)^2/\varphi(R_a)-\mathcal E_a$, where $\mathcal E_a$ is the residue-collision energy of the $2\tau(a)$ divisors of $pa$ modulo $R_a$: for $p\equiv1\pmod4$, $p\ge2^{20}$, $\sum_{a\in(p/4,p/2)}L_a\le-\tfrac1{24}p\log p$. So the lower bound that keeps only the principal character, summed over the shells, is negative for large $p$, and that bound alone does not show occupancy.
\item \emph{Item 7} (audited). For $p\equiv1\pmod 8$, bounds in terms of the least quadratic nonresidue; at the hard primes a middle state $(a,u)$ on a shell $p/4<a<p/2$ has $a\le3(p+3)/11$, and an exterior state satisfies $22(p+1)\le23(p-2a+1)^2$. The workbench says this localises solutions and does not prove the region occupied.
\item \emph{Item 14} (audited). For $h'=4j-1>t$ the divisors $W\mid j^2$ with $W\equiv t\pmod{h'}$ are canonical, companion or finite templates; the finite ones satisfy the sharp bound $h'\le t^2-3t+1$, with equality exactly in one parametrised family, attained at infinitely many primes $p\equiv1\pmod{840}$. For every $t\ge4$ infinitely many hard primes have an exterior state whose same-grade middle box is empty.
\item \emph{Item 21} (audited). For a proper rational trace source the complete set of same-word residual returns is $\{k\mid R: d\mid k,\ k\equiv1\pmod{4K(u)}\}$; a uniform return exists exactly for the nine words $u\mid36$.
\item \emph{Item 22} (audited, certificate). For every solution at $p\equiv1\pmod{12}$ an explicit symmetric polynomial $\mathfrak N(p,x,y,z)$ of degree $8$ satisfies $\mathfrak N<-\frac{125873811}{262144}p^8$ ($\approx-480.17\,p^8$), and $\mathfrak N<-6960.19\,p^8$ when $p\equiv1\pmod{24}$; the proof reduces to finite polynomial positivity certificates, recomputed independently. It is a statement about existing solutions, not about their existence.
\item \emph{Item 26} (audited in part: the least example and the deduction from it; its prime-square trace theorem at $p\equiv1\pmod{24}$ and its terminality statement read and correct as read, with the identity $16(4q^3+1)-(4q+1)(16q^2-4q+1)=15$ checked; the factor-sum, Pell and fixed-tail-sum sections not checked). $p=67369$ is the least prime $p\equiv1\pmod{24}$ whose first \emph{trace shell} ($a=16849$, $R=27$) precedes its first full shell ($a=16850$, $R=31$); so no universal map from rational traces to solutions avoids increasing $a$.
\item \emph{Item 4} (conditional on a computer enumeration of $1{,}381{,}117{,}764$ rooted profiles, $5{,}922{,}259$ of them range-compatible, run by two implementations of the workbench and not re-run here). At every hard prime, from any nonresidue prime, factoring at most 25 numbers $(p+\sigma qt^2)/4$ with $t\in\{1,3,5,7,9\}$ yields at least six distinct nonresidue primes. A direct test of all 108 hard primes below $40{,}000$ agrees.
\end{itemize}
Items 2 and 3 are method counterexamples, and their certificates are audited: the seven-cycle sink of item~2 at $p=2521$ and the closed triangle of item~3; their closure and escape theorems are located. For item~8, the middle identity $p=RQ-(Q+1)^2/(4A)$ and its two bounds were reproduced on all $1{,}613$ oriented middle states at primes $p\equiv1\pmod8$ below $4000$, and its atlas is located. The $22$-digit prime $7510085481569082811681$ of item~3 is proved prime here by a Pocklington chain ($p-1=2^5\cdot3\cdot5\cdot q_1$, $q_1-1=2\cdot3\cdot5\cdot351707\cdot q_2$, $q_2-1=2\cdot3\cdot17^2\cdot5051\cdot169307$, witnesses below $12$ at every level, the leaves by trial division). Item 23 is audited in part (its certificates; the deduction read in part); items 11--13, 15 and 20 use imported material (Section~\ref{sec:overlaps}); the remaining 17 items, including these five, are located. The results bench R01--R10 and the arguments X1--X7 are checked in \note{43a\_} §3; one source passage error (X5: $289^3\equiv169$, not $-1$, modulo $840$) is recorded by the workbench itself.

\subsection{The \texorpdfstring{$(3,4,\infty)$}{(3,4,infinity)} monodromy covers}\label{ssec:es-s6}

A package of the workbench (ES-09) imports three integer matrices $A_1,A_2,A_\infty$ with $A_1^3=A_2^4=A_1A_2A_\infty=I$ from the $S^6$ manuscript \cite{AlpogeS6}, written out in the package's \texttt{src/core.tex}, and nothing else from it. They are the duals $A_j=(T_j^{-1})^{\mathsf T}$ of the local monodromies $T_1,T_2$ of the $S^6$ period system, acting on the dual lattice $\Lambda=\operatorname{Hom}(V,\ZZ)$ with dual basis $(\hat\gamma,\hat u,\hat w,\hat\delta)$ \cite{SixSpherePaper}. The package uses the induced affine action on the hyperplane
\[
H_D=\{\lambda\in\Lambda/D\Lambda:\ \text{the $\hat\gamma$-coordinate of $\lambda$ is }1\}\cong(\ZZ/D)^3,
\]
and builds from its orbits covers of the projective line branched at three points.
It presents the cover as a complete failure object: the equation fails at $p$ if and only if a denominator-labelled cover has no section. That equivalence is correct, and it is Theorem~\ref{thm:shell} in cover-theoretic form: a section exists exactly when the tail denominator $R_a/\gcd(R_a,4u+1)$ or $R_a/\gcd(R_a,u+a)$ equals $1$.

For odd levels $D$, which is the package's standing hypothesis and all that the application to the equation uses (there $D$ divides the odd number $R_a$), the package's orbit and genus statements are correct. The action is transitive when $3\nmid D$, and the cover is then connected of genus $(5D^3-12D^2-17D+24)/24$. When $3\mid D$ it has two orbits, of sizes $D^3/9$ and $8D^3/9$, which give components of genera $(5D^3-12D^2-81D+216)/216$ and $(5D^3-12D^2-9D+27)/27$. At even levels these statements do not hold as stated: at $D=2$ the orbits have sizes $2$ and $6$ (found by the referee of the audit notes).
\status{Audited: the proof read; the relations, the orbits and the genera recomputed from cycle counts for $D=1,3,\dots,15$. The equivalence with solvability is a re-encoding of Theorem~\ref{thm:shell}.}

The orbits at every level are determined by one invariant. Write $\lambda=\hat\gamma+b\hat u+c\hat w+d\hat\delta\in H_D$ and $e=\gcd(D,6)$.

\begin{theorem}[the orbits at every level]\label{thm:orbits}
Let the monodromy group $G=\langle T_1,T_2\rangle$ act on $H_D$ dually, through $A_1,A_2$. For every $D\ge1$ its orbits coincide with those of the full stabilizer $G_{\ZZ}$ of the isotropic vector $\gamma$ in the symplectic group of the period lattice, acting in the same way. The orbit of $\lambda$ is
\[
\{\lambda'\in H_D:\ \gcd(b',c',e)=\gcd(b,c,e)\},
\]
of size $D^3\cdot\#\{v\in(\ZZ/e)^2:\ \gcd(v,e)=\gcd(b,c,e)\}/e^2$.
\end{theorem}
\status{New in version~4. Proved in the companion paper on the $S^6$ period system \cite[Theorem~6.1]{SixSpherePaper}, which describes $G$ as the kernel of a character $\Phi$ of order $12$ on $G_{\ZZ}$; the orbit argument is repeated below. Checked as full orbit partitions for every $D\le14$ by the referee script \note{ver52\_referee\_claims.py}, for every $D\le24$ by the referee's own scripts, and for $D=2^k$, $k\le5$, by \note{es09\_even\_levels.py}.}
\begin{proof}
We use from \cite[§§2--3]{SixSpherePaper}: every element of $G_{\ZZ}$ is uniquely $L_Mh(\alpha,\beta,s)$ with $M\in SL_2(\ZZ)$ and $h(\alpha,\beta,s)$ in a Heisenberg group of integer triples; $G$ is the kernel of the homomorphism $\Phi(L_Mh(\alpha,\beta,s))=\operatorname{ab}(M)+s-6(\alpha^2+\alpha\beta+\beta^2)\bmod12$, where $\operatorname{ab}$ is the abelianisation of $SL_2(\ZZ)$ onto $\ZZ/12$ with $\operatorname{ab}\begin{psmallmatrix}1&1\\0&1\end{psmallmatrix}=1$; and $G_{\ZZ}=G\langle t_\gamma\rangle$, where the transvection $t_\gamma=h(0,0,1)$ is central with $\Phi(t_\gamma)=1$.

\emph{The dual action.} The dual action of $h(\alpha,\beta,s)$ is $(b,c,d)\mapsto(b+6\beta,\,c-6\alpha,\,d-s-\alpha b-\beta c)$, and $L_M$ acts on $(b,c)$ by $M^{-\mathsf T}$ and fixes $d$.

\emph{The orbits of $G_{\ZZ}$.} The coordinate $d$ can be moved freely, by $s$, and $(b,c)$ matters only modulo $6(\ZZ/D)^2$. Since $(\ZZ/D)^2/6(\ZZ/D)^2\cong(\ZZ/e)^2$ and $SL_2(\ZZ)$ maps onto $SL_2(\ZZ/e)$, the $G_{\ZZ}$-orbits correspond to the $SL_2(\ZZ/e)$-orbits on $(\ZZ/e)^2$, which are classified by the content $\gcd(v,e)$.

\emph{$G$ has the same orbits.} Since $G$ is normal and $G_{\ZZ}=G\langle t_\gamma\rangle$ with $t_\gamma$ central, it suffices to show $t_\gamma\lambda\in G\lambda$ for every $\lambda$; if this holds at $\lambda$, it holds at $k\lambda$ for every $k\in G_{\ZZ}$, because $t_\gamma k\lambda=kt_\gamma\lambda\in kG\lambda=Gk\lambda$. Every point is $G_{\ZZ}$-equivalent to one with $c=0$, because $SL_2(\ZZ)$ maps onto $SL_2(\ZZ/D)$, which moves every vector of $(\ZZ/D)^2$ to one of the form $(g,0)$. At a point with $(b,c)=(g,0)$, take $M=(T^{\mathsf T})^{-1}$ with $T=\begin{psmallmatrix}1&1\\0&1\end{psmallmatrix}$. Then $M^{-\mathsf T}=T$ fixes $(g,0)$, so $L_M$ fixes $\lambda$; and $M=ST S^{-1}$ with $S=\begin{psmallmatrix}0&-1\\1&0\end{psmallmatrix}$, so $\Phi(L_M)=\operatorname{ab}(T)=1$. Hence $t_\gamma L_M^{-1}$ lies in $G$ and moves $\lambda$ exactly as $t_\gamma$ does.
\end{proof}

\begin{corollary}\label{cor:orbits}
For $D=2^k$, $k\ge1$, there are exactly two orbits, of sizes $D^3/4$ and $3D^3/4$. For odd $D$ there is one orbit if $3\nmid D$, and there are two, of sizes $D^3/9$ and $8D^3/9$, if $3\mid D$. For $6\mid D$ there are four orbits; at $D=6$ their sizes are $6$, $18$, $48$ and $144$, and at $D=12$ they are $48$, $144$, $384$ and $1152$. In general the orbits at $D=2^km$ with $m$ odd are the products of those at $2^k$ and at $m$.
\end{corollary}
\begin{proof}
Count the vectors of $(\ZZ/e)^2$ by content. For $e=2$ the contents $2$ and $1$ occur $1$ and $3$ times; for $e=3$, $1$ and $8$ times; for $e=6$, the contents $6,3,2,1$ occur $1,3,8,24$ times. Multiply by $D^3/e^2$. The product statement is the Chinese remainder theorem, since $e=\gcd(2^k,2)\gcd(m,3)$.
\end{proof}

At the levels $D=2^k$ with $1\le k\le5$ the genera of the two components are $(0,0)$, $(1,3)$, $(17,53)$, $(177,537)$ and $(1569,4721)$ (found by the referee pass of version~1 and recomputed here from cycle counts). By Theorem~\ref{thm:orbits} the components of these covers at every level are indexed by the invariant $\gcd(b,c,\gcd(D,6))$; their genera at all levels are Question~\ref{q:genera}.

\section{The author's notes, April--June 2026}\label{sec:notes}

At the author's request this paper also treats the author's notes on the equation, written from April to June 2026. There are two collections:
\begin{itemize}
\item one of twenty \TeX{} notes and one 52-page paper;
\item the Erdős--Straus part of a second collection, about fifty-five further notes, among them the source note of the terminal core.
\end{itemize}
They are cited by title and theorem number.

The notes were read in six independent passes by separate Claude instances. Each pass tabulated every numbered statement and checked it with its own programs. Every result proved in this section was then re-derived and re-checked with this paper's own programs (Appendix~\ref{app:verification}). The statements that rest only on a reading pass, on the referee of version~3, or on a cited source say so. Drafts that the author's latest record supersedes are treated only for the results they contain that were verified.

\paragraph{What the notes contain.}
Their Erdős--Straus arithmetic consists of the shell calculus of Section~\ref{sec:core} and the classical identity families of Rosati, Yamamoto, Elsholtz--Tao, Salez and Lopez. To it they attach structural layers, which are correct where they were checked:
\begin{itemize}
\item the Niemeier lattice $A_5^4D_4$, which is also archive Theorem~25.8, and its umbral group $GL_2(3)$, the automorphism group of the glue code (checked by a reading pass);
\item the Hauptmodul of $\Gamma_0(6)$ and the Monster class $6B$;
\item Ogg's primes;
\item Lorentzian and Apollonian geometry;
\item Cayley--Dickson algebras;
\item the split-zero semiring, where one isomorphism needs an additional relation (Proposition~\ref{prop:dossier});
\item mock theta functions, where the notes' identification with the minimal model $M(2,7)$ holds in the form of Theorem~\ref{thm:m27}.
\end{itemize}
None of these structures takes a prime, a shell or a certificate as input, so a statement about solvability drawn from them needs a typed map from the shell data to them; Section~\ref{sec:overlaps} states, for each, what is established about such a map.

\paragraph{What this section records.}
\begin{itemize}
\item what the notes add to Sections~\ref{sec:core} and~\ref{sec:families}: the exact form of the terminal core, a sieve of quadratic-residue conditions on a counterexample, and facts about single shells, identities and locks;
\item corrections to the notes, each with its proof or counterexample;
\item the mock theta functions of the notes, with the theorem that identifies their shadow with the modular representation of $M(2,7)$.
\end{itemize}

\subsection{The terminal core}\label{ssec:termcore}

Put $M_{23}=840\cdot11\cdot19\cdot23=4{,}037{,}880$. A \emph{base class} is a class $n\equiv r\pmod{840}$ with $r\in S_{840}=\{1,121,169,289,361,529\}$, together with units modulo $11$, $19$ and $23$. There are $23{,}760$ base classes.

The source note, \emph{Finite central-cover reductions for residual Erdős--Straus shells}, calls $(R,m,u)$ an \emph{elementary $j=1$ certificate} when:
\begin{itemize}
\item $R\equiv3\pmod4$ and $u\mid m^2$;
\item $n\equiv-R\pmod{4m}$, $n\equiv-4u\pmod R$, and $\gcd(n,R)=1$.
\end{itemize}
Then $m\mid a=(n+R)/4$, $u\mid a^2$ and $u\equiv-a\pmod R$, so the shell is occupied through the middle channel $M_a$. The primes that may divide $m$ are those of $\{2,3,5,7,11,19,23\}$ not dividing $R$ whose divisibility of $a$ the class forces.

The note makes two computations:
\begin{enumerate}[label=(\roman*)]
\item The $32$ shells $R\mid M_{23}$ with $R\equiv3\pmod4$ cover all but $4951$ base classes. This is a covering statement.
\item Over all $R$ built from $3,5,7,11,19,23$, it tests the certificate congruence modulo $\gcd(M_{23},R)$, so that a class counts as hit when some lift of it modulo $\operatorname{lcm}(M_{23},R)$ is certified. Exactly $2970$ classes are then never hit, namely $S_{840}\times Q_{11}\times Q_{19}\times Q_{23}$, where $Q_q$ is the group of non-zero squares modulo $q$.
\end{enumerate}
The note calls this set the \emph{support-zero core}. It uses the set to bound from below the residue of any obstruction built from such certificates.

\begin{lemma}\label{lem:j1square}
Let $(R,m,u)$ be an elementary $j=1$ certificate with $\gcd(m,R)=1$, and let $n$ satisfy its congruences. Put $\varepsilon_\ell(n)=(n/\ell)$ for odd primes $\ell$. Put $\varepsilon_2(n)=1$ if $n\equiv1\pmod8$ and $\varepsilon_2(n)=-1$ if $n\equiv5\pmod8$. Then the Jacobi symbol satisfies
\[
\Big(\frac nR\Big)=-\prod_{\ell\mid u}\varepsilon_\ell(n)^{v_\ell(u)}.
\]
In particular no $n$ in the class is a unit square modulo $\operatorname{lcm}(4m,R)$.
\end{lemma}
\begin{proof}
From $n\equiv-4u\pmod R$ we get $(n/R)=(-1/R)(4/R)(u/R)=-(u/R)$.

For an odd prime $\ell\mid m$, reciprocity with $R\equiv3\pmod4$ gives $(\ell/R)=(-R/\ell)$. This equals $(n/\ell)$ because $n\equiv-R\pmod\ell$.

If $2\mid m$, then $n\equiv-R\pmod8$. So $(2/R)=1$ exactly when $R\equiv7\pmod 8$, that is, when $n\equiv1\pmod8$.

If $n$ were a square modulo $\operatorname{lcm}(4m,R)$, every $\varepsilon_\ell(n)$ with $\ell\mid m$ would be $1$. Then $(n/R)=-1$, which is impossible for a square.
\end{proof}

\begin{theorem}[the terminal core]\label{thm:termcore}
\begin{enumerate}[label=(\alph*)]
\item A base class has no elementary $j=1$ certificate with $R$ built from $3,5,7,11,19,23$ on any lift if and only if it is a square modulo $11$, $19$ and $23$. So the support-zero core is exactly $S_{840}\times Q_{11}\times Q_{19}\times Q_{23}$, which has $2970$ classes.
\item The base classes that no single identity among Salez's seven modular equations \cite{Salez} with modulus dividing $M_{23}$ covers number $3520$. They are the $2970$ square classes and $550$ classes that are non-squares modulo $11$, $19$ or $23$.
\item So the classes left open at level $M_{23}$ strictly contain the core. For example the prime $3361\equiv1\pmod{840}$, $\equiv6\pmod{11}$, lies outside the core, and none of the $64{,}149$ families with modulus dividing $M_{23}$ covers its class. The primes $18{,}481$ and $2{,}840{,}041$ behave in the same way.
\end{enumerate}
\end{theorem}
\status{(a) The \emph{if} direction is proved here. The \emph{only if} direction is a finite computation: the source note's script and an independent program of this paper both give $4951$ and then $2970$. (b) A finite computation, calibrated at Salez's modulus $G_7=840\cdot11\cdot13\cdot17\cdot19\cdot23$, where it leaves $147{,}348$ classes, the number Salez reports. Salez states that his seven equations form a complete set of modular equations. The primes in (c) are of course solvable, for instance $4/3361=1/841+1/110139970+1/950330$.}
\begin{proof}
(a), \emph{if}. Let $n$ lie over a square class.
\begin{itemize}
\item $S_{840}$ is the group of unit squares modulo $840$, so $n$ is a quadratic residue modulo every prime $\ell\in\{3,5,7,11,19,23\}$.
\item Also $n\equiv1\pmod8$.
\item So $(n/R)=1$ for every $R$ built from these primes.
\item The primes that may divide $m$ lie in $\{2,3,5,7,11,19,23\}$, so Lemma~\ref{lem:j1square} forces $(n/R)=-1$ for any certificate.
\end{itemize}
Hence no lift is certified. The remaining statements are the computations in the status line.
\end{proof}

\begin{remark}\label{rem:termcore}
The notes' later papers work with the core as the set of hard classes that remain: the 52-page \emph{Terminal residual coordinates for the Erdős--Straus problem} (26 May 2026) and its June replacement \emph{Finite support algebra, cubic middle geometry, and orientation descent}. Theorem~\ref{thm:termcore} locates the core inside the set that single identities leave open. At level $M_{23}$ that set has $3520$ classes. The core is the part of it that quadratic reciprocity forces, and the $550$ further classes are the part that remains from using one identity per class.

In the vocabulary of the split-zero semiring (Section~\ref{ssec:splitzero}):
\begin{itemize}
\item the $2970$ core classes are \emph{unsupported};
\item the $550$ further classes are \emph{supported but not covered}: some lift of each is certified, but no identity of Salez's seven families holds on the whole class.
\end{itemize}
Salez's statement that his seven equations are complete concerns identity families of his linear type; the archive's audit of Salez (§13.2) records this scope. Whether combinations of identities cover the $550$ classes at some finite level is Question~\ref{q:550}.

By Dirichlet's theorem the core carries $2970/\varphi(M_{23})=1/256$ of all primes, and the classes left open at level $M_{23}$ carry $3520/760{,}320=1/216$.
\end{remark}

\subsection{Quadratic-residue conditions on a counterexample}\label{ssec:qrconditions}

Throughout this subsection, $p\equiv1\pmod{24}$ is prime. A certificate is a shell $R\equiv3\pmod4$ with $\gcd(R,pa)=1$ and a divisor $u\mid a^2$ with $R\mid4u+1$ (channel $E$) or $R\mid u+a$ (channel $M$). Theorem~\ref{thm:shell} then gives the solution.

\begin{theorem}[eight shifts]\label{thm:eightshifts}
Let $q$ be an odd prime dividing
\[
F(p)=(p+1)(p+2)(p+3)(p+4)(p+7)(p+8)(2p+1)(3p+1)
\]
with $(p/q)=-1$. Then the shell and divisor of Table~\ref{tab:eightshifts} give a certificate. Equivalently, a counterexample $p\equiv1\pmod{24}$ is a quadratic residue modulo every odd prime factor of $F(p)$.
\end{theorem}

\begin{center}\small
\begin{longtable}{@{}lllll@{}}
\caption{The certificates of Theorem~\ref{thm:eightshifts}, with $a=(p+R)/4$.}\label{tab:eightshifts}\\
\toprule
shift & $q$ with $(p/q)=-1$ & shell $R$ & $u$ & channel\\
\midrule
$p+1$ & $q\equiv3\pmod4$ & $q$ & $a$ & $E$\\
$p+4$ & $q\equiv3\pmod4$ & $q$ & $1$ & $M$\\
$p+2$ & $q\equiv5,7\pmod8$ & $q$ if $q\equiv7$, $3q$ if $q\equiv5\pmod 8$ & $a/2$ & $E$\\
$p+8$ & $q\equiv5,7\pmod8$ & the same & $2$ & $M$\\
$2p+1$ & $q\equiv5,7\pmod8$ & the same & $2a$ & $E$\\
$p+3$ & $q\equiv2\pmod3$ & $3$ & $q$ & $E$\\
$p+7$ & $q\equiv5$, $6$, $3\pmod7$ & $7$ & $q$, $aq$, $2aq$ & $E$, $M$, $M$\\
$3p+1$ & $q\equiv2\pmod3$ & $(4q+1)/3$ & $q$ & $E$\\
\bottomrule
\end{longtable}
\end{center}

\status{Proved here.
\begin{itemize}
\item \emph{Recorded rows.} $p+1$ and $p+4$ are archive Theorem~24.50. $p+3$ and $p+7$ are Propositions~\ref{prop:r3} and~\ref{prop:r7}, read through reciprocity. $3p+1$ is archive Theorem~24.171 with $k=3$.
\item \emph{Salez's charts.} Each of the rows $p+1$, $p+2$, $2p+1$, $p+4$, $p+8$ is the union, over the free divisor $R$, of one of Salez's charts \cite{Salez} with small fixed constants: (15b) with $(B,C)=(1,1),(1,2),(2,1)$ and (14c) with $(B,D)=(1,1),(1,2)$.
\item \emph{Not in the archive.} The quadratic-residue conditions from $p+2$ and $2p+1$, and the stated condition from $p+8$, are not stated in the archive, in versions~1--2 of this paper, or in \note{43\_}/\note{43a\_}. No literature search was made.
\item \emph{Checks.} Every certificate for every prime $p\equiv1\pmod{24}$ below $10^7$ was built and verified: $780{,}700$ of them, none bad. $247$ primes survive all eight shifts.
\end{itemize}}
\begin{proof}
\emph{The characters.} For $q\mid p+k$ with $q\nmid k$ we have $(p/q)=(-k/q)$; similarly $(p/q)=(-2/q)$ for $q\mid 2p+1$ and $(p/q)=(-3/q)$ for $q\mid3p+1$. These characters are $-1$ exactly on the classes of $q$ listed in the table; for $p+7$ this uses $(-7/q)=(q/7)$.

\emph{Each row is a certificate.} For each row the stated divisibility is an identity:
\begin{itemize}
\item $4a+1=p+R+1$;
\item $4(a+1)=p+R+4$;
\item $2(2a+1)=p+R+2$;
\item $4(a+2)=p+R+8$;
\item $8a+1=2(p+R)+1$;
\item $4q+1\equiv q+1\pmod3$;
\item for $p+7$: $4q+1\equiv0$, $a(q+1)$ and $a(2q+1)$ modulo $7$.
\end{itemize}

\emph{The rows $p+2$, $p+8$ and $2p+1$.} Since $p\equiv1\pmod3$, $3$ divides each of these shifts, so $R$ divides the shift. Here $q\neq3$ because $(p/3)=1$. Also $R\equiv7\pmod8$, so $8\mid p+R$ and $a$ is even.

\emph{The row $p+7$.} $8\mid p+7$, so $a$ is even, and $q\mid a$.

\emph{The row $3p+1$.} Put $N_3=(3p+1)/4$, which is odd and $\equiv1\pmod3$. Then $h=N_3/q\equiv2\pmod3$. With $a=q(h+1)/3$ one checks $4a-p=(4q+1)/3=R$.

\emph{Coprimality.} A common prime of $R$ and $a$ divides $4a-R=p$, and $p\nmid R$ in every row.
\end{proof}

At the strict record $p_*=8{,}803{,}369$, whose least occupied shell is $107$, two shifts give four certificates:
\begin{itemize}
\item $p_*+2=3\cdot223\cdot13159$, where $223\equiv7$ and $13159\equiv7\pmod8$, giving the shells $223$ and $13{,}159$;
\item $2p_*+1=3\cdot677\cdot8669$, where $677\equiv5$ and $8669\equiv5\pmod8$, giving the shells $2031$ and $26{,}007$.
\end{itemize}
None of the other six shifts gives a certificate at $p_*$. The least primes certified by $p+2$ alone, by $2p+1$ alone and by $p+8$ alone are $2521$, $9601$ and $33{,}961$.

\begin{theorem}[$p+5$]\label{thm:pplus5}
Let $N=(p+5)/2$, and let $p\equiv9\pmod{40}$, or $p\equiv121\pmod{840}$, or $p\equiv1,361\pmod{840}$ with $p\equiv4\pmod9$. Suppose some prime $q\mid p+5$ has $(p/q)=-1$. Then $N$ has a divisor $d\equiv-p\pmod{20}$ or a divisor $h\equiv-1\pmod{20}$, and either gives a certificate by Proposition~\ref{prop:locks} with $\mu=5$.

So a counterexample in these classes, which contain the hard classes $121$, $169$, $289$ and $529$, is a quadratic residue modulo every odd prime factor of $p+5$.
\end{theorem}
\status{Proved here. The case $p\equiv9\pmod{40}$ is in the note \emph{Full divisor locks, QR subtori, and survivor idempotents} in substance (its Theorem~3.7(ii), Corollary~3.8(ii) and Theorem~4.2). The extension to the classes $121$, $1$ and $361$ was found by a reading pass. The theorem was checked on all $41{,}419$ primes $p\equiv9\pmod{40}$ below $10^7$ and on $5688$ primes of the other classes. At $p=12601\equiv1\pmod{840}$, with $p\equiv1\pmod9$, it fails: $p+5=2\cdot3\cdot11\cdot191$ has two non-residue factors, yet $N$ has no divisor in either class.}
\begin{proof}
For $q\mid N$ we have $(p/q)=(-5/q)$. This equals $1$ exactly when $q\equiv1,3,7,9\pmod{20}$, a subgroup $H$ of index $2$.

\emph{The case $p\equiv9\pmod{40}$.} Here $N\equiv7\pmod{20}$, and the targets are $11$ and $19$. Let $q\notin H$ divide $N$:
\begin{itemize}
\item if $q\equiv11$ or $19$, then $q$ itself is the divisor;
\item if $q\equiv13$, then $N/q\equiv19$;
\item if $q\equiv17$, then $N/q\equiv11$.
\end{itemize}

\emph{The case $p\equiv1\pmod{40}$.} Here $N\equiv3\pmod{20}$, and a divisor $\equiv17$ or $19$ suffices, because its cofactor is then $\equiv19$. Let $q\notin H$ divide $N$:
\begin{itemize}
\item $3\mid N$ always, so $q\equiv13$ gives $3q\equiv19$;
\item $7\mid N$ when $p\equiv121\pmod{840}$, so $q\equiv11$ gives $7q\equiv17$;
\item $9\mid N$ when $p\equiv4\pmod9$, so $q\equiv11$ gives $9q\equiv19$.
\end{itemize}

\emph{Conversely}, the targets lie outside $H$, so a divisor in a target class has a prime factor outside $H$.
\end{proof}

\begin{theorem}[the visible-box principle]\label{thm:visiblebox}
Let $R\equiv3\pmod4$ be prime and $c$ a class of hard primes modulo $10080=2^5\cdot3^2\cdot5\cdot7$. Let
\[
\mathrm{Box}(c)=\Big\{\prod_{f\in\{2,3,5,7\}}f^{i_f}\bmod R:\ 0\le i_f\le2e_f\Big\},
\]
where $e_f$ is the least value of $v_f((p+R)/4)$ that $c$ forces. If $\mathrm{Box}(c)$ contains every non-zero square modulo $R$, then for every prime $p\in c$ the shell $R$ is occupied exactly when some odd prime $q\mid p+R$ has $(p/q)=-1$.
\end{theorem}
\status{Proved here; found by a reading pass. The archive's census refinements (Theorems~24.228 and~24.233) already give certificates in the shells $R=31$, $47$ and $71$ on parts of these classes; the equivalence is not stated there. The search up to $R=700$ is complete, because $|\mathrm{Box}(c)|\le7\cdot5\cdot3\cdot3=315$. Among the $72$ hard classes modulo $10080$ it finds $72$ good classes for $R=3$ and $R=7$, $48$ for $R=11$, $12$ for $R=19$, $48$ for $R=23$, and $30$, $20$, $4$, $12$, $1$ for $R=31$, $47$, $59$, $71$, $311$. The last five shells are new. All $90{,}897$ tests with $p<10^7$ agree, with $50{,}545$ certificates verified.}
\begin{proof}
Every element of $\mathrm{Box}(c)$ is the residue of a divisor of $a^2$.

If $q\mid a$ is a prime with $(q/R)=-1$, then $(p/q)=(-R/q)=(q/R)=-1$. Also $q\cdot\mathrm{Box}(c)$ is the non-square coset, so the divisors reach every unit, including the target $-4^{-1}$ of channel $E$.

The condition $(p/R)=-1$ needs no separate treatment. Let $m$ be the odd part of $p+R$. Since $p\equiv1\pmod8$, reciprocity gives $(p/m)=(m/p)=(R/p)(2/p)^{-v_2(p+R)}=(p/R)$. So if $(p/R)=-1$, some prime factor $q$ of $m$ has $(p/q)=-1$.

Conversely, suppose that no odd prime $q\mid p+R$ has $(p/q)=-1$. For an odd prime $\ell\mid a$, reciprocity gives $(\ell/R)=(-R/\ell)=(p/\ell)=1$; the prime $2$ divides $a$ only when $R\equiv7\pmod8$, and then $(2/R)=1$. So every prime factor of $a$ is a square modulo $R$, and Proposition~\ref{prop:jacobi} shows that the shell is empty.
\end{proof}

\begin{proposition}[$\mu=17$]\label{prop:mu17}
Let $p\equiv361$ or $529\pmod{840}$, $p\equiv1\pmod9$ and $p\bmod17\in\{8,15,16\}$. Then $(p+17)/2$ has a divisor $d\equiv-p$ or $h\equiv-1\pmod{68}$, giving a certificate by Proposition~\ref{prop:locks}, if and only if some prime $q\mid p+17$ has $(p/q)=-1$.
\end{proposition}
\status{Proved here; found by a referee pass. Checked on all $433$ such primes below $10^7$.}
\begin{proof}
Put $N=(p+17)/2$. Then $63\mid N$. The character $(p/q)=(-1/q)(q/17)$ lives modulo $68$, and both targets $-p$ and $-1$ lie outside its kernel. A finite check of the six classes shows that for every non-kernel class $r$ the sets $r\cdot W$ and $N(rW)^{-1}$, with $W=\{1,3,7,9,21,63\}$, meet the targets modulo $68$. The converse holds as in Theorem~\ref{thm:pplus5}.
\end{proof}

The primes $p\equiv1\pmod{24}$ that satisfy the conditions of Theorems~\ref{thm:eightshifts}, \ref{thm:pplus5} and~\ref{thm:visiblebox} and Proposition~\ref{prop:mu17} are counted in Table~\ref{tab:survivors}.

\begin{center}\small
\begin{longtable}{@{}lrrr@{}}
\caption{Primes $p\equiv1\pmod{24}$ satisfying the conditions of Theorems~\ref{thm:eightshifts}, \ref{thm:pplus5} and~\ref{thm:visiblebox} and Proposition~\ref{prop:mu17}, cumulatively. The column $10^8$ is from a referee pass.}\label{tab:survivors}\\
\toprule
conditions & below $10^6$ & below $10^7$ & below $10^8$\\
\midrule
hard classes & $2370$ & $20{,}513$ & $179{,}468$\\
and the eight shifts & $54$ & $247$ & $1431$\\
and $p+5$ on its classes & $29$ & $146$ & $816$\\
and $\mu=17$ & $29$ & $145$ & $806$\\
and the visible-box shells & $7$ & $31$ & $165$\\
\bottomrule
\end{longtable}
\end{center}

The first primes that satisfy all the conditions of the table are $43{,}201$, $65{,}521$ and $196{,}561$.

The referee of version~3 found certificates of three further shapes, on parts of the hard classes. Let $R\equiv3\pmod4$ with $a=(p+R)/4$ and $\gcd(R,pa)=1$, and let $s,t\ge1$.
\begin{itemize}
\item $u=s$ is an $M$-certificate when $s\mid a^2$ and $R\mid p+4s$, since $4(s+a)=p+4s+R$.
\item $u=a/t$ is an $E$-certificate when $t\mid a$, $\gcd(t,R)=1$ and $R\mid p+t$, since $t(4a/t+1)=p+R+t$.
\item $u=ta$ is an $E$-certificate when $t\mid a$ and $R\mid tp+1$, since $4ta+1=t(p+R)+1$.
\end{itemize}
Applied to the shifts $p+6$ ($t=6$), $p+12$, $p+16$, $p+20$, $p+24$, $p+36$ ($s=3,4,5,6,9$), $5p+1$ and $6p+1$ ($t=5,6$), they give conditions of the same kind as those above. The referee verified $46{,}441$ such certificates below $10^7$, and its script was re-run here. Below $10^7$, exactly $23$ primes $p\equiv1\pmod{24}$ have no certificate of any of the kinds of this subsection. The prime $43{,}201$, for example, has one from $p+24$: $4/43201=1/10914+1/1036824+1/1885982856$.

The survivors are rare, and already the first condition removes almost all primes.

\begin{proposition}[density zero]\label{prop:densityzero}
The primes $p\equiv1\pmod{24}$ that satisfy the condition of the row $p+1$ of Theorem~\ref{thm:eightshifts} have relative density zero among the primes $p\equiv1\pmod{24}$. A fortiori, so do the primes that satisfy all the conditions of this subsection.
\end{proposition}
\status{Proved here; new in version~4.}
\begin{proof}
If a prime $q\equiv3\pmod4$ divides $p+1$, then $(p/q)=(-1/q)=-1$, so the row $p+1$ gives a certificate. Hence a prime satisfying the condition has $p\not\equiv-1\pmod q$ for every prime $q\equiv3\pmod4$. (The prime $q=3$ never divides $p+1$ when $p\equiv1\pmod{24}$.) Let $q_1,\dots,q_k$ be the first $k$ primes $\equiv3\pmod4$ greater than $3$. By Dirichlet's theorem for the progressions modulo $24q_1\cdots q_k$, the primes $p\equiv1\pmod{24}$ with $p\not\equiv-1\pmod{q_i}$ for all $i\le k$ have relative density $\prod_{i\le k}\bigl(1-1/(q_i-1)\bigr)$ among the primes $p\equiv1\pmod{24}$. Since $\sum1/q$ over the primes $q\equiv3\pmod4$ diverges, this product tends to $0$ as $k$ grows, and the relative upper density of the set in question is at most every such product.
\end{proof}

Whether infinitely many primes survive all the conditions is Question~\ref{q:qrsurvivors}. A negative answer would prove the conjecture at all but finitely many primes.

\begin{remark}[the Collatz reading]\label{rem:collatz}
For $p\equiv1\pmod8$ the number $(3p+1)/4$ is the odd Collatz successor of $p$. So a counterexample $p\equiv1\pmod{24}$ has an odd Collatz successor built only from primes $\equiv1\pmod3$.

The odd Collatz predecessors of an odd $n$ with $3\nmid n$ satisfy $\Pi_{k+1}=4\Pi_k+1$, which is the notes' map $a\mapsto4a+1$. So $\Pi_k\equiv5\pmod8$ for $k\ge1$. Every prime $P\equiv5\pmod8$ is solved by the shell $R=3$ with $u=2$. Hence only the first member of a predecessor chain can be a hard prime.

This is the classical Syracuse predecessor structure \cite{Lagarias}, and it links this workbench with the Collatz workbench \cite{WorkbenchReader}. The link uses single steps of the Collatz map.
\end{remark}

\subsection{Single shells}\label{ssec:singleshells}

\begin{proposition}[non-coprime shells]\label{prop:noncoprime}
Let $p$ be prime and $R\equiv-p\pmod4$. Then $\gcd(R,pa)>1$ exactly when $p\mid R$; write $R=kp$.
\begin{enumerate}[label=(\alph*)]
\item If $p\nmid k$, the one-congruence criterion, that some $d\mid N^2$ has $d\equiv-N\pmod R$, is equivalent to occupancy, and every such $d$ gives integral $y$ and $z$.
\item If $p^2\mid R$, it is not. At $p=73$, $R=7\cdot73^2$, the divisor $d=4\cdot73^3$ satisfies the congruence and gives $y=60$ but $z=1920/73$, and the shell has no solution.
\end{enumerate}
\end{proposition}
\status{Proved here; found by a reading pass. Four of the notes state the criterion without the coprimality hypothesis: the defect-completion article (Lemma~2.1), \emph{Split-zero residual moonshine} (Theorem~3.1), the $A_8$ snowflake note (Theorem~2.1) and \emph{Terminal-core lock covers} (Theorem~2.1). The error does not affect the notes' applications, which are at hard primes with $R<3p^2$ (a reading pass): for $p\equiv1\pmod4$, $p^2\mid R$ forces $R\ge3p^2$. For $p\equiv3\pmod4$ it allows $R=p^2$.}
\begin{proof}
(a) Put $m=(k+1)/4$, so that $a=pm$ and $N=p^2m$. A divisor $d$ with $d\equiv-N\pmod{kp}$ is divisible by $p$; write $d=pd_1$, so that $k\mid d_1+pm$.
\begin{itemize}
\item If $p^3\nmid d_1$, then $d_1\mid(pm)^2$ is a certificate for the reduced shell $(k,pm)$, and the formulas agree.
\item If $d_1=p^3w$ with $w\mid m^2$, then $k\mid4p^2w+1$, so $(w/k)=(-1/k)=-1$.
\item But every prime $\ell\mid m$ has $(\ell/k)=1$: for odd $\ell$ because $k\equiv-1\pmod\ell$ and reciprocity applies; for $\ell=2$ because $k\equiv7\pmod8$. So $(w/k)=1$, a contradiction.
\end{itemize}
(b) Direct computation, including a search over all $y$.
\end{proof}

\begin{proposition}[ramified shells]\label{prop:ramified}
Let $(R,p)$ be a coprime shell.
\begin{enumerate}[label=(\alph*)]
\item If $R\mid p-1$, then $E_a=M_a$.
\item If $R\mid p^2-1$, then $u\mapsto a^2/u$ maps $E_a$ to itself, and $|E_a|$ is odd exactly when $R\mid p+1$.
\end{enumerate}
\end{proposition}
\status{Proved here; checked on all $947$ coprime shells with $R\mid p^2-1$ and $p\equiv1\pmod4$ below $2000$.}
\begin{proof}
(a) $4a\equiv p\equiv1\pmod R$, so $4(u+a)\equiv4u+1$.

(b) The complement sends the class $-4^{-1}$ to $-4a^2\equiv-p^2/4$, which is $-4^{-1}$ again when $p^2\equiv1\pmod R$. The only possible fixed point is $u=a$. Finally $a\in E_a$ exactly when $R\mid4a+1=p+R+1$.
\end{proof}

\begin{proposition}[the pair-count series]\label{prop:pairzeta}
Let $A_R(n)$ be the number of triples $(X,Y,Z)$ with $XYZ=n$, $\gcd(X,Y)=1$ and $R\mid X+Y$. For a coprime shell, $A_R(pa)$ is the number of ordered certificates. For $R\ge3$,
\[
\sum_{(n,R)=1}\frac{A_R(n)}{n^s}=\frac1{\varphi(R)}\sum_{\chi\bmod R}\chi(-1)\,\frac{\zeta_R(s)\,L(s,\chi)\,L(s,\bar\chi)}{\zeta_R(2s)},
\]
where $\zeta_R$ omits the primes dividing $R$.
\end{proposition}
\status{Proved here. This is Theorem~4.3 of the defect-completion article; $8757$ coefficients were checked by a referee pass. No positivity statement follows.}
\begin{proof}
Write $1_{X\equiv-Y}=\varphi(R)^{-1}\sum_\chi\chi(-1)\chi(X)\bar\chi(Y)$. The variable $Z$ contributes $\zeta_R(s)$. Möbius inversion over the common divisors of $X$ and $Y$ contributes $L(s,\chi)L(s,\bar\chi)/L(2s,\chi\bar\chi)$, and $\chi\bar\chi$ is principal.
\end{proof}

The divisor-pair note asks for finitely many central and edge templates covering $S_{840}$. As a congruence cover this is impossible, because no template class contains a square modulo its modulus.
\begin{itemize}
\item For the edge templates this is Proposition~\ref{prop:nosquare}.
\item A central template is the class $n\equiv-R\pmod{4c}$ with $s\mid c^2$, $\gcd(c,R)=1$, $R\equiv3\pmod4$ and $R\mid s+c$. If $-R$ were a unit square modulo $4c$, then:
\begin{itemize}
\item every odd prime $\ell\mid c$ would have $(-R/\ell)=1$, hence $(\ell/R)=1$;
\item if $2\mid c$, then $-R\equiv1\pmod8$, so $(2/R)=1$.
\end{itemize}
So $(c/R)=(s/R)=1$. But $R\mid s+c$ gives $(s/R)=(-c/R)=-1$. (A reading pass found this argument.)
\end{itemize}
By Dirichlet's theorem, primes $p\equiv1$ modulo the product of $840$ and the template moduli exist, and they escape every template.

\subsection{Identities, locks and types}\label{ssec:locks}

\begin{proposition}[Rosati's family]\label{prop:rosati}
Let $k\ge1$, $c\mid k$ and $m=4k-1$. If $n\equiv-c\pmod m$, then
\[
\frac4n=\frac1{kn}+\frac1{k(n+c)/m}+\frac1{kn(n+c)/(mc)}.
\]
\end{proposition}
\status{Proved here; Salez's chart (14a) with $(B,C,D)=(c,1,k/c)$. The case $k=22$, $c=11$ solves every $n\equiv76\pmod{87}$. The two \emph{direct $29$-channels} of \emph{Terminal residual coordinates} (Theorems~7.1--7.2), $n\equiv917\pmod{1276}$ and $n\equiv1605\pmod{2204}$, are correct. They are Salez's chart (15b) with $(B,C,F)=(1,29,11)$ and $(1,29,19)$. The first lies inside this family with $(k,c)=(80,40)$, $m=319$. The second lies inside chart (14a) with $(B,C,D)=(2,23,3)$, $m=551$. The notes' claim that these are the only direct channels holds only for lock divisors that are single primes from $\{5,11,19,23\}$. The lock with $\mu=5$ and $d=31$ solves $n\equiv429\pmod{620}$, a class containing the core prime $33{,}289$.}
\begin{proof}
Over the common denominator $kn(n+c)$ the three terms have numerators $n+c$, $mn$ and $mc$, which add up to $(1+m)(n+c)=4k(n+c)$. The denominators are integers because $m\mid n+c$ and $c\mid k$.
\end{proof}

\begin{proposition}[the full-divisor locks]\label{prop:locks}
Let $p$ be an odd prime, $\mu$ odd and $N=(p+\mu)/2$.
\begin{enumerate}[label=(\alph*)]
\item If $d\mid N$ is odd with $d\equiv-p\pmod{4\mu}$, then the shell $R=d$ carries the $E$-certificate $u=a/\mu$.
\item If $h\mid N$ with $h\equiv-1\pmod{4\mu}$, put $e=N/h$ and $t=(h+1)/(4\mu)$. Then $a=2\mu et$ and $R=\mu+2e$ satisfy $4a-p=R$, and $u=\mu^2t$ is an $M$-certificate.
\end{enumerate}
\end{proposition}
\status{Proved here. These are the locks of the notes, which a reading pass identifies with the faces $A=1$ of Type~I and $C=1$ of Type~II of \cite{ElsholtzTao}; lock (b) is Lopez's Type~B \cite{Lopez}. In the notes' proof of (a) the numerator should read $ps+p+\mu$, not $ps+p+\mu s$.}
\begin{proof}
(a) Here $\mu\mid a$, and $4u+1=(p+d+\mu)/\mu$ is divisible by $d$.

(b) $4a-p=2e(h+1)-2eh+\mu=R$ and $u+a=\mu tR$.
\end{proof}

\paragraph{Computations on these certificates.}
\begin{itemize}
\item Every hard prime below $10^7$ has a lock certificate. A reading pass found the same below $10^8$, where the hard primes number $179{,}468$. Not every prime has one: $409\equiv1\pmod{24}$, which is not hard, has none for any odd $\mu$, by an exhaustive search. It is the only prime $p\equiv1\pmod4$ below $3\cdot10^6$ without one (found by the referee of this version).
\item Every prime $p\equiv1\pmod{24}$ below $10^7$ has a Type~II solution, that is, an $M$-certificate, in a shell $R\le107$. The record $p_*$ is the only one that needs $107$. A reading pass found the same below $10^8$.
\item The \emph{complement-square} certificates $u=Q^2$ with $Q\mid a$ (Lopez's Type~C \cite{Lopez}) are not universal. The least prime $p\equiv1\pmod{24}$ without one is $97$, and the least hard prime without one is $3361$.
\end{itemize}

\subsection{Corrections}\label{ssec:noteserrors}

Besides the scope of the one-congruence criterion (Proposition~\ref{prop:noncoprime}), the numerator in the notes' proof of lock (a) (Proposition~\ref{prop:locks}) and the scope of the direct $29$-channels (Proposition~\ref{prop:rosati}), the checked corrections are the following. Each is stated with its proof or counterexample.
\begin{itemize}
\item \emph{Star--Kneser positivity} states a converse of Star--Kneser forcing. The archive records that it fails at $p_*$, $R=107$, and retracts it (§15.3; located).
\item Four notes identify the umbral group of $A_5^4D_4$, the automorphism group of the glue code, with $GL_2(3)$, by an argument that does not exclude the binary octahedral group. The conclusion holds: $GL_2(3)$ has $13$ involutions, while the binary octahedral group has exactly one element of order $2$ (a reading pass, and \note{misc\_checks.py}).
\item A decomposition of the shell count over the six hard classes, of the kind the notes call a Coxeter--Fourier decomposition, needs data finer than the class of $p$ modulo $840$, because the count is not a function of that class. With $A_R$ as in Proposition~\ref{prop:pairzeta} and $a=(p+11)/4$, $A_{11}(pa)$ takes the values $0,4,6,8,10$ at the primes $2521$, $5881$, $4201$, $15121$ and $13441$, all $\equiv1\pmod{840}$ (\note{misc\_checks.py}).
\item The map $(\ZZ/18)^\times\to S_{840}$ of the terminal papers is a bijection and not a homomorphism; the decomposition that uses it holds as stated (Proposition~\ref{prop:torsor}).
\item One isomorphism of the split-zero dossier, its Theorem~7.2, needs the relation $[\tau]\sim0$ (Proposition~\ref{prop:dossier}).
\item The Kneser theorem used is that of \cite{Kneser}, of 1955, not the density paper of 1953.
\end{itemize}

\subsection{The mock theta functions of the notes and the minimal model \texorpdfstring{$M(2,7)$}{M(2,7)}}\label{ssec:mock}

The notes attach to the cubic middle of the shell problem an explicit rank-three modular object. They follow the framework of 3d modularity \cite{CCFGH} and its revision \cite{CCKPS}. There the $\widehat Z$ invariants of plumbed $3$-manifolds are quantum modular, and their companions under orientation reversal are mock or mixed mock modular. For the Brieskorn sphere $\Sigma(2,3,7)$ the relevant objects are Ramanujan's order-$7$ mock theta functions \cite{CFS,CCKPS}.

The note \emph{An explicit mixed-mock Virasoro support realization} works with $m=42$ and the Atkin--Lehner group $K=\{1,6,14,21\}$. It computes the projector $P_7=P^+(6)P^+(14)P^+(21)P^-(42)$ on $\C[\ZZ/84]$ and its image basis
\begin{align*}
v_1&=e_1-e_{13}-e_{29}+e_{41}-e_{43}+e_{55}+e_{71}-e_{83},\\
v_5&=e_5+e_{19}+e_{23}+e_{37}-e_{47}-e_{61}-e_{65}-e_{79},\\
v_{11}&=e_{11}+e_{17}+e_{25}+e_{31}-e_{53}-e_{59}-e_{67}-e_{73},
\end{align*}
and writes regularised indefinite theta functions whose shadows lie on these three vectors.

First, the explicit mixed mock functions of the note are Ramanujan's order-$7$ mock theta functions, in Zagier's vector $M_7$, whose components are $q^{-1/168}F_0$, $q^{-25/168}F_1$ and $q^{47/168}F_2$ up to signs \cite{Zagier}. The $q$-expansions agree through $q^{40}$, up to one sign (a reading pass). For $m=30$ the same construction gives the order-$5$ functions. The first order-$7$ function is the invariant $q^{1/2}\widehat Z_0$ of the Brieskorn sphere $-\Sigma(2,3,7)$ \cite[eq.~(3.20)]{CFS}. So the notes' construction lands exactly on the mock theta functions of $3$-manifold invariants.

Second, the note identifies the shadow of this vector with the characters of the minimal model $M(2,7)$, whose central charge is $c=-68/7$ and whose conformal weights are $h=0,-\tfrac27,-\tfrac37$. The following theorem establishes this identification at the level of modular representations.

For $r\in\{1,5,11\}$ put
\[
G_r(\tau)=\sum_{\ell\in\ZZ}v_r(\ell\bmod84)\,\ell\,q^{\ell^2/168},
\]
the weight-$\tfrac32$ unary theta series on the three support vectors. For a vector $F$ of weight $k$ define matrices $\rho_F(S)$, $\rho_F(T)$ by $F(-1/\tau)=(-i\tau)^k\rho_F(S)F(\tau)$ and $F(\tau+1)=\rho_F(T)F(\tau)$. Let $\rho=\rho_G$, and let $\rho_{2,7}$ be the representation of $SL_2(\ZZ)$ on the characters $\chi_s=\chi^{(2,7)}_{1,s}$, $s=1,2,3$, given by the Rocha-Caridi formula \cite{RochaCaridi}. Let $\varepsilon$ be the multiplier of Dedekind's $\eta$, so that $\varepsilon(S)=1$ and $\varepsilon(T)=e(1/24)$ in this convention, and write $c_k=(2/\sqrt7)\sin(k\pi/7)$.

\begin{theorem}[the order-$7$ shadow carries the $M(2,7)$ representation]\label{thm:m27}
In the basis $(G_1,G_5,G_{11})$,
\[
\rho(S)=\begin{pmatrix}-c_1&c_2&c_3\\ c_2&c_3&-c_1\\ c_3&-c_1&c_2\end{pmatrix},\qquad \rho(T)=\operatorname{diag}\bigl(e(\tfrac1{168}),e(\tfrac{25}{168}),e(\tfrac{121}{168})\bigr).
\]
Let $P$ be the signed permutation sending $(G_1,G_5,G_{11})$ to $(\chi_2,\chi_3,-\chi_1)$. Then
\[
\overline{\rho(S)}=P\rho_{2,7}(S)P^{-1},\qquad \overline{\rho(T)}\,e(1/8)=P\rho_{2,7}(T)P^{-1}.
\]
Hence $\bar\rho\otimes\varepsilon^3\cong\rho_{2,7}$: the shadow of the order-$7$ mock theta vector, conjugated and twisted by the character of $\eta^3$, transforms exactly as the characters of $M(2,7)$.
\end{theorem}
\status{Proved here; new in version~4. The $S$-identity is verified exactly in $\Q(\zeta_{168})$ and numerically to $38$ digits (\note{exact\_shadow\_reps.py}, \note{m7\_vs\_M27.py}). Whether this isomorphism is stated in the literature was not checked.}
\begin{proof}
\emph{The $S$-matrix of $\rho$.} The theta decomposition gives
\[
\theta_{m,\mu}(-1/\tau,z/\tau)=\sqrt{\tau/2mi}\,e(mz^2/\tau)\sum_\nu e(-\mu\nu/2m)\,\theta_{m,\nu}(\tau,z)
\]
\cite[§5]{EichlerZagier}. Differentiating in $z$ at $z=0$ gives
\[
\theta^1_{m,\mu}(-1/\tau)=(-i\tau)^{3/2}\,\frac i{\sqrt{2m}}\sum_\nu e(-\mu\nu/2m)\,\theta^1_{m,\nu}(\tau)
\]
for the derivative series $\theta^1_{m,\mu}$. The span of $v_1,v_5,v_{11}$ is the image of the projector $P_7$, whose factors are Atkin--Lehner involutions; these commute with the finite Fourier transform above \cite{EichlerZagier,CDH}, so the transform of each $G_a$ lies in the span of $G_1,G_5,G_{11}$, as the direct evaluation below also confirms. Since the vectors $v_r$ are odd with disjoint supports, $\rho(S)_{ab}=(i/\sqrt{84})\sum_rv_a(r)e(-rb/84)$. This formula agrees with a direct evaluation of the series at sample points to $29$ digits.

\emph{The $S$-matrix of $\rho_{2,7}$.} The Rocha-Caridi formula gives $\rho_{2,7}(S)_{s\sigma}=(2/\sqrt7)(-1)^{s+\sigma}\sin(2\pi s\sigma/7)$.

\emph{The $S$-identity.} Both sides lie in $\Q(\zeta_{168})$, using $\sqrt7=-i\sum_{k=1}^6\bigl(\tfrac k7\bigr)\zeta_7^k$ and $\sqrt3=-i(\zeta_3-\zeta_3^2)$. The identity was verified exactly, by reduction modulo the cyclotomic polynomial $\Phi_{168}$.

\emph{The $T$-identity.} It follows from the exponents. The values $-r^2/168+21/168$ for $r=1,5,11$ are $20/168$, $-4/168$ and $68/168$ modulo $1$, that is $5/42$, $-1/42$ and $17/42$. These are the values of $h-c/24$ for $s=2,3,1$: $c/24=-17/42$, and $h_{1,s}=0,-\tfrac27,-\tfrac37$ for $s=1,2,3$.
\end{proof}

The note reaches the $M(2,7)$ exponents by its cubic phase transport $T\mapsto e(-1/24)T^3$. That map sends the labels $1,5,11$ to $-1/42$, $17/42$ and $5/42$: the correct set of exponents, but a different matching of labels to characters. With the note's matching, $1\mapsto\chi_3$, $5\mapsto\chi_1$ and $11\mapsto\chi_2$, the diagonal $S$-entries differ in absolute value: $\rho(S)_{11}=-c_1$, while $\rho_{2,7}(S)_{33}=c_3$. So no signed permutation and scalar realise that matching, and the pair $(\rho(S),e(-1/24)\rho(T)^3)$ does not satisfy the relation $(ST)^3=S^2$ (\note{m7\_vs\_M27.py}). Theorem~\ref{thm:m27} realises the identification through conjugation and the $\eta^3$ twist instead, and with it the note's claim that the explicit mixed mock functions carry a $c=-68/7$ character shadow holds at the level of modular representations, after conjugation and the $\eta^3$ twist. The note's third-order modular differential equation also checks: in the monic form $\vartheta^3f+\mu_1E_4\vartheta f+\mu_2E_6f=0$, with $\vartheta$ the Serre derivative, the indicial polynomial is $x(x-\tfrac16)(x-\tfrac13)+\mu_1x+\mu_2$, and the roots $\tfrac{17}{42},\tfrac5{42},-\tfrac1{42}$ force $\mu_1=\tfrac1{28}-\tfrac1{18}=-\tfrac5{252}$ and $\mu_2=\tfrac{85}{74088}$, the note's coefficients.

\begin{proposition}[the order-$5$ block carries the Fibonacci data]\label{prop:fibonacci}
Let $m=30$, $K=\{1,6,10,15\}$, and take the note's vectors
\[
w_1=e_1+e_{11}+e_{19}+e_{29}-e_{31}-e_{41}-e_{49}-e_{59},\qquad w_7=e_7+e_{13}+e_{17}+e_{23}-e_{37}-e_{43}-e_{47}-e_{53}.
\]
The weight-$\tfrac32$ series $G_r(\tau)=\sum_\ell w_r(\ell\bmod60)\,\ell\,q^{\ell^2/120}$, $r=1,7$, satisfy
\[
\rho(S)=\frac2{\sqrt5}\begin{pmatrix}\sin\frac\pi5&\sin\frac{2\pi}5\\ \sin\frac{2\pi}5&-\sin\frac\pi5\end{pmatrix},\qquad \rho(T)\,e(-1/8)=\operatorname{diag}\bigl(e(-\tfrac7{60}),e(\tfrac{17}{60})\bigr).
\]
These are the Fibonacci modular data, the data of the level-one $G_2$ WZW model at $c=14/5$ \cite{RSW}. The Lee--Yang data of $M(2,5)$ are not a twist of them by any signed permutation and scalar, since the absolute values of their diagonal $S$-entries are $(2/\sqrt5)\sin(2\pi/5)$, not $(2/\sqrt5)\sin(\pi/5)$.
\end{proposition}
\status{Proved here by the method of Theorem~\ref{thm:m27}, with the exact verification in $\Q(\zeta_{120})$; new in version~4. Scripts: \note{exact\_shadow\_reps.py} and \note{m5\_vs\_M25.py}.}

The rank-two note matches the Lee--Yang exponents $-\tfrac1{60}$ and $\tfrac{11}{60}$ through the cubic phase transport. By Proposition~\ref{prop:fibonacci}, the order-$5$ shadow carries the Fibonacci representation rather than the Lee--Yang one. Rank-two modular data are classified in \cite{RSW}, where Lee--Yang is related to Fibonacci by Galois conjugation; that relation was not verified separately here.

\paragraph{What these results add.}
Theorem~\ref{thm:m27} places Ramanujan's order-$7$ mock theta functions, and with them the $\widehat Z$ invariant of $-\Sigma(2,3,7)$, inside the modular representation of the $(2,7)$ minimal model: after conjugation and a twist by the character of $\eta^3$, their shadow transforms exactly as the $M(2,7)$ characters. The classical relation of this kind for order $5$ is the mock theta conjectures, which express the order-$5$ functions through the Rogers--Ramanujan functions, the characters of $M(2,5)$ \cite{AndrewsGarvan,Hickerson}. Proposition~\ref{prop:fibonacci} shows that for order $5$ the shadow's own modular data are the Fibonacci data, a Galois companion of Lee--Yang.

My assessment: the two results together suggest a pattern for the Brieskorn spheres $\Sigma(2,3,q)$, namely that the shadow representation of the associated mock vector is a twisted minimal-model representation or one of its Galois conjugates. I think so because the two cases $q=5,7$ fall under this description, and the pattern can be tested on further $q$ with the same scripts (Question~\ref{q:brieskorn}). For the Erdős--Straus side, the notes' map from the cubic middle to the three support labels is the link that would carry this structure over. That map is stated in the note's theorem \emph{Coordinate realization of middle support} and is not verified here.

\subsection{Further structural facts}\label{ssec:structfacts}

Three further structural statements of the notes were checked by program in reading passes, and one of them again here.
\begin{itemize}
\item \emph{The Niemeier lattice.} The Niemeier lattice with root system $A_5^4D_4$ and glue of order $72$ appears in the notes and as archive Theorem~25.8 \cite{ConwaySloane}. Its umbral group, the automorphism group of the glue code, is $GL_2(3)$ \cite{CDH}; as noted in Section~\ref{ssec:noteserrors}, $GL_2(3)$ and the binary octahedral group are distinguished by their involutions.
\item \emph{$\Gamma_0(6)$ and the class $6B$.} $T_6=\eta(\tau)^5\eta(3\tau)/(\eta(2\tau)\eta(6\tau)^5)$ is a Hauptmodul of $\Gamma_0(6)$, and $T_6+5+72/T_6$ is the McKay--Thompson series of the Monster class $6B$ \cite{ConwayNorton}.
\item \emph{Cayley--Dickson algebras.} The statements on sedenion zero divisors were checked in the conventions stated in the notes \cite{Moreno,BaezOctonions}.
\end{itemize}
The Busy-Beaver notes propose a pattern in residues modulo $107$, the residual of $p_*$. A test in an earlier reading pass compared these residues with a null model fixed in advance and found their frequencies within it (held-out $p=0.38$); the test was not re-run here, and it concerns that one proposed pattern only. The Lorentzian notes are treated in Section~\ref{ssec:lorentz}, where the relation between the modular flow and the Lorentz boost of their section on finite-dimensional modular flow is made precise (Proposition~\ref{prop:flowboost}).

\section{Negative results, with exact scope}\label{sec:negative}

Each row states a proposed implication that fails, with its counterexample or derivation, and the scope of the failure: what the counterexample does not show. Every certificate prime in the table is itself solvable. So these results limit specific mechanisms, not the conjecture, and several of them name the further input a mechanism needs.

\begin{center}\small
\begin{longtable}{@{}>{\raggedright\arraybackslash}p{0.19\linewidth}>{\raggedright\arraybackslash}p{0.29\linewidth}>{\raggedright\arraybackslash}p{0.21\linewidth}>{\raggedright\arraybackslash}p{0.21\linewidth}@{}}
\toprule
Source & What fails & Certificate & What it does not show\\
\midrule
\endhead
bounded transport & every occupied shell has a solution with unary layer $A\le2$ & $p=37$, $a=12$: only $u=8$ & the global cutoff ($a=10$ works at $37$)\\
first-two-shell sieve & occupancy of the $R=7$ shell gives a first-two-layer hit & $p=373$, $a=95$ & anything about other shells\\
item 2 & every canonical closed component contains a local hit & $p=2521$: a 7-cycle sink & $2521$ is solvable\\
item 3 & strict canonical descent ends at a hit & a closed triangle at a 22-digit prime & that prime is solvable\\
item 6 & the principal-only bound, summed over shells, is positive for large $p$ & $\sum L_a\le-\frac1{24}p\log p$ & anything about the true counts\\
item 6 & a nonnegative reweighting of it certifies occupancy & $p=87481$: all $21{,}870$ shells have $L_a\le0$ & only this certificate family fails\\
item 14 & same-grade middle repair from exterior states with $\alpha=-4t$, $t\ge4$ & $p=7840561$ & those primes are solvable\\
item 21 & a same-word residual return for words $u\nmid36$ & $3{,}041$ of the $3{,}654$ sources with $u\nmid36$ at primes below $2600$ & endpoint solutions exist\\
item 26 & a trace-to-solution map never increases $a$ & $p=67369$ & $67369$ is solvable\\
record shells (archive §19.3) & the deficit-box conclusions without strict recordness & $p=806{,}521$: full difference set, eight points missed & anything about strict records\\
author's notes (Section~\ref{sec:notes}) & the one-congruence shell criterion without coprimality & $p=73$, $R=7\cdot73^2$ & anything about shells with $p^2\nmid R$\\
author's notes & every prime $p\equiv1\pmod4$ has a lock certificate & $p=409$ & anything about the hard classes\\
Jacobi barrier (archive §10) & its converse: a non-residue prime factor of $a$ makes the shell occupied & $p=5$, $R=7$, $a=3$ & the barrier itself (Proposition~\ref{prop:jacobi})\\
no-two-shears theorem (supplement §5) & the theorem at the primes $p\equiv2\pmod3$ with $R(R+4)\mid p+4$ for a suitable residual $R$, for instance at every $p\equiv131\pmod{180}$; so no congruence hypothesis weaker than $p\not\equiv2\pmod3$ suffices (Propositions~\ref{prop:twoshears} and~\ref{prop:sheargraph}) & two-edge paths at $p=131$ and $p=1613$ & anything at $p\not\equiv2\pmod3$, where the theorem holds; it holds at most primes $p\equiv2\pmod3$ (Remark~\ref{rem:sheardensity}); anything about solvability\\
carrier census (archive §24) & forcing every row of the census domain by fixed-divisor carriers & the $2{,}759{,}345{,}613{,}832{,}500$ square rows (Corollary~\ref{cor:censuslimit}) & anything about the primes in those rows; carriers that use the factorisation of $(p+R)/4$\\
\bottomrule
\end{longtable}
\end{center}
All of these are statements about methods or certificate families, and none is a counterexample to the conjecture. Where a mechanism is sharp, the table says so: $p\not\equiv2\pmod3$ is the weakest congruence hypothesis under which the no-two-shears theorem holds at every prime, and it holds at most primes $p\equiv2\pmod3$ as well.

The archive records further negative results of the same kind. They are \textsf{located}, not checked here:
\begin{itemize}
\item the withdrawn constructions of §5 (p.~22), each with its point of failure;
\item the retracted converse of the Star--Kneser forcing theorem (§15.3, p.~122);
\item the failed universal bridge of Dyachenko's bounded identities (§12.2, p.~103);
\item a finite counterexample to one projection step in López Bonilla's 2026 preprint (§19, provenance, p.~148): in $C_6$, the set $\{0,1\}$ has trivial stabiliser, but its reduction modulo $2$ is all of $C_2$. The archive restricts this to that step and does not assess the whole paper.
\end{itemize}

\section{Bridges to other areas, and the directions they open}\label{sec:overlaps}

The research behind this paper explored many directions quickly, and in the author's view the bridges it built are among its most valuable outcomes. This section describes them. It does not claim that any of them is a complete research programme, and it does not treat a direction that is still incomplete as invalid: much of mathematics consists of programmes that are incomplete and are nevertheless valid mathematics. The audit did not have the compute to examine every direction to the depth needed to decide how promising it is. For each direction the section states what is established, with references to the proved results; what is not verified here; and my assessment.

The assessment addresses two separate questions. The first is the direction's interest on its own merit, for example as a bridge between two areas. The second is its bearing on the Erdős--Straus problem. When a direction relates the equation to a structure that does not take primes as input, the second question is what typed map connects the two, and how far that map is established. The assessments are mine and are marked as such; they are views held at this stage and may change as the directions are explored further.

\subsection{The split-zero support language}\label{ssec:splitzero}

The notes describe shells, classes and certificates through the split-zero semiring $G(R)=R\sqcup\{\tau\}$. In it, $\tau$ is an additive identity distinct from the ring zero $0_R$, and $\tau$ absorbs multiplication. The same semiring is the coefficient algebra of the zeta programme \cite{ZetaReader}, and it is archive §21. The note \emph{A split-zero repair dossier} develops the algebra of $G(R)$: ideals, prime ideals, spectrum, congruences, localisation and semimodules. Its classification statements were checked in a reading pass. One theorem of the dossier needs an additional relation, and the correction keeps to the construction's own convention of keeping $\tau$ and $0_R$ apart.

\begin{proposition}[the dossier's Theorem~7.2, with its relation]\label{prop:dossier}
Let $S=G(R)$, let $M_S$ be its multiplicative monoid, and let $\mathcal R$ be the congruence on the free commutative semiring $\NN[M_S]$ generated by $[a]+[b]\sim[a+b]$ for $a,b\in R$ and $[x]+[\tau]\sim[x]$ for $x\in S$.
\begin{enumerate}[label=(\alph*)]
\item The class of the empty sum $0$ in $\NN[M_S]/\mathcal R$ is $\{0\}$.
\item The natural map $\bar\Theta:\NN[M_S]/\mathcal R\to S$ sends both $0$ and $[\tau]$ to $\tau$, so it is not injective.
\item After adjoining the relation $[\tau]\sim0$, $\bar\Theta$ is an isomorphism.
\end{enumerate}
\end{proposition}
\status{Proved here; found by a reading pass of version~3, new as a proposition in version~4.}
\begin{proof}
(a) Every generating relation relates two non-empty formal sums. Adding a sum, or multiplying by a non-zero sum, preserves non-emptiness, and multiplying by $0$ gives $0\sim0$. So no chain of relations connects $0$ to a non-empty sum.
(b) A semiring homomorphism sends $0$ to $0_S=\tau$, and $\bar\Theta([\tau])=\tau$.
(c) With $[\tau]\sim0$, every formal sum reduces to $[\tau]$ or to $[r]$ for some $r\in R$: the terms $[\tau]$ can be removed by $[x]+[\tau]\sim[x]$, the empty sum is $\sim[\tau]$, and the remaining terms $[r_1],\dots,[r_k]$ combine to $[r_1+\dots+r_k]$. These representatives are pairwise inequivalent, because $\bar\Theta$ sends them to the distinct elements $\tau$ and $r\in R$ of $S$ (recall $\tau\ne0_R$). So $\bar\Theta$ is bijective, and it is a homomorphism by construction.
\end{proof}

The support states also have a typed meaning in the shell arithmetic. For a shell $a$, the three states $\tau$, $0$ and $+$ correspond to three situations:
\begin{itemize}
\item neither $-1$ nor $-4$ lies in the group $H_a$, and the barriers of Section~\ref{ssec:jacobi} apply;
\item $-1$ or $-4$ lies in $H_a$, but no divisor of $a^2$ reaches a target;
\item the shell is occupied.
\end{itemize}
By Proposition~\ref{prop:groupconverse}, under the exponent condition the middle state does not occur. Applied to classes of primes, as in Remark~\ref{rem:termcore}, the same vocabulary calls a class supported when some certificate applies on a lift of it.

My assessment: I think this vocabulary is useful for the problem, because it separates the two reasons for which a shell can be empty. One is group-theoretic, and the barriers handle it; the other is a matter of size, in which the exponents of $a$ are too small to reach a target inside the group. They call for different tools. The algebra of $G(R)$ developed in the dossier is also an object of independent interest.

\subsection{Mock and quantum modularity}

What is established is in Section~\ref{ssec:mock}: the notes' explicit mixed mock functions are Ramanujan's order-$7$ mock theta functions, and the first is $q^{1/2}\widehat Z_0(-\Sigma(2,3,7))$; their shadow carries the $M(2,7)$ representation after conjugation and the $\eta^3$ twist (Theorem~\ref{thm:m27}); and the order-$5$ block carries the Fibonacci data (Proposition~\ref{prop:fibonacci}). Not verified here: the note's support isomorphism from the cubic middle of the shell problem to the three support labels, its mass-gap coordinate for the cubic middle, and its Virasoro support-gap corollary.

My assessment: on its own merit, I think this is the most developed of the structural directions. It has produced a result, Theorem~\ref{thm:m27}, that stands independently of the Erdős--Straus problem and connects Ramanujan's order-$7$ functions, the $3$-manifold invariant of $-\Sigma(2,3,7)$ and the $(2,7)$ minimal model. For its bearing on the equation, the typed map needed runs from the shell data, which depend on $p$ through the factorisation of $(p+R)/4$, to the modular object. The notes' map goes through the finite support labels. An analytic object in the notes whose coefficients depend on that factorisation is the pair-count series of Proposition~\ref{prop:pairzeta}. A typed map between the pair-count series and the mock vector would carry information in both directions, and constructing one is Question~\ref{q:modularobject}.

\subsection{Moonshine, Niemeier lattices and Ogg's primes}

Checked in reading passes (Section~\ref{ssec:structfacts}): the Niemeier lattice with root system $A_5^4D_4$, its umbral group $GL_2(3)$, the $\Gamma_0(6)$ Hauptmodul and the series of the class $6B$. A constructive-algebra preprint of the notes also connects the extension prime $29$ with the Heegner discriminant $163$. The factorisation it uses is classical and correct: $j\bigl((1+\sqrt{-163})/2\bigr)=-640320^3$ with $640320=2^6\cdot3\cdot5\cdot23\cdot29$; by contrast, $5280=2^5\cdot3\cdot5\cdot11$ for the discriminant $67$. Not verified here: the preprint's rank-$15$ Ogg module, its genus-zero and supersingular realisations, and its phase decomposition, apart from the map treated next.

\begin{proposition}[the torsor map]\label{prop:torsor}
The map $\Theta:(\ZZ/18)^\times\to S_{840}$ of \emph{Terminal residual coordinates} §10, given by $1\mapsto1$, $5\mapsto289$, $13\mapsto361$, $17\mapsto169$, $7\mapsto121$, $11\mapsto529$, is a bijection but not a homomorphism. The map $5^k\mapsto289^k$ is an isomorphism, and it differs from $\Theta$ exactly by exchanging the images of $7$ and $13$.
\end{proposition}
\status{Proved here; new as a proposition in version~4.}
\begin{proof}
Both groups are cyclic of order $6$: $5$ generates $(\ZZ/18)^\times$, with powers $5,7,17,13,11,1$, and $289$ generates $S_{840}$, with powers $289,361,169,121,529,1$ modulo $840$. So $5^k\mapsto289^k$ sends $5,7,17,13,11,1$ to $289,361,169,121,529,1$. It agrees with $\Theta$ except at $7$ and $13$, whose images are exchanged. In particular $\Theta(5)^2=289^2\equiv361\pmod{840}$, while $5^2\equiv7\pmod{18}$ and $\Theta(7)=121$.
\end{proof}

The note's odd Ogg phase decomposition (its Theorem~10.1) uses only that $\Theta$ is a bijection, so it holds as stated; with the isomorphism in place of $\Theta$, two of its sectors exchange their contents (a reading pass).

My assessment: the moonshine data organise the residue layer of the problem, meaning the six hard classes and the core, in a way that matches known structures: the umbral group of $A_5^4D_4$, the $\Gamma_0(6)$ Hauptmodul and Ogg's primes. I think this match is of independent interest as an observation. Its bearing on the hard classes depends on the typed map that is used. What is proved is that polynomial identities and fixed-divisor certificates do not reach the square classes (Section~\ref{sec:problem}, Proposition~\ref{prop:nosquare}). I have not seen a map from the shell data to the moonshine data in the notes, and I do not have a view yet on whether one exists.

\subsection{Lorentzian and Apollonian geometry}\label{ssec:lorentz}

The Descartes form $F(x)=2\sum x_i^2-(\sum x_i)^2$ has signature $(3,1)$. The Apollonian group is a subgroup of $O_F(\ZZ)$ of infinite index that is Zariski dense in $O_F$, a thin group, and it acts on hyperbolic $3$-space \cite{SarnakLetter,GLMWY}. The notes identify a Lorentz form of their shell geometry with the Descartes form. The same Lorentzian geometry of circle configurations underlies Kocik's circle-geometric reading of the Koide relation \cite{Kocik}.

The note \emph{Lorentzian fixed fibers, Descartes sheets, and origin-resolved positivity in residual Erdős--Straus shells}, in its section \emph{Finite-dimensional modular flow as a Lorentz boost}, states that the congruence action $X\mapsto A_\rho XA_\rho$, with $A_\rho=\operatorname{diag}(e^{\rho/2},e^{-\rho/2})$, is a Lorentz boost on the Hermitian $2\times2$ matrices with the form $\det$. This is correct. The relation between this boost and the modular flow of the section's title is the following.

\begin{proposition}[modular flow and boost]\label{prop:flowboost}
Let $H_r=\operatorname{diag}(e^r,e^{-r})$ and $X=\begin{psmallmatrix}h+w&u+iv\\u-iv&h-w\end{psmallmatrix}$, with the Minkowski form $\det X=h^2-u^2-v^2-w^2$. For $z\in\C$ the map $X\mapsto H_r^zXH_r^{\bar z}$ lies in $SO^+(1,3)$.
\begin{enumerate}[label=(\alph*)]
\item At $z=it$ it is the modular automorphism $\sigma_t=\operatorname{Ad}(H_r^{it})$ of the state $\varphi_{H_r}$ with density proportional to $H_r$. It acts as the rotation by the angle $2rt$ in the $(u,v)$-plane and fixes $h$ and $w$.
\item At real $z=s$ it is the boost of rapidity $2rs$ in the $(h,w)$-plane and fixes $u$ and $v$. The note's $A_\rho$ is the case $s=\rho/(2r)$.
\end{enumerate}
So the modular flow and the note's boost are the compact and the non-compact real forms of one complex one-parameter subgroup $z\mapsto H_r^z$ of $SL_2(\C)$, acting on Minkowski space by $X\mapsto gXg^*$.
\end{proposition}
\status{Proved here; new in version~4. The coordinate formulas are checked symbolically in \note{misc\_checks.py}.}
\begin{proof}
$\det H_r^z=1$ and $(H_r^z)^*=H_r^{\bar z}$, so the map is the standard action of $SL_2(\C)$ on Hermitian matrices, which lies in $SO^+(1,3)$ because $SL_2(\C)$ is connected. At $z=it$, $H_r^{\bar z}=H_r^{-it}$, so the map is $\operatorname{Ad}(H_r^{it})$, the modular group of $\varphi_{H_r}$; it multiplies the off-diagonal entry $u+iv$ by $e^{2irt}$ and fixes the diagonal, which gives $(h,\,u\cos2rt-v\sin2rt,\,u\sin2rt+v\cos2rt,\,w)$. At $z=s$ it multiplies the diagonal entries $h+w$ and $h-w$ by $e^{2rs}$ and $e^{-2rs}$ and fixes $u+iv$, which gives $(h\cosh2rs+w\sinh2rs,\,u,\,v,\,w\cosh2rs+h\sinh2rs)$.
\end{proof}

My assessment: on its own merit, I think the Descartes setting is a natural place to ask how such Lorentz one-parameter families interact with the Apollonian thin group, and the question is one about thin groups in its own right. A companion note answers it \cite{ApollonianNote}. The orientation-preserving subgroup $A^+$ of the Apollonian group meets every compact subgroup of $SO_F(\R)$ only in the identity, so under any identification of the Descartes space with the Hermitian matrices no nontrivial element of the Apollonian group is a modular automorphism of a faithful state (the modular automorphisms lie in compact subgroups of $SO^+(1,3)$, and the reflections $S_i$ are not in $SO_F$). Boosts occur in the stabilisers of the circles of a packing, which are congruence groups (in the strip packing, the image of $\Gamma^0(4)$ in $PSL_2(\ZZ)$), and also as boosts that fix no circle, such as $S_1S_2S_1S_3S_4S_3$; every other loxodromic element of $A^+$ is a screw motion whose rotation angle is an irrational multiple of $2\pi$. For the bearing on the equation, the notes' map from shell data to Descartes quadruples is not verified here. I have not identified a typed map from shell occupancy to the Apollonian group, and I have no view yet on whether one exists.

\subsection{Additive combinatorics: Kneser's theorem}

The notes apply Kneser's theorem on sumsets in abelian groups \cite{Kneser} to the residues of divisors, producing Star--Kneser forcing, which the archive states as a per-shell sufficient condition for occupancy (Theorems~15.6 and~20.7; located, not verified here). Its converse fails at $p_*$, $R=107$, as the archive records (Proposition~15.7).

My assessment: I think additive combinatorics is a natural toolset for the second reason for which a shell can be empty (Section~\ref{ssec:splitzero}), where the target lies in $H_a$ but may not be reached. In additive notation the residues of the divisors of $a^2$ form a sumset of arithmetic progressions, one for each prime $\ell\mid a$ with $2v_\ell(a)+1$ terms, so sumset growth is exactly what decides whether they fill $H_a$.

\subsection{Computability: the Busy-Beaver notes}

The statement that $4/n=1/x+1/y+1/z$ is solvable for every $n\ge2$ is $\Pi_1$. Each instance is decided by a bounded search: in a solution with $x\le y\le z$ one has $x\le3n/4$, then $y\le2/(4/n-1/x)$, and $z$ is determined. So there is a Turing machine that halts if and only if the conjecture fails, and if it has $k$ states, the value of the Busy-Beaver function at $k$ decides the conjecture. This is the standard bridge between $\Pi_1$ statements and the Busy-Beaver function; such machines have been constructed explicitly for Goldbach's conjecture and for the Riemann hypothesis \cite{Aaronson}. The archive uses Busy-Beaver values, quoted from the literature and not checked here ($S(4)=107$, $\mathrm{Space}(5)=12289$), to select the shells of its §§3 and 20, and its §25 maps Busy-Beaver extrema to shell carriers (located). The proposed pattern in residues modulo $107$ is discussed in Section~\ref{ssec:structfacts}.

The $\Pi_1$ bridge is exact, and the least number of states of a machine that halts exactly when the conjecture fails is a well-posed target (Question~\ref{q:pi1}).

\subsection{Further directions in the same collections}

The collections also contain notes on Cayley--Dickson algebras and sedenion zero divisors, on the Koide relation, on a mass gap and on anomalies, and a split-zero zeta note. The Cayley--Dickson statements were checked in a reading pass (Section~\ref{ssec:structfacts}). The Koide, mass-gap and anomaly notes lie outside the scope of this paper; they were not read for it and are not verified here.

\subsection{The author's other workbenches}\label{ssec:wbedges}

Material crosses between this workbench and the author's other workbenches. The audit note \note{47\_} records every crossing that the audit found as a typed edge, of one of four types:
\begin{itemize}
\item \emph{P}: a proof-level import;
\item \emph{D}: a dependency or citation;
\item \emph{I}: the same identity or object on both sides;
\item \emph{A}: an analogy, with no typed map yet.
\end{itemize}
The table lists the edges that touch this workbench, and adds three that this paper found in the archive (marked \emph{new}).

\begin{center}\small
\begin{longtable}{@{}>{\raggedright\arraybackslash}p{0.13\linewidth}>{\raggedright\arraybackslash}p{0.44\linewidth}>{\raggedright\arraybackslash}p{0.07\linewidth}>{\raggedright\arraybackslash}p{0.28\linewidth}@{}}
\toprule
Edge & What crosses & Type & Status\\
\midrule
\endhead
NS $\to$ ES & The auxiliary torus endomorphism $J=\begin{pmatrix}3&1\\1&5\end{pmatrix}$ (determinant $14$) of the Navier--Stokes manuscript. It enters 21 statements of the supplement (§8) as an operator only. The core identity is $\{(u,v)\in G^2:u^3v=1=uv^5\}=\{(u,u^{-3}):u^{14}=1\}$ & P & audited (\note{21\_} §1); no fluid theorem is used\\
S$^6\to$ ES & The integral monodromies $A_1,A_2,A_\infty$ of the $(3,4,\infty)$ period system, in ES-09 & P & audited for odd levels; the orbits at every level are Theorem~\ref{thm:orbits} (§\ref{ssec:es-s6})\\
ES $\leftrightarrow$ S$^6$ & The Niemeier lattice with root system $A_5^4D_4$ and glue index $72$ (ES result R10) & I & ES's glue code verified; the two codes not compared\\
Jacobian map & Alpöge's map $F:\C^3\to\C^3$ with $\det DF=-2$ and the three-point fibre over $(-\tfrac14,0,0)$ (archive §§4, 7). The same map is in the Yang--Mills workbench and the zeta programme. ES's four-dimensional extension $P$ ($\det DP=-2$) crosses to the zeta programme's conductor & P & $\det DF$, $\det DP$ and the fibre verified (zeta paper \cite{ZetaReader}, Lemma~6.6; \note{21\_}); the conductor maps located\\
ES $\leftrightarrow$ zeta & The split-zero semiring $G(R)=R\sqcup\{\tau\}$ with its Boolean support character (archive §21; ES R01) & I & audited (\note{21\_} §6)\\
ES $\leftrightarrow$ zeta & A finite Weyl phase-space basis & I & ES side verified on $C_2\times C_3$\\
zeta $\to$ ES & NG20--NG21, used by item 20; the SplitZero packages of 13--16 September & D & located\\
ES $\to$ zeta & The 488-page reader's source, vendored into the zeta repository & D & located\\
ES $\to$ zeta & The 124-page ES/Fable reader: solutions as binary quartics, then the Fable cover and a weighted conductor (§\ref{ssec:crosswalk}) & P & the quartic relation proved here (Proposition~\ref{prop:quartic}); the rest located\\
ES $\leftrightarrow$ Collatz & $3(4n+1)+1=4(3n+1)$, i.e.\ $\theta T=4\theta$ (ES R05, diagrams by u/CivQ17) & I & audited\\
Collatz $\leftrightarrow$ ES $\leftrightarrow$ zeta & The homogenisation $(x,h)\mapsto(ax+bh,h)$ (ES X4) & I & ES's diagram read\\
Busy Beaver (new) & Busy-Beaver extrema quoted by the archive ($S(4)=107$ and $\mathrm{Space}(5)=12289$; for the five-state values it cites the formal enumeration of the bbchallenge collaboration) select the shells of archive §§3 and 20; §25 maps extrema to shell carriers. The zeta programme's detectability argument \cite{ZetaReader} uses a different machine for a different purpose & A & located; the two use different machines for different purposes\\
Connes (new) & Supplement §6 reads Connes' primitive intersections; archive §24 opens with Tomita--Takesaki transport. The zeta programme uses Connes' summation map and Meyer's quotient & A & located; no common statement was found\\
Fable (new) & \emph{Fable} names three different objects across the repositories: the three-dimensional Keller map, the $S^6$ construction in older files, and a conductor cover (\note{47a\_} §10). Archive §5 withdraws one Fable construction & --- & disambiguate at each use\\
ES $\leftrightarrow$ YM & A shared term, \emph{gap} (ES X7), and a proposed success projector & A & located; no typed map is written\\
\bottomrule
\end{longtable}
\end{center}

The crossings with typed content are the torus operator, the monodromy matrices, the Keller maps, the split-zero semiring and two identities; each is audited or recorded as located. The other connections are proposals awaiting typed maps, and the archive describes its late Fable, operator, Connes and moonshine material in the same way (p.~4): as a candidate interface, pending an explicit typed map and compatibility check.

\section{Questions, and what was not checked}\label{sec:open}

\subsection{Questions}

\begin{enumerate}[label=\arabic*.,ref=\arabic*]
\item\label{q:global} \emph{The global statement.} Show that for every prime $p\equiv1,121,169,289,361,529\pmod{840}$ some shell is occupied. This is the conjecture itself, restated by Theorem~\ref{thm:shell}, and the archive names it as the remaining obligation (p.~4). By Proposition~\ref{prop:groupbarrier}, a proof must produce, for some $R$, a shell whose group $H_a$ contains $-1$ or $-4$; when $R$ is prime this means a prime factor of $(p+R)/4$ that is a non-residue modulo $R$. That condition alone does not suffice, and Proposition~\ref{prop:groupconverse} says when it does.
\item\label{q:typeII} \emph{Type II universality} (Section~\ref{ssec:locks}). Does every prime $p\equiv1\pmod{24}$ have a solution with $p$ dividing exactly two denominators? A positive answer proves the conjecture. Below $10^7$ every such prime has one in a shell $R\le107$, and only $p_*=8{,}803{,}369$ needs $107$; a reading pass found the same below $10^8$.
\item\label{q:qrsurvivors} \emph{The quadratic-residue survivors} (Table~\ref{tab:survivors}). Are there infinitely many primes that satisfy every condition of Section~\ref{ssec:qrconditions}? They have density zero (Proposition~\ref{prop:densityzero}), and $23$ of them lie below $10^7$. A negative answer would prove the conjecture at all but finitely many primes.
\item\label{q:550} \emph{The $550$ supported classes} (Theorem~\ref{thm:termcore}). Is each of the $550$ non-square classes left open modulo $840\cdot11\cdot19\cdot23$ covered by finitely many single identities at some finite level? If one is not, no identity sieve reduces the problem to the square core.
\item\label{q:census} \emph{The census survivors.} $2.394\%$ of the final hard domain survives the carrier chain (§\ref{ssec:census}), and the square rows among them can never be forced by its carriers (Corollary~\ref{cor:censuslimit}). Can the non-square survivor rows all be forced by carriers of the same kind, if more moduli are added? Which carriers, using the factorisation of $(p+R)/4$ and not only residues of $p$, force the square rows?
\item\label{q:deficit} \emph{The strict-record deficit-box conjecture} (§\ref{ssec:records}; archive §19.2). It holds at the two strict records with non-empty deficit below $8{,}803{,}369$, and strict recordness cannot be dropped from it ($p=806{,}521$). A further test needs a strict record above $8{,}803{,}369$, and there is none below $3\cdot10^9$ (Proposition~\ref{prop:records}).
\item\label{q:growth} \emph{Growth of the least class.} Is $h$ unbounded on the primes? Along the strict records up to $3\cdot10^9$ it takes the values $1,2,3,6,8,15,27$; below $3\cdot10^9$ the next largest values are $21$, $20$ and $19$, at nine primes; and every record after $73$ lies in a hard class (Corollary~\ref{cor:hle2}). The workbench states no conjecture on the growth of $h(p)$, and none is stated here.
\item\label{q:shear} \emph{The density of the shear paths.} What is the exact relative density, among the primes $p\equiv2\pmod3$, of the primes with a path of two shear edges? It is at least $1/24$, from the class $131\bmod180$, and at most $0.335$ (Remark~\ref{rem:sheardensity}); below $10^7$ the proportion is $5.24\%$.
\item\label{q:genera} \emph{Genera of the monodromy covers.} What are the genera of the components of the $(3,4,\infty)$ covers at every level? Theorem~\ref{thm:orbits} gives the orbits, and the genera follow from the cycle counts of $A_1$, $A_2$ and $A_\infty$ on each orbit (§\ref{ssec:es-s6}).
\item\label{q:brieskorn} \emph{Brieskorn spheres $\Sigma(2,3,q)$.} For which $q$ does the shadow representation of the associated mock vector coincide, after conjugation and a twist, with a minimal-model representation, as for $q=7$ (Theorem~\ref{thm:m27}), or with a Galois conjugate of one, as for $q=5$ (Proposition~\ref{prop:fibonacci})?
\item\label{q:modularobject} \emph{A modular object for the shell counts.} Is there a modular, mock or quantum modular object whose coefficients are the shell counts $2|E_a|+|M_a|$, or a typed map from the pair-count series of Proposition~\ref{prop:pairzeta} to the mock vector of Section~\ref{ssec:mock}?
\item\label{q:thin} \emph{A typed map to the Descartes geometry.} Is there a typed map from shell occupancy, or from the notes' Descartes quadruples, to the Apollonian group? The companion note \cite{ApollonianNote} determines how the rotations and boosts of Proposition~\ref{prop:flowboost} meet that group; the map from the shell data is what is missing.
\item\label{q:pi1} \emph{The least $\Pi_1$ machine.} What is the least number of states of a Turing machine that halts if and only if the conjecture fails?
\end{enumerate}

\subsection{What this paper did not check}

\begin{itemize}
\item \emph{The archive.} Apart from the parts marked proved or reproduced in §\ref{ssec:archive}, its 619 pages were mapped, not audited. In particular, the structural constructions (§§2, 4, 6--8, 9 after its opening, 21--23, the operator part of §24, the structural material of §24.4, and §25) were not audited, and the literature audits (§§11--14 and 17) were not compared with the papers they audit.
\item \emph{Lean.} The source archive contains 67 files with the extension \texttt{.lean} (seven of them \texttt{.axioms.lean}); they were not rebuilt. The archive's own account of what each Lean module checks is quoted where it matters (§\ref{ssec:census}).
\item \emph{Section~\ref{sec:notes}.} Several statements there rest on the reading passes and are marked so:
\begin{itemize}
\item the lattice and modular facts of Section~\ref{ssec:structfacts}, apart from the involution count;
\item the comparison of the notes' $q$-expansions with Zagier's vector $M_7$;
\item the Busy-Beaver residue test, and the Busy-Beaver values that the archive quotes from the literature;
\item the Cayley--Dickson identities;
\item the counts below $10^8$.
\end{itemize}
Salez's completeness statement was not proved here. Drafts that the author's latest record supersedes were not treated.
\item \emph{The census.} The enumerations of §\ref{ssec:census} were not re-run. Only the arithmetic consistency of the reported counts was checked, together with the base rows and the carrier sets used in Corollary~\ref{cor:censuslimit}.
\item \emph{The supplement.} The general proofs of the supplement's Sections 1--7 were not re-derived, except that of the no-two-shears theorem of Section 5 (Proposition~\ref{prop:twoshears}); their finite content was reproduced in part (§\ref{ssec:supplement}). Section 8, the auxiliary torus, was not checked.
\item \emph{The 124-page reader.} Only Proposition~\ref{prop:quartic} was checked.
\item \emph{The continuation.} Items 1, 5, 9--13, 15--20, 24, 25, 27 and 28 were located, not checked. So were item 4's enumeration, the deduction in item 23, and ES-04 to ES-08 beyond the ES-08 witness.
\item \emph{Novelty.} No literature search for novelty was made here, except where a status line says so. The workbench claims no novelty for the classical coordinates and identities it uses. Several results were found by referee passes (Propositions~\ref{prop:sheargraph}, \ref{prop:groupbarrier}, \ref{prop:groupconverse}, \ref{prop:fixeddivisor} and~\ref{prop:nosquare}, and Corollaries~\ref{cor:hle2} and~\ref{cor:censuslimit}); some of them are forms of statements of the archive or the notes, as their status lines say. The results marked new in version~4 are not in the workbench, the notes or the earlier versions of this paper, as far as these were read; whether they are in the literature was not checked, and none is claimed as new mathematics beyond that.
\end{itemize}

\appendix
\section{Catalogue of statements}\label{app:catalogue}

This catalogue lists every statement of the workbench that this paper states, proves, reproduces or locates. For each it gives where the statement is in the workbench and how far it was checked here, in the labels defined in the introduction. \emph{WB} quotes the workbench's own label where it adds something. Locators are as in the front matter: \emph{archive} is the 619-page PDF, \emph{supplement} the 86-page supplement, and \emph{item $n$} the $n$-th entry of the README's \emph{Start reading} list.

\begin{center}\small
\begin{longtable}{@{}>{\raggedright\arraybackslash}p{0.36\linewidth}>{\raggedright\arraybackslash}p{0.22\linewidth}>{\raggedright\arraybackslash}p{0.22\linewidth}>{\raggedright\arraybackslash}p{0.11\linewidth}@{}}
\toprule
Statement & Where in the workbench & Status here & Here\\
\midrule
\endhead
\multicolumn{4}{@{}l}{\emph{The shell reduction}}\\*
The ordered pairs completing a first denominator $a$ number $2|E_a|+|M_a|$; solvable iff some shell is occupied (ES-01, ES-02) & archive Theorems~9.1 (p.~38) and~9.64 (p.~90), table (394); supplement §1 (for $p=12h+1$) & Proved here; the bijection and the count are the archive's, for every odd prime; this paper adds the swap analysis and (c) & Thm~\ref{thm:shell}\\
The least denominator lies in $(p/4,p/2)$, except for one solution at $p\equiv3\pmod4$ & a reduction step of item 22 & Proved here, for every odd prime & Thm~\ref{thm:shell}(a)\\
A counterexample prime has $E_a=M_a=\emptyset$ at every shell (the all-shell no-hit sieve) & supplement §1 & Audited (\note{43a\_} §1.2); it is Thm~\ref{thm:shell}(c) & §\ref{sec:core}\\
The residual-3 shell: criterion and count & supplement §1; archive Theorem~10.3 & Proved here; count extended to $p\equiv5\pmod{12}$ & Prop.~\ref{prop:r3}\\
The residual-7 shell: criterion & supplement §1; archive §10 & Proved here, for every $p\equiv1\pmod4$ & Prop.~\ref{prop:r7}\\
Both first shells empty only if $p$ lies in a hard class modulo $840$ & archive Corollary~10.7 (p.~95); the first pass observed the congruence modulo $24$ & Proved here & Prop.~\ref{prop:survivors}\\
Reduction to the six hard classes modulo $840$ & archive Theorem~10.6, p.~95 (with shells $R\le7$); classical (five identities) & Proved here & Props.~\ref{prop:survivors}, \ref{prop:mod840}\\
Jacobi-symbol barrier, and the failure of its converse & archive Theorem~10.9 and Proposition~10.10, pp.~96--97 & Proved here & Prop.~\ref{prop:jacobi}\\
The group barrier $-1\notin\langle4,\ell:\ell\mid a\rangle$, and a converse under an exponent condition & implied by archive §10 & Proved here (found by the referee pass) & Props.~\ref{prop:groupbarrier}, \ref{prop:groupconverse}\\
\midrule
\multicolumn{4}{@{}l}{\emph{Least classes, records and deficits}}\\*
$h(p)=(R+1)/4$ for the least occupied shell & archive §§17.4, 19 & Proved here & Lemma~\ref{lem:leastclass}\\
$h(p)\le2$ outside the six hard classes; every strict record after $73$ is hard & implied by archive Theorem~10.6 & Proved here (found by the referee pass); checked to $3\cdot10^9$ & Cor.~\ref{cor:hle2}\\
The seven strict records up to $8{,}803{,}369$; $1201\in\mathcal C_6\setminus\mathcal C_5$; $h(806521)=15$ & archive Proposition~17.6, pp.~131--132 (recovering Ionascu--Wilson's printed list and correcting one label); §19.2, p.~152 & Reproduced & Prop.~\ref{prop:records}\\
No strict record in $(8{,}803{,}369,\,3\cdot10^9]$; the primes with $h=19,20,21$ there & --- & Reproduced (two independent programs) & Prop.~\ref{prop:records}\\
At $p_*=8{,}803{,}369$ the least occupied shell has $R=107$, and it carries one unordered solution & archive §3 & Reproduced & §\ref{ssec:records}\\
The mod-$107$ deficit progression, the covering, and the two middle fibres (WB: credited to u/CommonCareful3149) & archive §16; bench R06 & Reproduced & §\ref{ssec:records}\\
The deficits at $118{,}801$ and $8{,}803{,}369$; the boundary example $806{,}521$ & archive §19.3 & Reproduced & §\ref{ssec:records}\\
The strict-record deficit-box conjecture & archive §19.2 (WB: conjecture) & Open & §\ref{ssec:records}, §\ref{sec:open}\\
The shell $a=3075$ at $p=12289$: counts $(3,6,3)$ & archive §20 & Counts reproduced (\note{esr\_referee\_\allowbreak checks.py}); the orbit and the quotient certificate located & §\ref{ssec:archive}\\
\midrule
\multicolumn{4}{@{}l}{\emph{Families and the census}}\\*
Fixed-divisor families: a divisor $u=sh$ or $u=s$ with $s\mid c^2$ solves the equation on a class modulo $4cR$ & archive Lemma~24.218, p.~422 (for primes) & Proved here, for all $n$ (found by the referee pass) & Prop.~\ref{prop:fixeddivisor}\\
Four families solving the equation on $248k+209$, $248k+89$, $17112k+11137$, $17112k+14113$ & archive Theorem~24.219, p.~423; Proposition~24.222 & Proved here; instances of Prop.~\ref{prop:fixeddivisor}; that they meet every hard class is shown here & Thm~\ref{thm:families}\\
The carrier census: $97.606\%$ of the final hard domain forced & archive pp.~3--5 and §24 (WB: bounded sufficient-carrier census) & Located; the arithmetic of the counts checked & §\ref{ssec:census}\\
No fixed-divisor carrier forces a square class; the census cannot force its $2{,}759{,}345{,}613{,}832{,}500$ square rows & --- & Proved here (found by the referee pass) & Prop.~\ref{prop:nosquare}, Cor.~\ref{cor:censuslimit}\\
The residual-eleven support ideal; the paired $R=7/R=11$ affine theorem & archive §18 & Located & §\ref{ssec:archive}\\
\midrule
\multicolumn{4}{@{}l}{\emph{The supplement and the 124-page reader}}\\*
The finite content of Sections 1--7 (gluing, ratio transport, shears, successors, orchard partition, boundary law, unary box, incidence, raw scale) & supplement §§1--7 & Reproduced in part & §\ref{ssec:supplement}\\
No two shear edges compose & supplement §5 (stated for $p\equiv1\pmod{12}$) & Generalised here to $p\not\equiv2\pmod3$; no weaker congruence hypothesis suffices; fails only on a minority of the $p\equiv2\pmod3$ & Props.~\ref{prop:twoshears}, \ref{prop:sheargraph}, Rem.~\ref{rem:sheardensity}\\
The Navier--Stokes auxiliary cover & supplement §8 & Not checked & §\ref{ssec:supplement}\\
The quartic of a solution: $p=-5u_4/u_3$ and the coefficient relation, and its converse & 124-page reader; item 11 & Proved here (also \note{21\_}); the converse found by the referee pass & Prop.~\ref{prop:quartic}\\
The Fable cover and the weighted conductor & 124-page reader & Located & §\ref{ssec:crosswalk}\\
\midrule
\multicolumn{4}{@{}l}{\emph{The September continuation}}\\*
Item 2: a seven-cycle sink at $p=2521$ (WB: method counterexample) & \texttt{research/incoming/} & The sink audited (certificate); the closure theorems located & §\ref{sec:items}, §\ref{sec:negative}\\
Item 3: a closed triangle at a 22-digit prime & \texttt{research/incoming/} & The triangle audited (certificate), its primality proved here; the escape theorem located & §\ref{sec:items}\\
Item 4: six nonresidue primes within 25 factorisations (WB: computer-assisted) & \texttt{research/incoming/} & Reductions read; conditional on an enumeration not re-run & §\ref{sec:items}\\
Item 6: the principal-only bound is negative in sum & \texttt{research/incoming/} & Audited & §\ref{sec:items}\\
Item 7: least-nonresidue localisation & \texttt{research/incoming/} & Audited & §\ref{sec:items}\\
Item 8: the middle identity and its bounds & \texttt{research/incoming/} & Reproduced on $1{,}613$ middle states; the atlas located & §\ref{sec:items}\\
Items 14 and 21: template bound; same-word returns & \texttt{research/incoming/} & Audited & §\ref{sec:items}\\
Item 22: the degree-8 invariant is negative on solutions & \texttt{research/incoming/} & Audited (certificate) & §\ref{sec:items}\\
Items 23 and 26 & \texttt{research/incoming/} & Audited in part & §\ref{sec:items}\\
Items 1, 5, 9--13, 15--20, 24, 25, 27, 28 & \texttt{research/incoming/} & Located (item 11's quartic proved here) & \note{43a\_} §2\\
ES-09: the $(3,4,\infty)$ monodromy covers & \texttt{received/\allowbreak es\_s6\_\allowbreak counterfactual/} & Audited for odd levels; at even levels the statements do not hold as stated; the orbits at every level proved here (new in version~4) & §\ref{ssec:es-s6}\\
ES-08: a forced solution at $p=1201$, $a=310$, $R=39$ & 12 September continuation & The witness checked; the forcing inequality not audited & \note{43a\_} §4\\
Results bench R01--R10, arguments X1--X7 & \texttt{RESULTS.md} & Checked in \note{43a\_} §3 (X1 not recomputed) & §\ref{sec:items}\\
\midrule
\multicolumn{4}{@{}l}{\emph{The author's notes (Section~\ref{sec:notes})}}\\*
The terminal core $S_{840}\times Q_{11}\times Q_{19}\times Q_{23}$ & source note of the core; the terminal-coordinates papers & Proved here (one direction; the other a finite computation, reproduced); the covering count $3520$ computed & Thm~\ref{thm:termcore}\\
The width-one barrier & pre-Niemeier note, Thm~3.1; archive Thm~15.5 & Proved here & Prop.~\ref{prop:groupbarrier}(c)\\
Eight quadratic-residue shifts; $p+5$; the visible-box shells; $\mu=17$ & implied by the notes' branches and locks & Proved here; $p+2$, $2p+1$, the extension of $p+5$ and five shells new & §\ref{ssec:qrconditions}\\
Non-coprime and ramified shells; the pair-count series & defect-completion article, Lemma~2.1 (which needs $p^2\nmid R$) and Thm~4.3 & Proved here & §\ref{ssec:singleshells}\\
The direct $29$-channels; Rosati's family; the full-divisor locks & terminal-coordinates paper, Thms~7.1--7.2; the locks note & Proved here; \emph{only two channels} holds only for single-prime divisors & Props.~\ref{prop:rosati}, \ref{prop:locks}\\
The survivors of the row $p+1$ have density zero & --- & Proved here (new in version~4) & Prop.~\ref{prop:densityzero}\\
The shadow of the order-$7$ mock theta vector and the characters of $M(2,7)$ & \emph{An explicit mixed-mock Virasoro support realization} & Proved here in the form of Thm~\ref{thm:m27} (new in version~4); the note's label matching does not realise it & Thm~\ref{thm:m27}\\
The order-$5$ block and the Fibonacci modular data & the rank-two note & Proved here (new in version~4) & Prop.~\ref{prop:fibonacci}\\
The split-zero dossier's Theorem~7.2 & \emph{A split-zero repair dossier} & Proved here with the relation $[\tau]\sim0$ & Prop.~\ref{prop:dossier}\\
The torsor map $(\ZZ/18)^\times\to S_{840}$ & terminal-coordinates paper, §10 & A bijection, not a homomorphism; the decomposition using it holds (a reading pass) & Prop.~\ref{prop:torsor}\\
Modular flow and Lorentz boost & the Lorentzian Descartes note & The note's boost proposition correct; the relation proved here (new in version~4) & Prop.~\ref{prop:flowboost}\\
\midrule
\multicolumn{4}{@{}l}{\emph{The rest of the archive}}\\*
Structural constructions: §§2, 4--8, 21--23, 25 and the operator part of §24 & archive & Located; §5 read; the fibre of §7 verified in the zeta paper & §\ref{ssec:archive}, §\ref{sec:overlaps}\\
Literature audits: §§11--14 and 17 & archive & Located; the classes of §17 replayed & §\ref{ssec:archive}\\
The Star--Kneser forcing theorem and its retracted converse & archive §15 & Located & §\ref{sec:negative}\\
\bottomrule
\end{longtable}
\end{center}

\section{How this paper was checked}\label{app:verification}

\paragraph{Sources.}
\begin{itemize}
\item The repository \texttt{KokunoYumeto/erdos-straus-foundation} at commit \texttt{d20b32e}, read-only.
\item The 619-page archive and its source archive, downloaded on 26 September 2026 from the Zenodo record 10.5281/zenodo.22678971:\\[2pt]
{\footnotesize\begin{tabular}{@{}ll@{}}
\texttt{00\_ERDOS\_STRAUSS\_Project\_Reader.pdf} & MD5 \texttt{c1ee553967e0e6f12d4976e171c11b15}\\
\texttt{01\_ERDOS\_STRAUSS\_Source\_Verification\_and\_Machine\_Memory.zip} & MD5 \texttt{4129ec50d0c02ff7daf61fde6c589de4}
\end{tabular}}\\[2pt]
The main source file of the archive (60{,}561 lines) and its fourteen satellite files were read where they are cited.
\item The supplement's source, in the repository folder\newline \texttt{research/expanded-2026-09-08/exact\_bounded\_transport\_72/}.
\item For the literature statements of §1: the source of Elsholtz and Tao's paper (arXiv 1107.1010v6), which the repository holds in its literature folders.
\end{itemize}

\paragraph{What was refereed earlier.}
Theorem~\ref{thm:shell} and Propositions~\ref{prop:r3}--\ref{prop:mod840}, the audited items of Section~\ref{sec:items}, and the audit of the package ES-09 in Section~\ref{ssec:es-s6} come from the Erdős--Straus section of the reader of four workbenches. There they went through five stages:
\begin{enumerate}[label=\arabic*.]
\item a first-pass audit by a separate Claude instance, which read the workbench and wrote independent checks (\note{43a\_});
\item a referee pass on the notes \note{43\_}--\note{47\_}, by an instance that had written none of them; it found, among other things, that the orbit statements of Section~\ref{ssec:es-s6} fail at even levels;
\item a referee pass on that reader;
\item two further passes that read it for understatement and missed generality as well as for overstatement.
\end{enumerate}
These passes produced the range $x<p/2$ in Theorem~\ref{thm:shell}, the first form of Proposition~\ref{prop:survivors}, Proposition~\ref{prop:mod840}, and the extension of the count of Proposition~\ref{prop:r3} to $p\equiv5\pmod{12}$. The checks are \note{checks/\allowbreak workbenches/\allowbreak workbench\_reader\_claims\_checks.py} and \note{checks/\allowbreak workbenches/\allowbreak generality\_checks.py}; the referee reports are in \note{checks/workbenches/}.

\paragraph{What version~1 added to the reader of four workbenches.}
\emph{New} here means new relative to that earlier reader, not a claim of new mathematics. Where the archive already has a statement, this paper cites it.
\begin{itemize}
\item Proposition~\ref{prop:jacobi}, which is the archive's Theorem~10.9 and Proposition~10.10, re-proved;
\item Lemma~\ref{lem:leastclass} and Proposition~\ref{prop:records}, with the reproduction of the deficit data in §\ref{ssec:records};
\item Theorem~\ref{thm:families}, and the statement that each family meets every hard class;
\item Proposition~\ref{prop:twoshears};
\item Proposition~\ref{prop:quartic}(a), whose relation was also checked in \note{21\_};
\item the map of the archive (§\ref{ssec:archive}), the consistency of the census counts (§\ref{ssec:census}), and the three new rows of the overlap table (Section~\ref{sec:overlaps}).
\end{itemize}
These went through the referee pass described next.

\paragraph{The referee pass of version~1.}
An independent Claude instance (\texttt{claude-opus-5-5}), which wrote none of the text, refereed version~1 on 27 September 2026. It read in both directions, as the author asked: for statements stronger than their proofs or sources, and for understatement, missed generality and missed or implied results. It reran the five scripts of version~1, whose outputs were identical, and wrote its own checks. It found every proof correct as written and every recomputed number in agreement. It reported five major and ten minor findings, and seven results that the reader's material implies. Every finding was verified here before it was applied: by reading the archive at the cited page, or by an independent script.
\begin{itemize}
\item \emph{Attributions.} Theorem~\ref{thm:shell}(b) is not a generalisation of the archive. The archive's Theorems~9.1 and~9.64 already hold for every odd prime; the paper now says so. The check here also found that the count $2|E_a|+|M_a|$ is in the archive's Theorem~9.64, which the referee had credited to version~1.
\item \emph{Understatements.} Proposition~\ref{prop:survivors} now concludes that the survivors lie in the hard classes modulo $840$, as the archive's Corollary~10.7 does. The census domain is stated. The page map of §24 is split, and the page counts are made consistent.
\item \emph{Two misdescriptions.}
\begin{itemize}
\item A remark of version~1, that the census carriers are not polynomial identities alone, was wrong. They are polynomial identities on their rows, and Proposition~\ref{prop:nosquare} and Corollary~\ref{cor:censuslimit} replace it.
\item Proposition~\ref{prop:twoshears}(b) had suggested that the theorem fails for every $p\equiv2\pmod3$. It fails only on a minority of them, and Proposition~\ref{prop:sheargraph} and Remark~\ref{rem:sheardensity} now describe that set.
\end{itemize}
\item \emph{Minor findings.} Among them:
\begin{itemize}
\item the residual $R=2$ at $p=2$;
\item the archive's correction of Ionascu and Wilson's label for $1201$;
\item the count of \texttt{.lean} files ($67$, not $84$);
\item the title-page date ($30$ August);
\item Elsholtz and Tao's definition of solvability by polynomials;
\item the full citation of Vaughan's paper.
\end{itemize}
\item \emph{Results adopted.} After verification, these are Corollary~\ref{cor:hle2}, Propositions~\ref{prop:groupbarrier}, \ref{prop:groupconverse}, \ref{prop:sheargraph}, \ref{prop:fixeddivisor} and~\ref{prop:nosquare}, Corollary~\ref{cor:censuslimit}, Proposition~\ref{prop:quartic}(b), the records to $3\cdot10^9$, and the even levels of §\ref{ssec:es-s6}.
\item \emph{A correction to the referee.} It gave $3{,}533{,}141$ as the prime with a shear path of length $5$. A complete enumeration here shows that the least such prime is $1{,}183{,}451$.
\item \emph{Not adopted.} The referee also reported an observation of Ionascu and Wilson on residues modulo $9240$. It could not be checked against their paper here, and is omitted.
\end{itemize}
The referee's report and its check programs are in \note{checks/es\_reader/referee/}; the verification scripts of version~1 are in \note{checks/es\_reader/referee\_verify/}.

\paragraph{The scripts.}
They import no workbench code and recompute every number from the definitions. They are republished in \texttt{contrib/claude-ab/rewrite/checks/es/}; the folders \note{referee/} and \note{referee\_verify/} are in the branch history (see below).

\begin{center}\small
\begin{tabularx}{\linewidth}{@{}>{\raggedright\arraybackslash}p{0.25\linewidth}>{\raggedright\arraybackslash}X>{\raggedright\arraybackslash}p{0.13\linewidth}@{}}
\toprule
Script & What it checks & Time, result\\
\midrule
\note{esr\_records.py} & $h(p)$ for every prime up to $8{,}803{,}369$ through Lemma~\ref{lem:leastclass} and Theorem~\ref{thm:shell}; the strict records; $h(806521)=h(118801)=15$ (for $p=2$ the script prints $R=4h-1=3$; the text's $R=2$ is $4x-p$ with the least denominator $x=1$) & 3 s; as stated\\
\note{esr\_lemma\_checks.py} & Proposition~\ref{prop:jacobi} on all $47{,}210$ shells meeting its hypothesis at odd primes below $3000$ (none occupied), and the failure of its converse; Lemma~\ref{lem:leastclass} against a brute-force enumeration of all solutions at the odd primes below $700$ & 25 s; all pass\\
\note{esr\_deficits.py} & the deficit data of §\ref{ssec:records} at $73$, $1129$, $1201$, $21169$, $118{,}801$, $8{,}803{,}369$ and $806{,}521$, their independence of the primitive root, the corollary of archive §16, and the counts $|E_a|=0$, $|M_a|=2$ at $p_*$ (22 checks) & 1 s; all pass\\
\note{esr\_families.py} & the four identities of Theorem~\ref{thm:families} as polynomial identities, exact fractions for $k\le200$, the hard classes met and the least prime in each & 1 s; all pass\\
\note{esr\_twoshears.py} & Proposition~\ref{prop:twoshears}: the full shear graph of every prime below $2000$, and a complete search for two-edge paths up to $10^7$ & 23 s; all pass\\
\note{referee\_verify/\allowbreak shears\_independent.py} & Proposition~\ref{prop:sheargraph}: degrees on the full graph below $1200$; complete enumeration of paths of length $2$ to $5$ below $10^7$; the class $131\bmod180$ & 21 s; as stated\\
\note{referee\_verify/\allowbreak shears\_density\_\allowbreak bound.py} & the bound $\sum_R\delta_R<0.335$ of Remark~\ref{rem:sheardensity} & 5 s; as stated\\
\note{referee\_verify/\allowbreak hsegment.c}, \note{hsegment\_A.c} & $h(p)$ for all primes below $3\cdot10^9$ by a segmented sieve; the records, the primes with $h\ge19$, and Corollary~\ref{cor:hle2} & 2 min; as stated\\
\note{esr\_referee\_checks.py} & Propositions~\ref{prop:survivors}, \ref{prop:groupbarrier}, \ref{prop:groupconverse}; the census domain; the base rows and the square-row count of Corollary~\ref{cor:censuslimit}; the shell at $p=12289$; the even levels of §\ref{ssec:es-s6} & 2 min; all pass\\
\bottomrule
\end{tabularx}
\end{center}

\paragraph{A partial second pass.}
A second-pass audit of the supplement was started with a separate Claude instance and stopped before it wrote its report. Its scripts and outputs are kept in \note{checks/es\_reader/second\_pass\_partial/}, and §\ref{ssec:supplement} reports what they show. The one reported mismatch was traced to the script's own expectation (§\ref{ssec:supplement}).

\paragraph{Version~3: the author's notes.}
Section~\ref{sec:notes} was added in version~3, on 27 September 2026.
\begin{itemize}
\item Six Claude instances, each of which wrote none of this paper, read the notes in full in separate passes. Each tabulated every numbered statement and checked it with its own programs.
\item claude-ab then re-derived every result proved in Section~\ref{sec:notes} and re-checked it with its own programs, which are in \note{checks/es50/}:
\begin{itemize}
\item \note{es50\_checks.py}: 25 items, all pass, in 49 s. Every certificate is turned into $(x,y,z)$ and tested by $4xyz=p(xy+yz+zx)$.
\item \note{es50\_support\_core.py}: an independent reproduction of the counts $4951$ and $2970$ of Theorem~\ref{thm:termcore}(a), in 4 s.
\item \note{es50\_survivors.py}: Table~\ref{tab:survivors}.
\item \note{salez/}: Theorem~\ref{thm:termcore}(b). The enumerator of Salez's seven modular equations was written by one of the reading passes. It gives $3520$ classes at $840\cdot11\cdot19\cdot23$, and $147{,}348$ at Salez's modulus, which is the number he reports.
\end{itemize}
\item The full record, with the passes' findings and the attributions, is note \note{50\_}.
\item \emph{The referee of version 3.} An independent Claude instance, which wrote none of the text, refereed note \note{50\_} and Section~\ref{sec:notes} on 27 September 2026.
\begin{itemize}
\item It re-ran every program, including the author's own scripts, and recomputed the main numbers by independent code. It found every proof correct.
\item It reported three major findings, all applied: one error charged to the wrong note, an open question false as stated (the prime $409$), and this section's own verification status, which was stated too broadly.
\item It reported fifteen minor findings, applied after checking. Among them: Salez's scope, attributions to the notes and the archive, and the reduced form of Theorem~\ref{thm:visiblebox}.
\item It found results that the material implies. The eight further shifts after Table~\ref{tab:survivors} and the equivalence in Proposition~\ref{prop:groupbarrier}(c) are adopted.
\end{itemize}
Its report and programs are in \note{checks/es50/referee/}.
\end{itemize}

\paragraph{Version~4: the new results.}
The results new in version~4 were proved by claude-ab and checked by programs written for them, which are in \texttt{contrib/claude-ab/rewrite/checks/es/}:
\begin{itemize}
\item \note{exact\_shadow\_reps.py}: Theorem~\ref{thm:m27} and Proposition~\ref{prop:fibonacci}, exactly in $\Q(\zeta_{168})$ and $\Q(\zeta_{120})$, together with the agreement of the $S$-matrix formula with a direct evaluation of the series (to $29$ and $30$ digits);
\item \note{m7\_vs\_M27.py} and \note{m5\_vs\_M25.py}: the same $S$-matrices reproduced numerically from the $q$-series to $38$ digits, the search over signed permutations, conjugation and twists, and the failure of $(ST)^3=S^2$ for the pair of the note's cubic transport;
\item \note{misc\_checks.py}: the values of $A_{11}(pa)$ on the primes $p\equiv1\pmod{840}$, the involutions of $GL_2(3)$, and the coordinate formulas of Proposition~\ref{prop:flowboost};
\item \note{es09\_even\_levels.py} and the scripts of the companion $S^6$ paper \cite{SixSpherePaper} (in \texttt{checks/s6/}): Theorem~\ref{thm:orbits} and Corollary~\ref{cor:orbits}.
\end{itemize}
Proposition~\ref{prop:densityzero}, Proposition~\ref{prop:dossier} and Proposition~\ref{prop:torsor} are proved in the text.

\paragraph{What has not been done.}
The results adopted from the referee of version~1 were proved by that referee and checked by claude-ab, and they have had no further independent reading. Version~4, including its new results, was read by one further independent referee, a Claude instance that wrote none of the text; its findings were verified and applied. The parts of the workbench marked \textsf{located} were not checked, and §\ref{sec:open} lists what else was left out.

\paragraph{Where everything is.}
This paper, its source and the scripts it rests on are on the branch \texttt{claude/claude-ab-grind-20260925} of \texttt{KokunoYumeto/zeta-function-research-reader}: the paper in \texttt{contrib/claude-ab/rewrite/papers/erdos\_straus/}, the scripts in \texttt{contrib/claude-ab/rewrite/checks/es/}, and the condensed Markdown record of the same material in \texttt{contrib/claude-ab/rewrite/condensed/}. The earlier versions, the audit notes \note{21\_}, \note{43\_}, \note{43a\_}, \note{47\_}, \note{47a\_} and \note{50\_}, and the folders \note{checks/es\_reader/} and \note{checks/es50/} with the referee reports, are in the branch's history at commit \texttt{a90d9e5b}, in \texttt{contrib/claude-ab/grind\_20260925/}.

\section{Changes in version 4}\label{app:changes}

Version~3 of this paper was withdrawn on 27 September 2026, together with the audit notes, because of its framing. Version~4 keeps its theorems, proofs and citations, corrects its framing, and adds new results. The changes are the following.

\paragraph{Framing.}
\begin{itemize}
\item The title, abstract and introduction are rewritten. The introduction states the main results as Theorems~I--V, explains what they mean, describes the bridges, credits the contributors, and sets out the four statuses and how the status labels map onto them.
\item Section~\ref{sec:notes} is retitled \emph{The author's notes, April--June 2026}, and the notes are described in the same way throughout.
\item Quotation marks around the terms of the workbench and the notes are replaced by italics.
\item In the map of the archive (§\ref{ssec:archive}) the category \emph{side} is renamed \emph{structural} and defined through the archive's own description of these sections as an interface still to be supplied by a typed map.
\item Statements that structures which do not take primes as input cannot bear on solvability are replaced by what is established about the typed maps that would connect them (Sections~\ref{sec:notes} and~\ref{sec:overlaps}). The negative results of Section~\ref{sec:negative} are stated with their scope, and the table's closing sentence now says which mechanisms are sharp.
\item The drafts that the author's latest record supersedes are treated only for the results they contain that were verified: the paragraph on one such draft, its row in the catalogue, and the row of Section~\ref{sec:negative} on the terminal core as the residual locus are removed. The terminal core is now described by what Theorem~\ref{thm:termcore} locates (Remark~\ref{rem:termcore}).
\item Statements that rest only on a reading pass over the notes are marked \emph{(a reading pass)} and count as not verified here.
\item Section~\ref{sec:overlaps} is rebuilt as a section on the bridges to other areas, with what is established, what is not verified here, and my assessment for each; the table of edges to the author's other workbenches is its last subsection.
\end{itemize}

\paragraph{Corrections of content.}
\begin{itemize}
\item \emph{The minimal model $M(2,7)$.} Version~3 said that the note's identification of the mock theta shadow with the characters of $M(2,7)$ fails, that the note's modified modular pair is not a representation of $SL_2(\ZZ)$, and that no relabelling matches both $S$ and $T$. Re-deriving the statement shows that the identification holds after conjugation and the twist by the character of $\eta^3$ (Theorem~\ref{thm:m27}). What fails is the note's matching of labels through its cubic phase transport: no signed permutation and scalar realise it, and the pair it produces does not satisfy $(ST)^3=S^2$; version~3's statements hold in that scope only.
\item \emph{The Lorentzian note.} Version~3 listed the note's section on modular flow as an exception to the correctness of the notes' structural links, stating that the modular flow is a rotation and not a Lorentz boost. That statement stands (Proposition~\ref{prop:flowboost}(a)). What changes is the verdict on the note: its proposition, that the congruence action of $A_\rho$ is a Lorentz boost, is correct, and Proposition~\ref{prop:flowboost} gives the exact relation between the two.
\item \emph{The sieve.} Version~3 said that a sieve bound of the form $x/(\log x)^\kappa$ follows from the quadratic-residue conditions and that the conditions do not approach a proof. The first statement was not proved there and is removed; Proposition~\ref{prop:densityzero} proves density zero, and Question~\ref{q:qrsurvivors} states what a negative answer would prove.
\item \emph{The shell count as a class function.} Version~3 said that $A_{11}(p)$ takes the values $0,4,6,8,10$ on primes $p\equiv1\pmod{840}$. The count meant is $A_{11}(pa)$ with $a=(p+11)/4$, the number of ordered certificates of the shell; $A_{11}(p)$ itself takes only the values $0$ and $2$ at primes. Section~\ref{ssec:noteserrors} now states it so, with the primes at which each value occurs.
\item \emph{The further certificates of the referee of version~3.} They are now stated with their proofs, and the $23$ primes below $10^7$ are described as the primes $p\equiv1\pmod{24}$ with no certificate of the kinds considered, rather than as survivors of conditions that version~3 did not state.
\item \emph{The Busy-Beaver notes.} Version~3's statement that the residues modulo $107$ are at chance level is replaced by the description of the test and its scope (Section~\ref{ssec:structfacts}), and its statement that the only link with the Busy-Beaver function is the classical one by the exact $\Pi_1$ bridge (Section~\ref{sec:overlaps}).
\end{itemize}

\paragraph{New results.}
Proposition~\ref{prop:densityzero} (density zero of the survivors of the row $p+1$); Theorem~\ref{thm:m27} (the order-$7$ shadow carries the $M(2,7)$ representation); Proposition~\ref{prop:fibonacci} (the order-$5$ block carries the Fibonacci data); Propositions~\ref{prop:dossier} and~\ref{prop:torsor}, which state as propositions, with proofs, two points that version~3 listed as errors; Proposition~\ref{prop:flowboost} (modular flow and boost); and Theorem~\ref{thm:orbits} with Corollary~\ref{cor:orbits} (the orbits of the $(3,4,\infty)$ covers at every level), which answer version~3's question on the covers at even levels. The coefficients of the note's modular differential equation are checked in Section~\ref{ssec:mock}. Section~\ref{sec:overlaps} also states, from the companion note \cite{ApollonianNote}, how rotations and boosts meet the Apollonian group, and, from a reading pass, that the note's odd Ogg phase decomposition holds with the map of Proposition~\ref{prop:torsor}.

\paragraph{Questions.}
Section~\ref{sec:open} is reordered. Version~3's question on the covers at even levels is answered and replaced by the question on their genera; Questions~\ref{q:brieskorn}--\ref{q:pi1} are new.

Apart from the corrections listed above, no statement of version~3 changed its mathematical content.

\addcontentsline{toc}{section}{References}
\begin{thebibliography}{99}
\small
\bibitem{AlpogeS6} \emph{A compact complex threefold fibred by tori over the projective line, and the six-sphere}, manuscript circulated by L.~Alpöge, 2026, \url{https://alpo.ge/s6.pdf}.
\bibitem{CCFGH} M.~C.~N.~Cheng, S.~Chun, F.~Ferrari, S.~Gukov, S.~M.~Harrison, \emph{3d modularity}, J. High Energy Phys. \textbf{2019}, no.~10, 010.
\bibitem{CFS} M.~C.~N.~Cheng, F.~Ferrari, G.~Sgroi, \emph{Three-manifold quantum invariants and mock theta functions}, Phil. Trans. R. Soc. A \textbf{378} (2020), 20180439, arXiv:1912.07997.
\bibitem{ElsholtzTao} C.~Elsholtz, T.~Tao, \emph{Counting the number of solutions to the Erdős--Straus equation on unit fractions}, J. Aust. Math. Soc. \textbf{94} (2013), 50--105, doi:10.1017/S1446788712000468.
\bibitem{IonascuWilson} E.~J.~Ionascu, A.~Wilson, \emph{On the Erdős--Straus conjecture}, arXiv:1001.1100 (2010).
\bibitem{Kneser} M.~Kneser, \emph{Ein Satz über abelsche Gruppen mit Anwendungen auf die Geometrie der Zahlen}, Math. Z. \textbf{61} (1955), 429--434.
\bibitem{Lagarias} J.~C.~Lagarias, \emph{The $3x+1$ problem and its generalizations}, Amer. Math. Monthly \textbf{92} (1985), 3--23.
\bibitem{Lopez} M.~A.~Lopez, \emph{A complete congruence system for the Erdős--Straus conjecture}, arXiv:2404.01508 (2024).
\bibitem{MihneaDumitru} Spiridon Mihnea, Dumitru C.~Bogdan (author names as listed on arXiv), \emph{Further verification and empirical evidence for the Erdős--Straus conjecture}, arXiv:2509.00128 (2025).
\bibitem{MontgomeryVaughan} H.~L.~Montgomery, R.~C.~Vaughan, \emph{The large sieve}, Mathematika \textbf{20} (1973), 119--134.
\bibitem{Mordell} L.~J.~Mordell, \emph{Diophantine Equations}, Academic Press, London, 1969.
\bibitem{RosserSchoenfeld} J.~B.~Rosser, L.~Schoenfeld, \emph{Approximate formulas for some functions of prime numbers}, Illinois J. Math. \textbf{6} (1962), 64--94 (Theorem~15).
\bibitem{Salez} S.~E.~Salez, \emph{The Erdős--Straus conjecture: new modular equations and checking up to $N=10^{17}$}, arXiv:1406.6307 (2014).
\bibitem{Vaughan} R.~C.~Vaughan, \emph{On a problem of Erdős, Straus and Schinzel}, Mathematika \textbf{17} (1970), 193--198, doi:10.1112/S0025579300002886.
\bibitem{Zagier} D.~Zagier, \emph{Ramanujan's mock theta functions and their applications [d'après Zwegers and Ono--Bringmann]}, Séminaire Bourbaki, exp.~986, Astérisque \textbf{326} (2009), 143--164.
\bibitem{ZetaReader} Claude (Anthropic), \emph{The split-zero programme around the Riemann zeta function: its results, with proofs, and its bridges}, version~3, 2026, in the Zenodo record \emph{Split-Zero cohomology and the zeta-function research programme}, concept DOI 10.5281/zenodo.22678085; source on the branch \texttt{claude/claude-ab-grind-20260925} of \texttt{KokunoYumeto/zeta-function-research-reader}, folder \texttt{contrib/claude-ab/rewrite/papers/zeta/} (Lemma~6.6 has the same number in versions 1--3).
\bibitem{WorkbenchReader} Claude (Anthropic), \emph{Four workbenches and where they meet: the Erdős--Straus, Collatz, Erdős Problem 817 and Yang--Mills workbenches}, 2026, in the same Zenodo record as \cite{ZetaReader}; source in \texttt{contrib/claude-ab/rewrite/papers/workbenches/} on the same branch.
\bibitem{SixSpherePaper} Claude (Anthropic), \emph{The $S^6$ record and the integral monodromy of its $(3,4,\infty)$ period system}, 2026, in the Zenodo record 10.5281/zenodo.22678442 and its later versions; source in \texttt{contrib/claude-ab/rewrite/papers/s6/} on the same branch.
\bibitem{ApollonianNote} Claude (Anthropic), \emph{Rotations, boosts and the Apollonian group}, 2026, a companion note to this paper; source in \texttt{contrib/claude-ab/rewrite/papers/apollonian/} on the same branch.
\bibitem{Aaronson} S.~Aaronson, \emph{The Busy Beaver frontier}, SIGACT News \textbf{51} (2020), no.~3, 32--54.
\bibitem{AndrewsGarvan} G.~E.~Andrews, F.~G.~Garvan, \emph{Ramanujan's ``lost'' notebook VI: the mock theta conjectures}, Adv. Math. \textbf{73} (1989), 242--255.
\bibitem{BaezOctonions} J.~C.~Baez, \emph{The octonions}, Bull. Amer. Math. Soc. \textbf{39} (2002), 145--205.
\bibitem{CCKPS} M.~C.~N.~Cheng, I.~Coman, P.~Kucharski, D.~Passaro, G.~Sgroi, \emph{3d modularity revisited}, arXiv:2403.14920 (2024).
\bibitem{CDH} M.~C.~N.~Cheng, J.~F.~R.~Duncan, J.~A.~Harvey, \emph{Umbral moonshine and the Niemeier lattices}, Res. Math. Sci. \textbf{1} (2014), Art.~3.
\bibitem{ConwayNorton} J.~H.~Conway, S.~P.~Norton, \emph{Monstrous moonshine}, Bull. London Math. Soc. \textbf{11} (1979), 308--339.
\bibitem{ConwaySloane} J.~H.~Conway, N.~J.~A.~Sloane, \emph{Sphere Packings, Lattices and Groups}, 3rd ed., Springer, 1999, ch.~16--18.
\bibitem{EichlerZagier} M.~Eichler, D.~Zagier, \emph{The Theory of Jacobi Forms}, Birkhäuser, 1985.
\bibitem{GLMWY} R.~L.~Graham, J.~C.~Lagarias, C.~L.~Mallows, A.~R.~Wilks, C.~H.~Yan, \emph{Apollonian circle packings: geometry and group theory I. The Apollonian group}, Discrete Comput. Geom. \textbf{34} (2005), 547--585.
\bibitem{Hickerson} D.~Hickerson, \emph{A proof of the mock theta conjectures}, Invent. Math. \textbf{94} (1988), 639--660.
\bibitem{Kocik} J.~Kocik, \emph{The Koide lepton mass formula and geometry of circles}, arXiv:1201.2067 (2012).
\bibitem{Moreno} G.~Moreno, \emph{The zero divisors of the Cayley--Dickson algebras over the real numbers}, Bol. Soc. Mat. Mexicana (3) \textbf{4} (1998), 13--28.
\bibitem{RochaCaridi} A.~Rocha-Caridi, \emph{Vacuum vector representations of the Virasoro algebra}, in: Vertex Operators in Mathematics and Physics, MSRI Publ.~3, Springer, 1985, 451--473.
\bibitem{RSW} E.~Rowell, R.~Stong, Z.~Wang, \emph{On classification of modular tensor categories}, Comm. Math. Phys. \textbf{292} (2009), 343--389.
\bibitem{SarnakLetter} P.~Sarnak, \emph{Letter to J.~Lagarias about integral Apollonian packings}, June 2007.
\end{thebibliography}

\end{document}
